\chapter[Fine-grained analysis:  $\ell_{\infty}$ and $\ell_{2,\infty}$ perturbation theory]{Fine-grained spectral analysis: \\ $\ell_{\infty}$ and $\ell_{2,\infty}$ perturbation theory}
\label{cha:Linf-theory}


In a growing number of applications, the $\ell_2$-type distance between subspaces, which is the central subject studied in Chapter~\ref{chap:application-L2}, turns out to be inadequate for performance characterization. Rather, what would be of interest is the entrywise behavior of the eigenvector and the matrix under consideration. This is especially important when the individual entries of the eigenvector or the matrix of interest carry pivotal operational meanings. For example, in a recommendation system, one might be interested in controlling the prediction error of a user's preference on a specific product, which  concerns  a specific entry in a user-product rating matrix; in sensor network localization, one might seek to control the ranging error w.r.t.~a pair of sensors, which corresponds to  entrywise prediction errors in a Euclidean distance matrix; and last but not least, in community recovery, the entries of the leading eigenvector of a certain data matrix might encode the community membership associated with each individual (as explained in Section~\ref{sec:community-detection}).

Tackling the preceding applications calls for development of fine-grained spectral analysis beyond classical $\ell_2$ perturbation theory.  To be more precise, consider once again the observation model
$$
\bm{M}=\bm{M}^{\star} + \bm{E}
$$  previously studied in Chapter~\ref{cha:matrix-perturbation} (cf.~\eqref{eq:perturbed-M}). The sort of fine-grained theory being sought after gravitates around the following questions concerned with $\ell_{\infty}$ and/or $\ell_{2,\infty}$ perturbation:
%
\begin{itemize}
	\item For a symmetric matrix $\bm{M}^{\star}$, how to characterize the effect of $\bm{E}$ on the $\ell_{\infty}$ perturbation  of the leading eigenvector,
		or the $\ell_{2,\infty}$ perturbation of the rank-$r$ leading eigenspace?
	\item For a general matrix $\bm{M}^{\star}$, how to pin down the $\ell_{\infty}$ perturbation of the leading singular vector, or the $\ell_{2,\infty}$ perturbation  of the rank-$r$ leading singular subspace, in response to the perturbation $\bm{E}$?
	\item How to assess the entrywise estimation error of the matrix estimate produced by the spectral method, and how is it affected by $\bm{E}$?
\end{itemize}
%
Unfortunately, a direct application of classical $\ell_2$ perturbation theory typically leads to overly crude bounds when coping with the above questions.  In particular, when the $\ell_2$ error is approximately evenly distributed across entries, naively upper bounding the entrywise error by the $\ell_2$ error is often loose by an order-of-magnitude.
In order to conquer such limitations,
this chapter introduces a modern suite of techniques that delivers tight $\ell_\infty$ and $\ell_{2,\infty}$ error control by leveraging the statistical nature of data models. 





%\section{An illustrative example: Rank-1 matrix denoising}
\section{Leave-one-out analysis: An illustrative example}
\label{sec:rank-1-denoising-LOO}



To paint a high-level picture of the core ideas empowering the $\ell_{\infty}$ and $\ell_{2,\infty}$ analysis, we find it helpful to first look at a pedagogical example of rank-1 matrix denoising, a special case of the formulation introduced in Section~\ref{sec:matrix-denoising-setup}.


\subsection{Setup and algorithm}
\label{sec:setup-rank1-denoising}
Suppose that we observe a noisy copy of an unknown rank-1 matrix $\bm{M}^{\star}$ as follows
%
\begin{equation}
	\bm{M}=\bm{M}^{\star}+\bm{E}=\lambda^{\star}\bm{u}^{\star}\bm{u}^{\star\top}+\bm{E} \in \mathbb{R}^{n\times n},
	\label{eq:defn-matrix-denoising-rank1}
\end{equation}
%
where $\lambda^{\star}>0$ and $\bm{u}^{\star}\in \mathbb{R}^{n}$
represent the largest eigenvalue  of $\bm{M}^{\star}$ and its associated eigenvector, respectively.  We assume the Gaussian noise model as in Section~\ref{sec:matrix-denoising-setup}, namely, $\bm{E}$ is a symmetric matrix whose upper triangular part comprises of i.i.d.~entries drawn from $\mathcal{N}(0,\sigma^2)$.
 In addition, we remind the readers of the following incoherence parameter $\mu$  (cf.~Definition~\ref{assump:mc-incoherence}):
%
\begin{align}
	%\|\bm{u}^{\star}\|_{\infty}= \sqrt{\mu/n} = \sqrt{\mu/n} \,\|\bm{u}^{\star}\|_2,
	\mu = n\|\bm{u}^{\star}\|_{\infty}^{2},
	\label{defn:incoherence-rank-1-denoising}
\end{align}
%
which satisfies $1 \leq \mu \leq n$ in this rank-1 case; see Remark~\ref{remark:range-mu-MC}.



Letting $\lambda$  be the leading eigenvalue of $\bm{M}$ (i.e., $\lambda = \lambda_{1}(\bm{M})$) and $\bm{u}\in \mathbb{R}^{n}$
 the associated eigenvector, the spectral method attempts to estimate $\bm{u}^{\star}$ using $\bm{u}$. In this section, we are particularly interested in controlling the entrywise error, defined in terms of the $\ell_{\infty}$ distance (modulo the global sign):
 \begin{align}
	\mathsf{dist}_{\infty}\big(\bm{u},\bm{u}^{\star}\big)
	\coloneqq \min\big\{ \|\bm{u}-\bm{u}^{\star}\|_{\infty}, \|\bm{u}+\bm{u}^{\star}\|_{\infty} \big\}.
	\label{eq:defn-dist-infty-vector}
\end{align}




\subsection{$\ell_{\infty}$ performance guarantees}



While Section~\ref{sec:performance-denoising-L2} delivers $\ell_2$ estimation guarantees for the spectral estimate $\bm{u}$,
it falls short of characterizing the entrywise behavior---except for the crude and highly suboptimal bound  $\mathsf{dist}_{\infty}(\bm{u},\bm{u}^{\star})\leq \mathsf{dist}(\bm{u},\bm{u}^{\star})$. Encouragingly,  this simple spectral method is provably  accurate in an entrywise fashion, as revealed by the following theorem.
%
 \begin{theorem}
	 Consider the settings in Section~\ref{sec:setup-rank1-denoising}.
	 There exists some sufficiently small constant $c_0>0$ such that if $\sigma \sqrt{n} \leq c_0 \lambda^{\star}$, then
 %
 \begin{align}
	 \mathsf{dist}_{\infty}\big(\bm{u},\bm{u}^{\star}\big)
	 \lesssim \frac{\sigma (\sqrt{\log n} + \sqrt{\mu} )}{\lambda^{\star}}
 \end{align}
  %
	 holds with probability exceeding $1-O(n^{-8})$.
  \end{theorem}
%

In particular, if the incoherence parameter obeys $\mu \lesssim \log n$ (the case where no entries of $\bm{u}^{\star}$ are significantly larger in magnitude than the average magnitude), then our $\ell_{\infty}$ bound reads
%
\begin{align}
	 \mathsf{dist}_{\infty}\big(\bm{u},\bm{u}^{\star}\big)
	 \lesssim \frac{\sigma\sqrt{\log n}}{\lambda^{\star}} ,
 \end{align}
%
which is about $\sqrt{n/\log n}$ times smaller than the $\ell_2$ error bound \eqref{eq:dist-U-Ustar-denoising-L2}, that is, $\mathsf{dist}\big(\bm{u},\bm{u}^{\star}\big)   \lesssim \frac{\sigma\sqrt{ n}}{\lambda^{\star}}$.
This implies that the estimation errors of $\bm{u}$ are dispersed  more or less evenly across all entries---a message that is previously unavailable from classical $\ell_2$ perturbation theory.




\subsection{Key ingredient and intuition: Leave-one-out estimates}
\label{sec:LOO-estimates-denoising}

To facilitate entrywise analysis, a crucial ingredient lies in the introduction of a set of leave-one-out auxiliary estimates, detailed below.


\paragraph{Construction of leave-one-out estimates.} For each $1\leq l\leq n$, let us construct an auxiliary matrix  $\bm{M}^{(l)}$ as follows
%
\begin{align}
	\bm{M}^{(l)} \coloneqq \lambda^{\star}\bm{u}^{\star}\bm{u}^{\star\top}+\bm{E}^{(l)} ,
	\label{eq:defn-Ml-denoising}
\end{align}
%
where the noise matrix $\bm{E}^{(l)}$ is generated according to
%
\begin{equation}
E_{i,j}^{(l)} \coloneqq
\begin{cases}
E_{i,j},\qquad & \text{if }i\neq l\text{ and }j\neq l, \\
0, & \text{else}.
\end{cases}
\label{eq:Eij-l-defn}
\end{equation}
%
In words, $\bm{M}^{(l)}$ (resp.~$\bm{E}^{(l)}$) is obtained by leaving out the randomness in the $l$-th row/column of $\bm{M}$ (resp.~$\bm{E}$). The pattern of the leave-one-out construction is illustrated in Figure~\ref{fig:loo_illustration}.
Let $\lambda^{(l)}$ and $\bm{u}^{(l)}$ denote respectively the leading eigenvalue and leading eigenvector of $\bm{M}^{(l)}$; 
these leave-one-out estimates are introduced solely for analysis purpose. 
It is important to recognize that by construction, $\bm{M}^{(l)}$ (and hence $\bm{u}^{(l)}$) is independent of the noise in the $l$-th row/column of $\bm{E}$, a fact that plays a pivotal role in controlling the perturbation of the $l$-th entry of $\bm{u}^{\star}$.

\begin{figure}[t]
\begin{center}
\includegraphics[width=0.85\textwidth]{figures/loo_illustration}
\end{center}
\caption{Illustration of the leave-one-out auxiliary matrix $\bm{M}^{(l)}$, which removes all the noise in the $l$-th row and the $l$-th column of $\bm{M}$.} \label{fig:loo_illustration}
\end{figure}


\paragraph{Intuition.} Before delving into the proof, let us first explain the rationale at an intuitive level.
%
\begin{enumerate}
	\item Given that $\bm{u}^{(l)}$ is obtained by dropping only a tiny fraction of the data, we expect $\bm{u}^{(l)}$ to be exceedingly close to $\bm{u}$, i.e.,
	%
	\begin{align}
		\bm{u} \approx \pm \bm{u}^{(l)}.
	\end{align}
	%
	In words,  $\bm{u}^{(l)}$ forms a reliable surrogate of $\bm{u}$, which motivates us to analyze $\bm{u}^{(l)}$ instead (if there are foreseeable benefits to do so).


\item The way we construct $\bm{u}^{(l)}$ makes it particularly convenient to analyze the behavior of the $l$-th entry, denoted by $u_l^{(l)}$.
More specifically, given that $(\lambda^{(l)},\bm{u}^{(l)})$ is an eigenpair of $\bm{M}^{(l)}$, one has (assuming for the moment that $\lambda^{(l)} \neq 0$)
%
	\begin{align}
		u_{l}^{(l)} & =\frac{1}{\lambda^{(l)}}\bm{M}_{l,\cdot}^{(l)}\bm{u}^{(l)}=\frac{1}{\lambda^{(l)}}\bm{M}_{l,\cdot}^{\star}\bm{u}^{(l)}=\frac{\lambda^{\star}}{\lambda^{(l)}}u_{l}^{\star}\bm{u}^{\star\top}\bm{u}^{(l)}   \label{eq:ul_l-expression-denoising} \\
		& \approx \pm u_{l}^{\star}.
	\end{align}
	%
	Here, the first line follows since, by design, the $l$-th rows of $\bm{M}^{(l)}$
		and $\bm{M}^{\star}$ coincide (both of which are given by $\lambda^{\star}u_{l}^{\star}\bm{u}^{\star\top}$), whereas the second line holds as long as the size $\sigma$ of the noise is sufficiently small, so that $ \lambda^{(l)} / \lambda^{\star} \approx 1$ and $\bm{u}^{\star\top}\bm{u}^{(l)}\approx \pm 1$ according to the $\ell_2$ perturbation theory.
	% where the second identity follows by construction of $\bm{M}^{(l)}$.
	
	%% this point below is justified more rigorously in the formal proof, and hence omitted from the intuition as it is more technical
	
%\item The above $\ell_2$-pertubation theory requires evaluation of $\|\bE^{(l)}\|$.
%	By construction, for any $n$-dimensional vector $\bx$ with restriction $x_l = 0$, we have $\bx^{\top} \bE \bx = \bx^{\top} \bE^{(l)} \bx$.  Therefore,
%$$
%		\| \bE \| \geq \sup_{\|\bx\|_2=1, x_l = 0} \bx^{\top} \bE \bx =  \sup_{\|\bx\|_2=1, x_l = 0}  \bx^{\top} \bE^{(l)} \bx =
%\| \bE ^{(l)}\|
%$$
\end{enumerate}
%
Combining the above observations suggests that $u_{l} \approx \pm u_{l}^{(l)} \approx \pm u_{l}^{\star}$.


\subsection{Leave-one-out analysis}
\label{sec:LOO-analysis-rank1-denoising}
Now we make rigorous the heuristic argument in the last subsection, which relies heavily on careful statistical analysis.


\subsubsection{Preparation: what we have learned from $\ell_2$ perturbation theory}
%
Let us start by collecting a few results derived from the $\ell_2$ perturbation theory in Section~\ref{sec:matrix-denoising-L2} for handy reference. Experienced readers can proceed directly to Step 1.


Specifically, suppose that $\sigma\sqrt{n}\leq\frac{1-1/\sqrt{2}}{5}\lambda^{\star}$. Then with probability at least $1-O(n^{-8})$,
%
\begin{subequations}
\label{eq:L2-perturbation-summary-denoising}
%
\begin{align}
  \|\bm{E} \|  &\leq 5\sigma \sqrt{n}    & \|\bm{E}^{(l)} \| &\leq \|\bm{E} \| \leq 5\sigma \sqrt{n}  && \\
  \mathsf{dist}( \bm{u} , \bm{u}^{\star} )  &\leq \frac{10\sigma\sqrt{n}}{\lambda^{\star}}  ~~
 & \mathsf{dist}( \bm{u}^{(l)} , \bm{u}^{\star} )  &\leq \frac{10\sigma\sqrt{n}}{\lambda^{\star}} && \\
	| \lambda - \lambda^{\star} |  &\leq 5\sigma \sqrt{n} & | \lambda^{(l)} - \lambda^{\star} |  &\leq 5\sigma \sqrt{n} &&\\
	\max_{j: j\geq 2} | \lambda_j(\bm{M})   |  &\leq 5\sigma \sqrt{n} & \max_{j: j\geq 2}| \lambda_j(\bm{M}^{(l)}) |  &\leq 5\sigma \sqrt{n}
\end{align}
%
\end{subequations}
%
hold simultaneously for all $1\leq l\leq n$. Here, the first line arises from~\eqref{eq:noise-matrix-E-bound-denoising}, and the remaining claims follow the same argument as in the proof of Corollary~\ref{cor:davis-kahan-conclusion-corollary}. 
 Consequently, there exist global signs $z, z_l \in \{1,-1\}$  obeying $\|z\bm{u} - \bm{u}^{\star}\|_2= \mathsf{dist}( \bm{u}, \bm{u}^{\star} ) \leq 10\sigma\sqrt{n} / \lambda^{\star}$ and $\|z_l\bm{u}^{(l)} - \bm{u}^{\star}\|_2= \mathsf{dist}( \bm{u}^{(l)}, \bm{u}^{\star} ) \leq 10\sigma\sqrt{n} / \lambda^{\star}$. To simplify presentation, we shall assume
 \begin{subequations}
\label{eq:ul-ustar-WLOG-denoising}
\begin{align}
	\|\bm{u} - \bm{u}^{\star}\|_2 &= \mathsf{dist}( \bm{u}, \bm{u}^{\star} ), \quad \\
	\big\| \bm{u}^{(l)} - \bm{u}^{\star} \big\|_2 &= \mathsf{dist}( \bm{u}^{(l)}, \bm{u}^{\star} ) , \quad 1 \leq l \leq n
\end{align}
\end{subequations}
%
without loss of generality.
As a simple yet useful byproduct: if ${20\sigma{\sqrt{n}}} < \lambda^{\star}$, then Condition \eqref{eq:ul-ustar-WLOG-denoising} necessarily implies
%
\begin{align}
	\big\| \bm{u} -  \bm{u}^{(l)} \big\|_2 = \mathsf{dist}\big( \bm{u}, \bm{u}^{(l)}\big), \qquad 1 \leq l \leq n.  	
	\label{eq:u-ul-rotation-denoising}
\end{align}
To see this, combine the triangle inequality and \eqref{eq:L2-perturbation-summary-denoising} to yield
%
\begin{align*}
	\big\|\bm{u}-\bm{u}^{(l)}\big\|_{2} & \leq\big\|\bm{u}-\bm{u}^{\star}\big\|_{2}+\big\|\bm{u}^{(l)}-\bm{u}^{\star}\big\|_{2}
	\leq\frac{20\sigma\sqrt{n}}{\lambda^{\star}}<1 ,
\end{align*}
%
which taken collectively with the fact $\|\bm{u} \|_{2} = \|\bm{u}^{(l)}\|_{2}=1$ gives
%
\[
\big\|\bm{u}+\bm{u}^{(l)}\big\|_{2}^{2}=
2\big\|\bm{u}\big\|_{2}^{2}+2\big\|\bm{u}^{(l)}\big\|_{2}^2
-\big\|\bm{u}-\bm{u}^{(l)}\big\|_{2}^{2}>1>\big\|\bm{u}-\bm{u}^{(l)}\big\|_{2}^{2}.
\]
%
This together with the definition \eqref{eq:dist_UUstar-spectral} of $\mathsf{dist}(\cdot,\cdot)$ validates  \eqref{eq:u-ul-rotation-denoising}.



\subsubsection{Step 1: bounding the proximity of leave-one-out $\&$ true estimates}

In this step, we seek to control the distance between the true estimate $\bm{u}$ and the leave-one-out estimate $\bm{u}^{(l)}$. Suppose for the moment that
\begin{equation}\label{eq:noise-small-gap}
\|\bm{M}-\bm{M}^{(l)}\| \leq (1 - 1 / \sqrt{2}) \Big(\lambda^{(l)}-  \max_{j\geq 2}\big| \lambda_{j}\big(\bm{M}^{(l)}\big) \big| \Big).
\end{equation}
We can then invoke the Davis-Kahan theorem (cf.~Corollary \ref{cor:davis-kahan-conclusion-corollary}) and the relation \eqref{eq:u-ul-rotation-denoising} to yield
%
\begin{align}
	\big\|\bm{u}-\bm{u}^{(l)}\big\|_{2} &
	%=\mathsf{dist}\big(\bm{u},\bm{u}^{(l)}\big)
	\leq\frac{2\|\big(\bm{M}-\bm{M}^{(l)}\big)\bm{u}^{(l)}\|_{2}}{\lambda^{(l)}- \max\limits_{j\geq 2}\big| \lambda_{j}\big(\bm{M}^{(l)}\big) \big| }
	\leq \frac{4\|\big(\bm{M}-\bm{M}^{(l)}\big)\bm{u}^{(l)}\|_{2}}{\lambda^{\star}}.
	\label{eq:u-ul-diff-denoising}
\end{align}
%
Here, the last inequality invokes \eqref{eq:L2-perturbation-summary-denoising} and $20\sigma\sqrt{n}\leq \lambda^{\star}$ to obtain
%
\begin{align}
	\lambda^{(l)}-  \max_{j\geq 2}\big| \lambda_{j}\big(\bm{M}^{(l)}\big) \big| & \geq  (\lambda^{\star}-  5\sigma \sqrt{n} ) - 5\sigma \sqrt{n} \geq\lambda^{\star}/2 .
	\label{eq:lambda-l-minus-lambda-j-LB1}
	%\\
	%& \geq \lambda^{\star}- \|\bm{E}^{(l)}\| \geq \lambda^{\star}-10\sigma\sqrt{n}\geq\lambda^{\star}/2,
\end{align}
A byproduct of this calculation is that $\lambda^{(l)}$ is positive for all $1\leq l \leq n$.
%
%provided that $20\sigma\sqrt{n}\leq \lambda^{\star}$.

It thus remains to control the term $\|\big(\bm{M}-\bm{M}^{(l)}\big)\bm{u}^{(l)}\|_{2}$ in \eqref{eq:u-ul-diff-denoising}, towards which certain statistical independence proves crucial. Specifically, we observe that (by construction of $\bm{M}^{(l)}$)
%
\begin{align*}
\big(\bm{M}-\bm{M}^{(l)}\big)\bm{u}^{(l)} & =\bm{e}_{l}\bm{E}_{l,\cdot}\bm{u}^{(l)}+u_{l}^{(l)}(\bm{E}_{\cdot,l} - E_{l,l}\bm{e}_{l}),
\end{align*}
%
where $\bm{e}_{l}$ is the $l$-th standard basis vector,
% $u_{l}^{(l)}$ is the $l$-th entry of $\bm{u}^{(l)}$,
and $\bm{E}_{l,\cdot}$ (resp.~$\bm{E}_{\cdot,l}$)
denotes the $l$-th row (resp.~column) of $\bm{E}$. By construction, $\bm{u}^{(l)}$ is statistically independent of $\bm{E}_{l,\cdot}$, thus indicating that
%
\begin{align}
	\bm{E}_{l,\cdot} \bm{u}^{(l)} ~\sim~ \mathcal{N}(0, \sigma^2 \|\bm{u}^{(l)}\|_2^2) = \mathcal{N}(0, \sigma^2 )
\end{align}
%
conditioned on $\bm{u}^{(l)}$. Hence, with probability at least $1-n^{-10}$,
%
\begin{align}
	\big| \bm{E}_{l,\cdot} \bm{u}^{(l)} \big| \leq 5\sigma \sqrt{\log n}, \qquad 1\leq l\leq n.
\end{align}
%
In addition, $\|\bm{E}_{\cdot,l} - E_{l,l}\bm{e}_{l}\|_{2} \leq \|\bm{E}_{\cdot,l}\|_2\leq \|\bm{E}\|\leq 5\sigma \sqrt{n}$ (cf.~\eqref{eq:L2-perturbation-summary-denoising}). Consequently,
%
\begin{align*}
 & \|\big(\bm{M}-\bm{M}^{(l)}\big)\bm{u}^{(l)}\|_{2}
	\leq \big|\bm{E}_{l,\cdot}\bm{u}^{(l)}\big|  +
	 \big\|\bm{E}_{\cdot,l}\big\|_{2} \cdot \big|u_{l}^{(l)}\big| \\
 & \qquad \leq5\sigma\sqrt{\log n}+\big\|\bm{E}_{\cdot,l}\big\|_{2}\big(\big|u_{l}\big|+\big\|\bm{u}-\bm{u}^{(l)}\big\|_{\infty}\big)\\
 & \qquad \leq5\sigma\sqrt{\log n}+5\sigma\sqrt{n}\|\bm{u}\|_{\infty}+5\sigma\sqrt{n}\big\|\bm{u}-\bm{u}^{(l)}\big\|_{2}.
\end{align*}
%
Substitution into \eqref{eq:u-ul-diff-denoising} yields
%
\begin{align*}
\big\|\bm{u}-\bm{u}^{(l)}\big\|_{2}
 & \leq \frac{20\sigma\sqrt{\log n}+20\sigma\sqrt{n}\|\bm{u}\|_{\infty}+20\sigma\sqrt{n}\big\|\bm{u}-\bm{u}^{(l)}\big\|_{2}}{\lambda^{\star}}\\
 & \leq\frac{20\sigma\sqrt{\log n}+20\sigma\sqrt{n}\|\bm{u}\|_{\infty}}{\lambda^{\star}}+\frac{1}{2}\big\|\bm{u}-\bm{u}^{(l)}\big\|_{2},
\end{align*}
%
provided that $40\sigma\sqrt{n}\leq \lambda^{\star}$.
Rearranging terms and taking the union bound, we demonstrate that with probability at least $1- O(n^{-8})$,
%
\begin{align}
\big\|\bm{u}-\bm{u}^{(l)}\big\|_{2} & \leq\frac{40\sigma\sqrt{\log n}+40\sigma\sqrt{n}\|\bm{u}\|_{\infty}}{\lambda^{\star}} , \qquad 1\leq l\leq n.
	\label{eq:u-ul-proximity-final-denoising}
\end{align}
%

\medskip
\begin{proof}[Proof of the relation \eqref{eq:noise-small-gap}] Apply the triangle inequality to see that
\begin{align*}
\|\bm{M}-\bm{M}^{(l)}\|&\leq\|\bm{M}-\bm{M}^{\star}\|+\|\bm{M}^{\star}-\bm{M}^{(l)}\|=\|\bm{E}\|+\|\bm{E}^{(l)}\| \\
&\leq10\sigma\sqrt{n}\leq\lambda^{\star}/10.
\end{align*}
%
Here, the equality arises from the definition of $\bm{M}$ and $\bm{M}^{\star}$,
the penultimate inequality uses (\ref{eq:L2-perturbation-summary-denoising}), while the last inequality holds as long
as $100\sigma\sqrt{n}\leq\lambda^{\star}$. This together with \eqref{eq:lambda-l-minus-lambda-j-LB1} establishes \eqref{eq:noise-small-gap}.
%
\end{proof}


\subsubsection{Step 2: analyzing leave-one-out estimates}

We now turn attention to bounding the size of the $l$-th entry $u_l^{(l)}$ of $\bm{u}^{(l)}$. Since $\lambda^{(l)}> 0$, by \eqref{eq:ul_l-expression-denoising}, we have 
%together with the fact $\lambda^{(l)}> 0$ allows us to obtain
%
\begin{align*}
u_{l}^{(l)}-u_{l}^{\star} & %=u_{l}^{\star}\Big(\frac{\lambda^{\star}}{\lambda^{(l)}}\bm{u}^
%{\star\top}\bm{u}^{(l)}-1\Big)
=u_{l}^{\star}\Big(\frac{\lambda^{\star}}{\lambda^{(l)}}\bm{u}^{\star\top}\bm{u}^{(l)}-\bm{u}^{\star\top}\bm{u}^{\star}\Big)\\
 & =u_{l}^{\star}\Big(\frac{\lambda^{\star}-\lambda^{(l)}}{\lambda^{(l)}}\bm{u}^{\star\top}\bm{u}^{(l)}\Big)+u_{l}^{\star}\bm{u}^{\star\top}\big(\bm{u}^{(l)}-\bm{u}^{\star}\big).
\end{align*}
%
%where the first line relies on \eqref{eq:ul-l-expression-denoising}
%and the fact  $\|\bm{u}^{\star}\|_2=1$. Therefore,
The triangle inequality and the Cauchy-Schwarz inequality then give
%
\begin{align}
\big|u_{l}^{(l)}-u_{l}^{\star}\big|
& \leq\big|u_{l}^{\star}\big|\cdot\frac{\big|\lambda^{\star}-\lambda^{(l)}\big|}{\big|\lambda^{(l)}\big|} \cdot \|\bm{u}^{\star}\|_{2} \cdot \|\bm{u}^{(l)}\|_{2}
	+ \big|u_{l}^{\star}\big|\cdot\|\bm{u}^{\star}\|_{2} \cdot \big\|\bm{u}^{(l)}-\bm{u}^{\star}\big\|_{2} \notag\\
 & \leq\big|u_{l}^{\star}\big|\cdot\frac{10\sigma\sqrt{n}}{\lambda^{\star}}+\big|u_{l}^{\star}\big|\cdot\frac{10\sigma\sqrt{n}}{\lambda^{\star}} \notag\\
 & \leq  \frac{20\sigma\sqrt{n}}{\lambda^{\star}} \big\|\bm{u}^{\star}\big\|_{\infty}.
	\label{eq:ul-ulstar-bound-denoising-rank1}
\end{align}
%
Here, the second line holds due to Condition (\ref{eq:L2-perturbation-summary-denoising}) and the fact $\big|\lambda^{(l)}\big|\geq\lambda^{\star}/2$; see~(\ref{eq:lambda-l-minus-lambda-j-LB1}).



\subsubsection{Step 3: putting all pieces together}


Putting \eqref{eq:u-ul-proximity-final-denoising} and \eqref{eq:ul-ulstar-bound-denoising-rank1} together, we arrive at
%
\begin{align}
\big\|\bm{u}-\bm{u}^{\star}\big\|_{\infty} & =\max_{l}\big|u_{l}-u_{l}^{\star}\big|\leq\max_{l}\Big\{\big|u_{l}^{(l)}-u_{l}^{\star}\big|+\big\|\bm{u}-\bm{u}^{(l)}\big\|_{2}\Big\} \notag\\
 & \leq \frac{20\sigma\sqrt{n}}{\lambda^{\star}} \big\|\bm{u}^{\star}\big\|_{\infty}
	+ \frac{40\sigma\sqrt{\log n}+40\sigma\sqrt{n}\|\bm{u}\|_{\infty}}{\lambda^{\star}}.
	\label{eq:u-ustar-inf-first-bound}
\end{align}
%
The above upper bound, however, involves the term $\|\bm{u}\|_{\infty}$, which can further be bounded by
$\|\bm{u}^{\star}\|_{\infty} + \|\bm{u}-\bm{u}^{\star} \|_{\infty}$.  Substituting this into \eqref{eq:u-ustar-inf-first-bound}, we have
\begin{align*}
\big\|\bm{u}-\bm{u}^{\star}\big\|_{\infty}
%\eqref{eq:u-ustar-inf-first-bound}
%	& \leq\frac{20\sigma\sqrt{n} \big\|\bm{u}^{\star}\big\|_{\infty}}{\lambda^{\star}}
%	+ \frac{40\sigma\sqrt{\log n}  + 40 \sigma \sqrt{n} \big( \|\bm{u}^{\star}\|_{\infty} + \|\bm{u}-\bm{u}^{\star} \|_{\infty} \big)}{\lambda^{\star}}\\
	& \leq \frac{40\sigma\sqrt{\log n} + 60 \sigma \sqrt{n}\, \|\bm{u}^{\star}\|_{\infty} }{\lambda^{\star}}+\frac{1}{2}\big\|\bm{u}-\bm{u}^{\star}\big\|_{\infty},
\end{align*}
%
provided that $80\sigma\sqrt{n}\leq\lambda^{\star}$. Rearranging terms yields
\begin{align*}
	 \big\|\bm{u}-\bm{u}^{\star}\big\|_{\infty}
	 &\leq \frac{ 80\sigma\sqrt{\log n} + 120 \sigma \sqrt{n}\, \|\bm{u}^{\star}\|_{\infty} }{\lambda^{\star}}
	 % \leq  \frac{200\sigma\sqrt{\log n}\, ( \sqrt{n} \|\bm{u}^{\star}\|_{\infty} )}{\lambda^{\star}} \\
	 = \frac{80\sigma\sqrt{\log n} + 120 \sigma \sqrt{\mu} }{\lambda^{\star}} ,
\end{align*}
%
where the last identity results from the definition \eqref{defn:incoherence-rank-1-denoising} of $\mu$.

%






\section{$\ell_{2,\infty}$ eigenspace perturbation under independent noise}
\label{sec:setup-general-theory-independent}



The appealing entrywise behavior of the eigenvector estimator in Section~\ref{sec:rank-1-denoising-LOO} hints at the promising performance of spectral methods for broader contexts. In this section, we set out to develop a more general framework about $\ell_{2,\infty}$ eigenspace perturbation  that covers a wide spectrum of scenarios.



\subsection{Setup and notation}\label{sec:setup-general-theory}


\paragraph{Ground truth.}
Consider a rank-$r$ symmetric matrix $\bm{M}^{\star}\in \mathbb{R}^{n\times n}$  with eigenvectors $\bm{u}_1^{\star},\cdots,\bm{u}_n^{\star}$ and associated eigenvalues $\lambda_1^{\star},\cdots,\lambda_n^{\star}$ obeying
%
\begin{align} \label{eq:assumption-lambda-r-positive-setup}
	|\lambda_1^{\star}| \geq |\lambda_2^{\star}| \geq \cdots \geq |\lambda_r^{\star}| >0
	\quad \text{and} \quad
	\lambda_{r+1}^{\star} = \cdots = \lambda_n^{\star}=0.
\end{align}
%
We shall write the eigendecomposition $\bm{M}^{\star}=\bm{U}^{\star}\bm{\Lambda}^{\star}\bm{U}^{\star\top}$ as usual, where $\bm{\Lambda}^{\star} \coloneqq \mathsf{diag}([\lambda_1^{\star},\cdots,\lambda_r^{\star}])$ and $\bm{U}^{\star}\coloneqq   [\bm{u}_1^{\star},\cdots, \bm{u}_r^{\star}]\in \mathbb{R}^{n\times r}$. Denote the condition number of $\bm{M}^{\star}$ as
%
\begin{align}
	\label{eq:defn-kappa}
	\kappa \coloneqq {|\lambda_1^{\star}|} \,/\, {|\lambda_r^{\star}|}.
\end{align}
%
Akin to Definition~\ref{assump:mc-incoherence}, the  incoherence parameter of $\bm{M}^{\star}$ is defined as
%
\begin{align}
	%\|\bm{U}^{\star}\|_{2,\infty} \coloneqq \max_i \big\| \bm{e}_i^{\top}\bm{U}^{\star} \big\|_{2} =  \sqrt{\frac{\mu r}{n}} .
	\mu \coloneqq \frac{n\|\bm{U}^{\star}\|_{2,\infty}^{2}}{r},
	\label{eq:defn-Ustar-incoherence}
\end{align}
%
%This incoherence parameter, which is crucial in our theoretical development,
%captures how well the energy of $\bm{U}^{\star}$ is spread out across all rows. It is readily seen from the fact $n^{-1/2}\|\bm{U}^{\star}\|_{\mathrm{F}} \leq \|\bm{U}^{\star}\|_{2,\infty}\leq \|\bm{U}^{\star}\|$ that
a parameter that captures how well the energy of $\bm{U}^{\star}$ is spread out across all rows and that obeys (see Remark~\ref{remark:range-mu-MC})
%
\begin{align}
	1 \leq \mu \leq n/r .
	\label{eq:mu-constraint-general}
\end{align}



\paragraph{Observed data.} What we observe is a corrupted version
%
\begin{align}
	\bm{M}=\bm{M}^{\star}+\bm{E} \in \mathbb{R}^{n\times n},
	\label{eq:M-noisy-copy-general}
\end{align}
%
where $\bm{E}$ is a symmetric noise matrix. We denote by $\{\lambda_i\}_{1\leq i\leq n}$ the set of eigenvalues of $\bm{M}$ obeying
%
\begin{align}
	|\lambda_1| \geq |\lambda_2| \geq \cdots \geq |\lambda_n|,
	\label{eq:lambda-1-n-ordering-magnitude}
\end{align}
%
and let $\bm{u}_i$ be the eigenvector of $\bm{M}$ associated with $\lambda_i$. We shall introduce the diagonal matrix $\bm{\Lambda}\in \mathbb{R}^{r\times r}$ as $\bm{\Lambda} \coloneqq \mathsf{diag}([\lambda_1,\cdots,\lambda_r])$.





\paragraph{Noise assumptions.}
This section aims to cover a fairly broad class of scenarios of independent noise.
In particular, the noise matrix considered herein is assumed to satisfy the mild conditions listed below.
%
\begin{assumption}
	\label{assumption-noise-general}
	%[Noise assumptions]
	The entries in the lower triangular part of $\bm{E}=[E_{i,j}]_{1\leq i,j\leq n}$ are independently generated obeying
%
\begin{align}
	\mathbb{E}[E_{i,j}] = 0, \quad \mathbb{E}[E_{i,j}^2] =: \sigma_{i,j}^2 \leq \sigma^2, \quad |E_{i,j}|\leq B, \quad \text{for all }i\geq j.
	\label{eq:noise-E-condition-general}
\end{align}
%
	In particular, $\sigma^2$ is taken to be the smallest choice satisfying \eqref{eq:noise-E-condition-general}. 
	Further, it is assumed that
	%there is some positive quantity $c_{\mathsf{b}}=O(1)$ such that
	%
	\begin{align}
		c_{\mathsf{b}} \coloneqq \frac{B}{  \sigma  \sqrt{n/(\mu\log n)} } = O(1).
		\label{eq:assumption-B-sigma}
	\end{align}
\end{assumption}
%

We emphasize that both $\sigma$ and $B$ are quantities that are allowed to scale with $n$. When $\mu$ is not too large, Condition \eqref{eq:assumption-B-sigma}
allows the maximum magnitude $B$ of each noisy entry to be substantially larger than the typical size $\sigma$.








\paragraph{Goal and algorithm.}
We seek to estimate $\bm{U}^{\star}$ based on $\bm{M}$. Towards this, a simple spectral  method computes the matrix $\bm{U}=[\bm{u}_1,\cdots,\bm{u}_r]\in \mathbb{R}^{n\times r}$ that comprises the top-$r$ leading eigenvectors of $\bm{M}$.






\subsection{$\ell_{2,\infty}$ and $\ell_{\infty}$ theoretical guarantees}
\label{sec:L2inf-Linf-eigen-theory}


The leave-one-out argument introduced before, when properly strengthened, enables powerful  $\ell_{2,\infty}$  performance guarantees for the spectral estimate $\bm{U}$, which concern row-wise perturbation of the eigenspace. Before continuing, we remind the readers of the global rotation ambiguity, namely, in general we cannot expect $\bm{U}$  to be close to $\bm{U}^{\star}$ unless suitable global rotation is taken into account. In light of this, we introduce the following notation that helps identify a proper rotation matrix.
%
\begin{definition} For any matrix $\bm{Z}$ with SVD $\bm{Z}=\bm{U}_{Z}\bm{\Sigma}_Z\bm{V}_{Z}^{\top}$ (where $\bm{U}_Z$ and $\bm{V}_Z$ represent respectively the left and right singular matrices of $\bm{Z}$, and $\bm{\Sigma}_Z$ is a diagonal matrix composed of the singular values), define
%
\begin{align}
	\mathsf{sgn}(\bm{Z}) \coloneqq \bm{U}_{Z}  \bm{V}_{Z}^{\top}
	\label{eq:defn-sgn-Z}
\end{align}
%
to be the matrix sign function of $\bm{Z}$.
\end{definition}
%
\begin{remark}
The matrix sign function is commonly encountered when aligning two matrices---classically known as the orthogonal Procrustes problem \citep{schonemann1966generalized}.
Consider any two matrices $\widehat{\bB},\bB\in \mathbb{R}^{n\times r}$ with $r\leq n$. Among all rotation matrices,
the one that best aligns  $\widehat \bB$ with $\bB$ is precisely $ \mathsf{sgn}(\widehat \bB^\top \bB)$ (see, e.g., \citep[Appendix D.2.1]{ma2017implicit}), namely,
$$
	\mathsf{sgn}(\widehat \bB^\top \bB) = \underset{\bm{O}\,\in \mathcal{O}^{r\times r}}{\arg\min} ~\| \widehat \bB {\bm O} - \bB\|_{\mathrm{F}}^2.
$$
\end{remark}
%
With this definition in place, we are ready to state an $\ell_{2,\infty}$ theory adapted from \citet{abbe2020entrywise}.
Compared to the original development in \citet{abbe2020entrywise}, 
the theorem and its proof provided herein are more streamlined versions tailored to the current setting. 
%
\begin{theorem}
\label{thm:UsgnH-Ustar-MUstar-general}
	Consider the settings and assumptions in Section~\ref{sec:setup-general-theory}. Define $\bm{H} \coloneqq \bm{U}^{\top}\bm{U}^{\star}$. With probability exceeding $1-O(n^{-5})$, one has
%
\begin{subequations}
\label{eq:UsgnH-Ustar-MUstar-bound-theorem-general}
\begin{align}
& \big\|\bm{U}\mathsf{sgn}(\bm{H})-\bm{U}^{\star}\big\|_{2,\infty}  \lesssim \frac{\sigma\kappa\sqrt{\mu r}+\sigma\sqrt{r\log n}}{ |\lambda_{r}^{\star} | },
\label{eq:UsgnH-Ustar-bound-theorem-general}\\
& \big\|\bm{U}\mathsf{sgn}(\bm{H})-\bm{M}\bm{U}^{\star}\big(\bm{\Lambda}^{\star}\big)^{-1}\big\|_{2,\infty} \notag\\
& \qquad\qquad \lesssim \frac{\sigma\kappa\sqrt{\mu r}}{ |\lambda_{r}^{\star} | }+ \frac{\sigma^{2}\sqrt{rn\log n}+\sigma B\sqrt{\mu r\log^{3}n}}{\big(\lambda_{r}^{\star}\big)^{2}}, 
	%\frac{\sigma\kappa\sqrt{\mu r}}{ |\lambda_{r}^{\star} | }+\frac{\sigma^{2}\sqrt{rn\log n}\,(1+c_{\mathsf{b}}\sqrt{\log n})}{\big(\lambda_{r}^{\star}\big)^{2}},
\label{eq:UsgnH-MUstar-bound-theorem-general}
\end{align}
\end{subequations}
%
provided that $\sigma \sqrt{n \log n}  \leq c_{\sigma} |\lambda_r^{\star}|$ for some sufficiently small constant $c_{\sigma}>0$.
\end{theorem}
%
The proof of this theorem can be found in Section~\ref{sec:proof-thm:UsgnH-Ustar-MUstar-general}. 
Note that under the assumption in the theorem, the bound on the right-hand side of 
\eqref{eq:UsgnH-MUstar-bound-theorem-general} is no larger than the one on the right-hand side of \eqref{eq:UsgnH-Ustar-bound-theorem-general}. 
In fact, \eqref{eq:UsgnH-MUstar-bound-theorem-general} could indeed be tighter than \eqref{eq:UsgnH-Ustar-bound-theorem-general} in some important scenarios like community recovery (see Section~\ref{sec:community-detection-linf}).


The $\ell_{2,\infty}$ perturbation theory in Theorem~\ref{thm:UsgnH-Ustar-MUstar-general} accommodates a broad family of noise matrices with independent entries.
In the sequel, we take a moment to interpret several key messages conveyed by this result.


\paragraph{De-localization of estimation errors.} For simplicity, let us concentrate on the case where $\mu,\kappa =O(1)$.   Note that the Davis-Kahan theorem introduced previously results in the following $\ell_2$ estimation guarantees (to be detailed in Section~\ref{sec:preliminary-facts-Linf-general})
%
\begin{align}
	\mathsf{dist}_{\mathrm{F}}(\bm{U},\bm{U}^{\star})  \leq \sqrt{r} \,\mathsf{dist}(\bm{U},\bm{U}^{\star}) \lesssim  \frac{\sigma \sqrt{nr}}{|\lambda_r^{\star}|}.
	\label{eq:simple-Euclidean-error-general}
\end{align}
%
In comparison, the $\ell_{2,\infty}$ bound derived in  Theorem~\ref{thm:UsgnH-Ustar-MUstar-general} simplifies to
%
\begin{align}
	\min_{\bm{R}\,\in \mathcal{O}^{r\times r}}\big\|\bm{U} \bm{R}-\bm{U}^{\star}\big\|_{2,\infty}
	 \leq \big\|\bm{U}\mathsf{sgn}(\bm{H})-\bm{U}^{\star}\big\|_{2,\infty}  \lesssim  \frac{\sigma \sqrt{r\log n}}{ |\lambda_{r}^{\star}| }
	 \label{eq:Linf-eigen-bound-simpler}
\end{align}
%
under the condition $\mu,\kappa = O(1)$,  which is about $O(\sqrt{n/\log n})$ times smaller than the Euclidean error bound \eqref{eq:simple-Euclidean-error-general}.
This implies that the estimation error of $\bm{U}$ is fairly de-localized and spread out across all rows.


\paragraph{First-order approximation.}
Informally, Theorem~\ref{thm:UsgnH-Ustar-MUstar-general} (and its analysis) unveils the goodness of the first-order approximation
%
\begin{align}
	\bm{U}\mathsf{sgn}(\bm{H}) \approx \bm{M}\bm{U}^{\star}(\bm{\Lambda}^{\star})^{-1} = \bm{U}^{\star} + \bm{E}\bm{U}^{\star}(\bm{\Lambda}^{\star})^{-1}
\end{align}
%
uniformly across all rows. An implication of Theorem~\ref{thm:UsgnH-Ustar-MUstar-general} is that $\bm{U}$ might be closer to the first-order approximation $\bm{M}\bm{U}^{\star}(\bm{\Lambda}^{\star})^{-1}$ than to the ground truth $\bm{U}^{\star}$ (namely, the upper bound on the right-hand side of  \eqref{eq:UsgnH-MUstar-bound-theorem-general} is smaller than the bound on the right-hand side of \eqref{eq:UsgnH-Ustar-bound-theorem-general} under the stated conditions). In principle, the linear term $\bm{E}\bm{U}^{\star}(\bm{\Lambda}^{\star})^{-1}$ can be viewed as a correction term that helps improve the approximation fidelity.
As we shall see momentarily in Section~\ref{sec:community-detection-linf}, this subtle difference leads to sharper performance guarantees in applications like community recovery.




\paragraph{Entrywise estimation errors.} There is no shortage of applications where one cares more about the fine-grained estimation accuracy of the matrix rather than that of the low-rank factors. Fortunately, the $\ell_{2,\infty}$ theory derived in  Theorem~\ref{thm:UsgnH-Ustar-MUstar-general} in turn enables entrywise performance guarantees when estimating the matrix $\bm{M}^{\star}$. This is summarized in the following corollary, with the proof deferred to Section~\ref{sec:proof-cor-entrywise-error}.
%
\begin{corollary}
	\label{cor:entrywise-error-general}
	Consider the settings and assumptions in Section~\ref{sec:setup-general-theory}, and assume  further that
	$\sigma \kappa \sqrt{n\log n} \leq c_1 |\lambda_r^{\star}|$ for some sufficiently small constant $c_1>0$. Then with probability at least $1-O(n^{-5})$, one has
	%
	\begin{align}
		\big\|\bm{U}\bm{\Lambda}\bm{U}^{\top}-\bm{M}^{\star}\big\|_{\infty} & \lesssim\sigma\kappa^{2}\mu r\sqrt{\frac{\log n}{n}}.
	\end{align}
	%
\end{corollary}
%
Once again, it is instrumental to explain the result by specializing it to the simpler regime where $\kappa, \mu, r =O(1)$.  In this case, the finding of Corollary~\ref{cor:entrywise-error-general} reduces to
%
\begin{align}
	\big\|\bm{U}\bm{\Lambda}\bm{U}^{\top}-\bm{M}^{\star}\big\|_{\infty} & \lesssim\sigma  \sqrt{\frac{\log n}{n}}.
	\label{eq:ULambdaU-entrywise-simple}
\end{align}
%
In comparison, the Euclidean error of this spectral estimate satisfies (which follows by combining \eqref{eq:ULambdaU-Mstar-norm-denoising} with \eqref{eq:X-spectral-norm-iid-special-ramon} and \eqref{eq:assumption-B-sigma})
%
\begin{align}
	\big\|\bm{U}\bm{\Lambda}\bm{U}^{\top}-\bm{M}^{\star}\big\|_{\mathrm{F}} \leq 2\sqrt{2} \|\bm{E}\| \lesssim \sigma\sqrt{n}
\end{align}
%
with high probability, which is on the order of $n/\sqrt{\log n}$ times larger than the entrywise error bound \eqref{eq:ULambdaU-entrywise-simple}. In other words, the energy of the estimation error of the unknown matrix is also dispersed more or less across all matrix entries, a message that cannot be derived from classical matrix perturbation theory alone.



\paragraph{Leave-one-out analysis.}
As alluded to previously, the proof of Theorem~\ref{thm:UsgnH-Ustar-MUstar-general} relies heavily upon the leave-one-out analysis framework to decouple delicate statistical dependency. While the core idea bears close resemblance to the exposition in Section~\ref{sec:LOO-analysis-rank1-denoising}, implementing this idea rigorously for the general case requires considerably more effort. We defer a complete proof to Section~\ref{sec:proof-thm:UsgnH-Ustar-MUstar-general}.





\section{$\ell_{2,\infty}$ singular subspace perturbation under independent noise}
\label{sec:general-L2inf-SVD-perturbation}

The general theory presented in Section~\ref{sec:setup-general-theory-independent} applies only to symmetric matrices. It is not uncommon, however, to encounter scenarios where the matrix of interest $\bm{M}^{\star}$ is  asymmetric.
This motivates the need of extending the $\ell_{2,\infty}$ perturbation theory to accommodate more general matrices, which is the main content of the current section.



\subsection{Setup and notation}\label{sec:general-theory-setup-asymm}
\paragraph{Ground truth.}

Consider an unknown rank-$r$ matrix $\bm{M}^{\star}\in\mathbb{R}^{n_{1}\times n_{2}}$. Let
 $\bm{M}^{\star}=\bm{U}^{\star}\bm{\Sigma}^{\star}\bm{V}^{\star\top}$
represent the SVD of $\bm{M}^{\star}$,
where $\bm{U}^{\star}\in\mathbb{R}^{n_{1}\times r}$ (resp.~$\bm{V}^{\star}\in\mathbb{R}^{n_{2}\times r}$)
entails the top-$r$ left (resp.~right) singular vectors of $\bm{M}^{\star}$,
and $\bm{\Sigma}^{\star}=\mathsf{diag}([\sigma_{1}^{\star},\sigma_{2}^{\star},\cdots,\sigma_{r}^{\star}])$
is formed by the (nonzero) singular values of $\bm{M}^{\star}$. We  arrange the singular values $\{\sigma_i^{\star}\}$ in descending order (i.e., $\sigma_{1}^{\star}\geq \sigma_{2}^{\star} \geq \cdots \geq \sigma_{r}^{\star} > 0$). 
%
%and define the condition number of the matrix $\bm{M}^{\star}$ as $\kappa\coloneqq  {\sigma_{1}^{\star}}/{\sigma_{r}^{\star}}$.
%
% In accordance with (\ref{eq:defn-Ustar-incoherence}),

\paragraph{Key parameters.}
As usual,  $\mu$ stands for the incoherence parameter  of  $\bm{M}^{\star}$ (see Definition~\ref{assump:mc-incoherence}),
and  the condition number of the matrix $\bm{M}^{\star}$ is defined as $\kappa\coloneqq  {\sigma_{1}^{\star}}/{\sigma_{r}^{\star}}$.
Without loss of generality, it is assumed that $$n_1 \leq n_2$$ and we set $n \coloneqq n_1 + n_2$.



\paragraph{Observations and noise assumptions. }
%
Assume we have access to corrupted observations of $\bm{M}^{\star}$ as follows:
%
\[
	\bm{M}=\bm{M}^{\star}+\bm{E}\in\mathbb{R}^{n_{1}\times n_{2}},
\]
%
where $\bm{E}=[E_{i,j}]$ stands for a  noise or perturbation matrix. We impose the following conditions on  $\bm{E}$, which is a natural
adaptation of Assumption~\ref{assumption-noise-general} to the general case and covers a diverse array of scenarios.

\begin{assumption}
\label{assump:noise-general-rect}
The entries of $\bm{E}$ are independently generated obeying
%
\begin{align}
	\mathbb{E}[E_{i,j}] = 0, \quad \mathbb{E}[E_{i,j}^2]\leq \sigma^2, \quad |E_{i,j}|\leq B \quad \text{for all }i, j.
\end{align}
%
	Further,  assume that
	\begin{align}
		c_{\mathsf{b}} \coloneqq \frac{B}{\sigma  \sqrt{n_1/(\mu\log n)}} = O(1).
	\end{align}\end{assumption}



\paragraph{Goal and algorithm.}
Again, we aim at estimating $\bm{U}^{\star}$ and $\bm{V}^{\star}$, based
on the observation $\bm{M}$, using a spectral method. Specifically, let
%
\begin{equation}
\bm{M}=\left[\begin{array}{cc}
\bm{U} & \bm{U}_{\perp}\end{array}\right]\left[\begin{array}{cc}
\bm{\Sigma}\\
 & \bm{\Sigma}_{\perp}
\end{array}\right]\left[\begin{array}{c}
\bm{U}^{\top}\\
\bm{V}^{\top}
\end{array}\right]\label{eq:general-M-SVD}
\end{equation}
be the SVD of $\bm{M}$, in which $\bm{U}\bm{\Sigma}\bm{V}^{\top}$
is the rank-$r$ SVD (i.e., the singular values in $\bm{\Sigma}\coloneqq\mathsf{diag}([\sigma_{1},\cdots,\sigma_{r}])$
are larger than those in $\bm{\Sigma}_{\perp}$). The spectral method then
	deploys ($\bm{U},\bm{V}$) as an estimate of ($\bm{U}^{\star},\bm{V}^{\star}$).


	

\subsection{$\ell_{2,\infty}$ and $\ell_{\infty}$ theoretical guarantees}
\label{sec:theory-Linf-SVD-general}

We now present a theorem that generalizes Theorem~\ref{thm:UsgnH-Ustar-MUstar-general} and Corollary~\ref{cor:entrywise-error-general} to accommodate general (asymmetric and possibly rectangular) matrices. This can be accomplished via a standard ``symmetric dilation'' trick; the details can be found in Section~\ref{sec:proof-inf-rect}.  
% Yuxin: moved to Section 4.3.1
%Recall the definition of incoherent parameter $\mu$ in definition~\ref{assump:mc-incoherence} and condition number $\kappa$ that generalizes \eqref{eq:defn-kappa}



\begin{theorem}
\label{thm:2-inf-asymm}
%
Consider the settings and
assumptions in Section~\ref{sec:general-theory-setup-asymm}, and
define $\bm{H}_{\bm{U}}\coloneqq\bm{U}^{\top}\bm{U}^{\star}$ and
$\bm{H}_{\bm{V}}\coloneqq\bm{V}^{\top}\bm{V}^{\star}$. With probability
at least $1-O(n^{-5})$, one has
%
\begin{align}
	& \max\Big\{\|\bm{U}\mathsf{sgn}(\bm{H}_{\bm{U}})-\bm{U}^{\star}\|_{2,\infty},\, \|\bm{V}\mathsf{sgn}(\bm{H}_{\bm{V}})-\bm{V}^{\star}\|_{2,\infty} \Big\} \nonumber\\
	& \qquad\qquad\qquad\qquad \lesssim\frac{\sigma\sqrt{ r}\big(\kappa\sqrt{\frac{n_{2}}{n_{1}}\mu}+\sqrt{\log n}\big)}{\sigma_{r}^{\star}},\label{eq:main-result-2-infty-general-asymm}
\end{align}
%
provided that $\sigma\sqrt{n\log n}\leq c_{1}\sigma_{r}^{\star}$
for some sufficiently small constant $c_{1}>0$. In addition, if $\sigma\kappa\sqrt{n\log n}\leq c_{2}\sigma_{r}^{\star}$
for some small enough constant $c_{2}>0$, then the following holds with
probability at least $1-O(n^{-5})$:
%
\begin{equation}
	\|\bm{U}\bm{\Sigma}\bm{V}^{\top}-\bm{M}^{\star}\|_{\infty}\lesssim\sigma\kappa^{2}\mu r\sqrt{\frac{(n_2/n_1) \log n}{n_{1}}}.\label{eq:main-result-infty-general-asymm}
\end{equation}
%
\end{theorem}



The messages conveyed in Theorem~\ref{thm:2-inf-asymm} largely parallel those in Theorem~\ref{thm:UsgnH-Ustar-MUstar-general} and Corollary~\ref{cor:entrywise-error-general}.
For simplicity, let us discuss the implications when $\kappa, \mu, r = O(1)$ and $n_1 \asymp n_2$ (i.e., the aspect ratio of the matrix is $n_2 / n_1 = O(1)$).
In this scenario, Theorem~\ref{thm:2-inf-asymm} implies that
%
\begin{align*}
\|\bm{U}\mathsf{sgn}(\bm{H}_{\bm{U}})-\bm{U}^{\star}\|_{2,\infty}+\|\bm{V}\mathsf{sgn}(\bm{H}_{\bm{V}})-\bm{V}^{\star}\|_{2,\infty} & \lesssim\frac{\sigma\sqrt{\log n}}{\sigma_{r}^{\star}},\\
\|\bm{U}\bm{\Sigma}\bm{V}^{\top}-\bm{M}^{\star}\|_{\infty} & \lesssim\sigma  \sqrt{\frac{\log n}{n}},
\end{align*}
%
both of which bear close similarities to our previous observations \eqref{eq:Linf-eigen-bound-simpler} and \eqref{eq:ULambdaU-entrywise-simple}. Akin to our discussions in Section~\ref{sec:L2inf-Linf-eigen-theory}, these findings tell us that the  singular subspace estimation errors (resp.~the matrix estimation errors)  are fairly spread out across all rows of the singular subspace (resp.~all entries of the matrix).





\section{Application: Entrywise guarantees for matrix completion}

To illustrate the utility of the fine-grained perturbation theory presented in previous sections, let us revisit the problem of matrix completion introduced in Section~\ref{sec:mc-l2} and apply our refined theory.


As a recap, the spectral method proposed for matrix completion proceeds by computing the best rank-$r$ approximation $\bm{U}\bm{\Sigma}\bm{V}^{\top}$ of the rescaled data matrix $\bm{M}=p^{-1}\mathcal{P}_{\Omega}(\bm{M}^{\star})$, where $p$ is the probability of each entry being observed, and $\mathcal{P}_{\Omega}(\cdot)$ denotes the Euclidean projection onto the set of matrices supported on the sampling set $\Omega$.
This time, we seek to characterize the $\ell_{2,\infty}$ error when estimating the true singular subspaces $\bm{U}^{\star}$ and $\bm{V}^{\star}$,
as well as the $\ell_{\infty}$ error when estimating the unknown matrix $\bm{M}^{\star}$, as stated below. As before, we set $n \coloneqq n_1 + n_2$.




\begin{theorem}
\label{thm:mc-inf}
%
Consider the settings and assumptions in Section~\ref{sec:problem-formulation-MC}, and
define $\bm{H}_{\bm{U}}\coloneqq\bm{U}^{\top}\bm{U}^{\star}$ and
$\bm{H}_{\bm{V}}\coloneqq\bm{V}^{\top}\bm{V}^{\star}$.
Suppose that $n_1\leq n_2$ and $n_{1}p\geq C\kappa^{4}\mu^{2}r^{2}\log n$ for some
sufficiently large constant $C>0$. Then with probability greater
than $1-O(n^{-5})$, we have
%
\begin{subequations}
%
\begin{align}
\max\{\|\bm{U}\mathsf{sgn}(\bm{H}_{\bm{U}})-\bm{U}^{\star}\|_{2,\infty},\,\,&\|\bm{V}\mathsf{sgn}(\bm{H}_{\bm{V}})-\bm{V}^{\star}\|_{2,\infty}\} \nonumber\\& \leq\kappa^{2}\sqrt{\frac{\mu^{3}r^{3}\log n}{n_{1}^{2}p}};
\label{eq:mc-inf-claim-1}\\
\|\bm{U}\bm{\Sigma}\bm{V}^{\top}-\bm{M}^{\star}\|_{\infty} & \lesssim\kappa^{2}\mu^{2}r^{2}\sqrt{\frac{\log n}{n_{1}^{3}p}}\|\bm{M}^{\star}\|.
\label{eq:mc-inf-claim-2}
\end{align}
%
\end{subequations}
%
\end{theorem}
%
\begin{proof}
	Recall our notation $\bm{E}= \bm{M}  -\bm{M}^{\star}= p^{-1}\mathcal{P}_{\Omega}(\bm{M}^{\star})-\bm{M}^{\star}$.
 It is straightforward to check
that $\bm{E}$ satisfies Assumption~\ref{assump:noise-general-rect} with
%
\begin{equation}
\sigma^{2}\coloneqq\frac{\|\bm{M}^{\star}\|_{\infty}^{2}}{p},\qquad\text{and}\qquad B\coloneqq\frac{\|\bm{M}^{\star}\|_{\infty}}{p}. \label{eq:mc-inf-noise}
\end{equation}
%
In addition, from the relation $B= c_{\mathsf{b}}\sigma\sqrt{n_1/({\mu}\log n)}$, it is seen that $c_{\mathsf{b}}= O(1)$
holds as long as $n_{1}p\gtrsim\mu\log n$.
With these preparations in place, the claims in Theorem~\ref{thm:mc-inf} follow
directly from Theorem~\ref{thm:2-inf-asymm} and the bound (\ref{eq:mc-entry-upper}) on $\|\bm{M}^{\star}\|_{\infty}$
(and hence on $\sigma$).
%
\end{proof}

%\begin{figure}
%\begin{center}
%	\includegraphics[width=0.55\textwidth]{figures/mc_inf_new.pdf} 
%\end{center}	
%	\caption{Numerical performance of spectral methods for matrix completion. Set $n=500$, $r=10$, and the underlying low-rank matrix $\bm{M}^{\star}$ as a product of two random orthonormal matrices. We vary the sampling probability $p$ from 0.25 to 0.9 with an equal space of 0.05. For each $p$, we draw 20 independent observation patterns and run the spectral method to recover the data matrix. The Euclidean error $\|\widehat{\bm{M}}-\bm{M}^{\star}\|_{\mathrm{F}}/\|\bm{M}^{\star}\|_{\mathrm{F}}$, the spectral norm error $\|\widehat{\bm{M}}-\bm{M}^{\star}\|/\|\bm{M}^{\star}\|$, and the $\ell_{\infty}$ norm error $\|\widehat{\bm{M}}-\bm{M}^{\star}\|_{\infty}/\|\bm{M}^{\star}\|_{\infty}$ are reported when averaged over 20 independent trials, where $\widehat{\bm{M}}=\bm{U}\bm{\Sigma}\bm{V}^{\top}$. 
%	\label{fig:mc-inf}}
%\end{figure}


In what follows, we compare the $\ell_{2,\infty}$ and $\ell_{2}$ performance
guarantees derived in the above theorem with those $\ell_{2}$ guarantees presented
in Section~\ref{sec:mc-l2-theory}; see Figure~\ref{fig:matrix-completion-motivation} in Section~\ref{sec:motivating_examples} for empirical performances. For the sake of brevity, we shall concentrate on the case where $\mu,\kappa,r=O(1)$.




\begin{itemize}

\item {\em $\ell_{2,\infty}$ singular space perturbation bounds. }
Comparing the $\ell_{2,\infty}$ performance guarantee (\ref{eq:mc-inf-claim-1}) with  Theorem~\ref{thm:mc-l2-subspace}, one sees that the $\ell_{2,\infty}$ perturbation bounds could be an order of $\sqrt{n_1}$ times smaller than the $\ell_{2}$ counterpart, showcasing the de-localization effect of the errors of the spectral estimates $\bm{U}$ and $\bm{V}$.

\item {\em Entrywise matrix estimation bounds. }
Furthermore, the entrywise error (\ref{eq:mc-inf-claim-2}) is about an order of $n_1$ times smaller than the corresponding Euclidean error predicted  in Theorem~\ref{thm:mc-l2-matrix}. This indicates that no entry in the resulting matrix estimate suffers from an error significantly higher than the average entrywise error.

\end{itemize}





\section{Application: Exact community recovery}
\label{sec:community-detection-linf}


Another application that benefits remarkably from the fine-grained eigenvector perturbation theory is community recovery. This section reexplores the stochastic block model studied in Section~\ref{sec:community-detection}, and develops significantly enhanced theoretical support for spectral clustering. 




\subsection{Performance guarantees: Exact recovery}
\label{sec:performance-exact-recovery}


The focus of this section is simultaneous recovery of the community memberships of all vertices, 
which is termed {\em exact recovery} or {\em strong consistency} in the community detection literature \citep{abbe2017community}. This imposes a much stronger requirement than the {\em weak consistency} studied in Section~\ref{sec:theory-community-detection}. 

For the sake of conciseness, the theorem below concentrates on the challenging regime where $p,q\asymp \frac{\log n}{n}$,  corresponding to the lowest possible edge densities that allow for exact recovery. This is because, if $p < \frac{\log n}{n}$, then with high probability, one can find isolated vertices that are not connected with any edge in the graph \citep{durrett2007random}; hence, there will be absolutely no means to infer the community membership of these isolated vertices.  The theoretical guarantee is as follows. 
%
\begin{theorem}
	\label{thm:community-recovery-linf}
	Fix any constant $\varepsilon>0$, and consider the setting of Section~\ref{sec:setup-community-detection}. Suppose  $p=\frac{\alpha \log n}{n}$ and $q=\frac{\beta \log n}{n}$ for some sufficiently large constants $\alpha > \beta >0$.\footnote{In the current proof, the constants $\alpha$, $\beta$ might depend on the fixed choice of $\varepsilon$. Encouragingly, this restriction can be lifted; see~\citet{abbe2020entrywise} for details.} In addition, assume that
	%
	\begin{align}
		\big( \sqrt{p} -\sqrt{q}\big)^2 \geq   2\left( 1+ \varepsilon  \right) \frac{\log n}{n} .
		\label{eq:H-pq-lower-bound-theorem}
	\end{align}
	%
	With probability  $1-o(1)$, the spectral method in Section~\ref{sec:algorithm-community-detection} yields
	%
	\[
		x_{i}=x_{i}^{\star}~~\text{for all }1\leq i\leq n, \quad~\text{or}\quad~ x_{i}=-x_{i}^{\star}~~\text{for all }1\leq i\leq n.
	\]
\end{theorem}
%
This theorem, which first appeared in \citet{abbe2020entrywise},  identifies a sufficient recovery condition in terms of the edge densities. 
The result substantially strengthens the $\ell_2$-based theory in Section~\ref{sec:theory-community-detection}, 
uncovering the capability of the spectral method in achieving not merely almost exact recovery in the average sense, 
but more appealingly,  exact community recovery that ensures correct labels of all vertices.  

A natural question arises as to whether the recovery condition \eqref{eq:H-pq-lower-bound-theorem} is improvable via more sophisticated algorithms. Answering this question requires information-theoretic thinking, that is, how to characterize a fundamental threshold---in terms of the difference of edge densities---below which exact recovery is deemed  infeasible. As has been demonstrated in \citet{abbe2014exact,mossel2015consistency,hajek2015achieving}, no algorithm whatsoever is able to achieve exact community recovery  if
%
\begin{equation}
	\big( \sqrt{p} -\sqrt{q}\big)^2 \leq  2(1-\varepsilon) \frac{\log n}{n} 
	\label{eq:IT-lower-bound-CD}
\end{equation}
%
for any constant $\varepsilon >0$. 
This fundamental lower bound, in conjunction with Theorem~\ref{thm:community-recovery-linf}, reveals a sharp phase transition behind the performance of the spectral method. In particular, its optimality is guaranteed all the way down to the information-theoretic threshold; see Figure~\ref{fig:sbm-inf} for numerical evidence.  

\begin{figure}[!t]
\begin{center}
	\includegraphics[width=0.55\textwidth]{figures/SBM-rec-new.pdf} 
\end{center}	
	\caption{Phase transition of spectral methods for exact community recovery. Set $n=300$, $p=(a \log n) / n$, and $q = (b \log n) / n$. We vary $a,b$ from 0 to 8 with an equal space of 0.1. For each configuration of $a \geq b$, we conduct 100 Monte Carlo trials and report the empirical success rate for recovering the entire community structure correctly. 
	%White areas indicate complete success, while dark ones indicate complete failure. 
	The empirical phase transition occurs near the information-theoretic threshold (see \eqref{eq:IT-lower-bound-CD}).
	\label{fig:sbm-inf}}
\end{figure}



Given that the above information-theoretic threshold is specified in terms of $( \sqrt{p} -\sqrt{q})^2$,
the reader might naturally wonder what the operational meaning of this quantity is. As it turns out, this metric is a sort of distance measure between the two edge probability distributions under consideration. In truth, in the setting of Theorem~\ref{thm:community-recovery-linf}, this metric is intimately related to the squared Hellinger distance between two Bernoulli distributions.
%
\begin{definition}[Squared Hellinger distance]
%
\label{defn:Hellinger-distance}
Consider two distributions $P$ and $Q$ over a finite alphabet $\mathcal{Y}$. The squared Hellinger distance $\mathsf{H}^2(P\,\|\,Q)$ between $P$ and $Q$  is defined as follows	%
\begin{equation}
\mathsf{H}^2(P\,\|\,Q)\coloneqq\frac{1}{2}\sum\nolimits_{y\in \mathcal{Y}}\Big(\sqrt{P(y)}-\sqrt{Q(y)}\Big)^{2}.
\label{eq:defn-Hellinger-PQ}
\end{equation}
%
\end{definition}
%
In particular, consider the squared Hellinger distance between two Bernoulli distributions of interest $\mathsf{Bern}(p)$ and $\mathsf{Bern}(q)$, where we denote by $\mathsf{Bern}(p)$ the Bernoulli distribution with mean $p$. It is seen that \citep{chen2016information}
%
\begin{align*}
	\mathsf{H}^2\big( \mathsf{Bern}(p),   \mathsf{Bern}(q) \big) 
	&\coloneqq \frac{1}{2} \big( \sqrt{p} -\sqrt{q}\big)^2 + \frac{1}{2} \big( \sqrt{1-p} -\sqrt{1-q}\big)^2 \nonumber\\
	&= (1+o(1)) \frac{1}{2} \big( \sqrt{p} -\sqrt{q}\big)^2,
\end{align*}
%
when $p=o(1)$ and $q=o(1)$.\footnote{To justify this approximation, the following calculation suffices: \[ \sqrt{1-q}-\sqrt{1-p}=\frac{p-q}{\sqrt{1-p}+\sqrt{1-q}}=(1+o(1))\big(\sqrt{p}-\sqrt{q}\big)\big(\sqrt{p}+\sqrt{q}\big)=o\big(\sqrt{p}-\sqrt{q}\big).\]} 
The phase transition phenomenon identified in \eqref{eq:H-pq-lower-bound-theorem} and \eqref{eq:IT-lower-bound-CD} can then be alternatively described as
%
\begin{align*}
\text{\text{spectral method works}}\quad & \text{if }\mathsf{H}^{2}\big(\mathsf{Bern}(p),\mathsf{Bern}(q)\big)\geq(1+\varepsilon)\frac{\log n}{n}\\
\text{no algorithm works}\quad & \text{if }\mathsf{H}^{2}\big(\mathsf{Bern}(p),\mathsf{Bern}(q)\big)\leq(1-\varepsilon)\frac{\log n}{n}
\end{align*}
%
for an arbitrary small constant $\varepsilon >0$. 


\subsection{Proof of Theorem~\ref{thm:community-recovery-linf}} 

We now turn to the proof of Theorem~\ref{thm:community-recovery-linf}. 
Without loss of generality, suppose that $x_i^{\star}=1$ for all $1\leq i\leq n/2$  and $x_i^{\star}=-1$ for all $i> n/2$, so that 
$\bm{u}^{\star} =\frac{1}{\sqrt{n}}  {\small
\left[\begin{array}{c}
\bm{1}_{n/2}\\
-\bm{1}_{n/2}
\end{array}\right] }$. 



Recalling the matrix $\bM$ given in \eqref{eq:M-data-matrix-SBM} and its mean $\bM^{\star}$ in \eqref{eq:M-data-matrix-expectation-SBM}, 
one can immediately see that $\kappa=\mu=r=1$ for $\bM^{\star}$ in this application. 
Theorem \ref{thm:UsgnH-Ustar-MUstar-general} (cf.~\eqref{eq:UsgnH-MUstar-bound-theorem-general}) readily implies 
the existence of some $z\in\{1,-1\}$ such that
%
\begin{align}
 & \Big\| z\bm{u}-\frac{1}{\lambda^{\star}}\bm{M}\bm{u}^{\star}\Big\|_{\infty}
\lesssim\frac{\sigma}{|\lambda^{\star}|}+\frac{\sigma^{2}\sqrt{n\log n}+\sigma B \,{\log^{3/2}n}}{(\lambda^{\star})^{2}} 
\label{eq:implication-Thm-symmetry-1st-CD}
\end{align}
%
with probability at least $1-O(n^{-5})$. Additionally, it has already been explained in Section~\ref{sec:community-detection} that  
%
\[
	 B=1,\quad\sigma^{2}\leq\max\{p,q\}=p, \quad \text{and} \quad \lambda^{\star}={n(p-q)}/{2}.
\]
%
Substitution into \eqref{eq:implication-Thm-symmetry-1st-CD} reveals that
%
\begin{align}
	\big\| z\lambda^{\star}\bm{u}-\bm{M}\bm{u}^{\star} \big\|_{\infty} & \lesssim\sigma+\frac{\sigma^{2}\sqrt{n\log n}}{\lambda^{\star}}+\frac{\sigma B\,\log^{3/2}n}{\lambda^{\star}}\nonumber \\
 	& \leq  C \Big( \sqrt{p}+\frac{p\sqrt{\log n}}{\sqrt{n}(p-q)}+\frac{\sqrt{p}\log^{3/2}n}{n(p-q)} \Big)
	\label{eq:z-lambda-u-MU-bound-CD}
\end{align}
%
holds for some universal constant $C>0$.
As a result, a crucial step boils down to controlling $\bm{M}\bm{u}^{\star}$ in an entrywise manner: each element is a difference between two independent random binomial random variables and  is accomplished through the following lemma. 
%
\begin{lemma}
	\label{lemma:M-ustar-lower-bound-CD}
	Suppose that
%
\begin{align}
	\mathsf{H}_{p,q}^{2} \coloneqq \big( \sqrt{p} -\sqrt{q}\big)^2 \geq\left( 1+ \varepsilon  \right) \frac{2\log n}{n} 
	\label{eq:H-pq-lower-bound-lemma}
\end{align}
% 
for some quantity $\varepsilon>0$.  Let $\varepsilon_0 \coloneqq \frac{\varepsilon\log n}{\sqrt{n}\log\frac{p(1-q)}{q(1-p)}}-\frac{1}{\sqrt{n}}$. 
Then with probability exceeding $1-n^{-\varepsilon/2}$, one has
%
\begin{align*}
	 \bm{M}_{l,\cdot}\bm{u}^{\star}   \geq \varepsilon_0 \,\,\,\text{for all }l\leq \frac{n}{2},
	\quad\text{and}\quad
	   \bm{M}_{l,\cdot} \bm{u}^{\star}  
	\leq - \varepsilon_0 \,\,\, \text{for all } l > \frac{n}{2}. 
	%\label{eq:M-ustar-bound-CD-lemma}
\end{align*}
%
\end{lemma}
%

We now return to analyze the entrywise behavior of $\bm{u}$. Note that
%
\begin{align*}
(z\lambda^{\star})u_{l} & \geq\bm{M}_{l,\cdot}\bm{u}^{\star}-\big| z\lambda^{\star}u_{l}-\bm{M}_{l,\cdot}\bm{u}^{\star}\big|  \quad\text{for all } l\leq n/2;\\
(z\lambda^{\star})u_{l} & \leq\bm{M}_{l,\cdot}\bm{u}^{\star}+\big| z\lambda^{\star}u_{l}-\bm{M}_{l,\cdot}\bm{u}^{\star}\big|  \quad\text{for all }l>n/2.
\end{align*}
%
This together with \eqref{eq:z-lambda-u-MU-bound-CD}, Lemma~\ref{lemma:M-ustar-lower-bound-CD} and $\lambda^{\star}>0$ yields that if
%
\begin{align}
	\frac{\varepsilon\log n}{\sqrt{n}\log\frac{p(1-q)}{q(1-p)}}  > \frac{1}{\sqrt{n}} 
	+  C \Big( \sqrt{p} + \frac{p\sqrt{\log n}}{\sqrt{n}(p-q)} + \frac{\sqrt{p}\log^{3/2}n}{n(p-q)} \Big), 
	\label{eq:epsilon-condition-CD-inf-123}
\end{align}
%
then it follows that
%
\[
zu_{l}>0\quad\text{for all }1\leq l\leq\frac{n}{2},\quad\text{and}\quad zu_{l}<0\quad\text{for all }l>\frac{n}{2},
\]
%
thus guaranteeing exact community recovery once the rounding procedure (based on the sign) is applied.  


To finish up, it remains to validate Condition~\eqref{eq:epsilon-condition-CD-inf-123}. 
Fixing $\varepsilon>0$ to be a constant, we make the following observations. 
%
\begin{itemize}
	\item From the assumptions $\varepsilon \asymp 1$ and $p,q\asymp \frac{\log n}{n}$ (or $\alpha,\beta \asymp 1$), one has
	%
	\[
		\frac{\varepsilon\log n}{\sqrt{n}\log\frac{p(1-q)}{q(1-p)}}=\frac{\varepsilon\log n}{\sqrt{n}\log\frac{(1+o(1))\alpha}{\beta}}
		\asymp \frac{\log n}{\sqrt{n}}\gg \sqrt{p}+\frac{1}{\sqrt{n}}.
	\]
	\item Turning to the term $\frac{p\sqrt{\log n}}{\sqrt{n}(p-q)}$,   we  observe that 
	%
	\begin{align}
		& \log\frac{p(1-q)}{q(1-p)}=\log\Big(1+\frac{p-q}{q(1-p)}\Big)\leq\frac{p-q}{q(1-p)}\leq\frac{2(\alpha-\beta)}{\beta}, \notag\\
		& \qquad\quad \Longrightarrow \quad 
		\frac{\varepsilon\log n}{\sqrt{n}\log\frac{p(1-q)}{q(1-p)}}\geq\frac{\varepsilon\log n}{2\sqrt{n}}\,\frac{\beta}{\alpha-\beta},
		\label{eq:LB-CD-12345}
	\end{align}
	%
	where in the last inequality of the  first line we have used the assumption $p=o(1)$ and hence $1/(1-p) \leq 2$. 
	Given that $\varepsilon, \alpha,\beta \asymp 1$, it is guaranteed that
	%
	\[
		\eqref{eq:LB-CD-12345} \gg \frac{\alpha\sqrt{\log n}}{\sqrt{n}(\alpha-\beta)}=\frac{p\sqrt{\log n}}{\sqrt{n}(p-q)} .
	\]
	% 
	\item We then move on to the term $\frac{\sqrt{p}\log^{3/2}n}{n(p-q)}$. 
	If $\alpha / \beta \leq 2$ and  $\beta \geq\frac{200C^{2}}{\varepsilon^{2}} \geq \frac{100C^{2}\alpha/\beta}{\varepsilon^{2}}$, then one  has $\beta \geq \frac{10C\sqrt{\alpha}}{\varepsilon}$ and hence
	%
	\[
		\eqref{eq:LB-CD-12345} \geq\frac{5C\sqrt{\alpha}\log n}{\sqrt{n}(\alpha-\beta)}=\frac{5C\sqrt{p}\log^{3/2}n}{n(p-q)} .
	\]
	%
	 In addition, if $\alpha / \beta > 2$, then it follows that
	%
	\begin{equation}
		\frac{\sqrt{p}\log^{3/2}n}{n(p-q)}=\frac{\sqrt{\alpha}\log n}{\sqrt{n}(\alpha-\beta)}\leq\frac{2\sqrt{\alpha}\log n}{\sqrt{n}\alpha}=\frac{2\log n}{\sqrt{n\alpha}},
		\label{eq:LB-CD-5678}
	\end{equation}
	%
	where the inequality holds true since $\alpha - \beta > \alpha - \alpha/2 = \alpha / 2$. 
	Using the basic inequality $\log x \leq \sqrt{x}$ further leads  to
	%
	\[
		\frac{\varepsilon\log n}{\sqrt{n}\log\frac{p(1-q)}{q(1-p)}}\geq\frac{\varepsilon\log n}{\sqrt{n}\log\frac{2\alpha}{\beta}}
		\geq\frac{\varepsilon\sqrt{\beta}\log n}{\sqrt{n}\sqrt{2\alpha}}\geq\frac{5C\sqrt{p}\log^{3/2}n}{n(p-q)}.
	\]
	%
	Here, the first inequality holds since $p,q=o(1)$ and hence $\frac{1-q}{1-p}\leq 2$, 
	whereas the last relation relies on \eqref{eq:LB-CD-5678} and holds with the proviso that $\beta\geq200C^{2}/\varepsilon^{2}$. 
	
 
\end{itemize}
%
The above calculations taken collectively establish Condition~\eqref{eq:epsilon-condition-CD-inf-123} under the assumptions of Theorem~\ref{thm:community-recovery-linf}, thus concluding the proof. 



\begin{remark}
	It is worth pointing out that the bound \eqref{eq:UsgnH-Ustar-bound-theorem-general} in Theorem \ref{thm:UsgnH-Ustar-MUstar-general} is not sufficiently tight when establishing this result. 
	Instead, one needs to resort to the more refined bound \eqref{eq:UsgnH-MUstar-bound-theorem-general} in Theorem \ref{thm:UsgnH-Ustar-MUstar-general}, 
	which allows us to sharpen the error bound by explicitly accounting for the first-order error term $(\bm{M}-\bm{M}^{\star})\bm{u}^{\star}$. 
\end{remark}




 


\subsection{Proof of auxiliary lemmas}
\label{sec:proof-auxiliary-CD}

Before embarking on the proof of Lemma~\ref{lemma:M-ustar-lower-bound-CD}, we first record non-asymptotic tail bounds concerning log-likelihood ratios and a sum of Bernoulli random variables, which make apparent the role of the squared Hellinger distance \citep{tsybakov2009nonparm}.
%
\begin{lemma}
	\label{lemma:LLR-Hellinger}
	Consider two distributions $P$ and $Q$ over a finite alphabet $\mathcal{Y}$, and suppose that $P(y)\neq 0$ for all $y\in \mathcal{Y}$. Generate an independent sequence $\{y_i\}_{1\leq i\leq n}$ obeying $y_i \sim P$. Then for any $\zeta\in \mathbb{R}$ one has
	%
\begin{equation}
\mathbb{P}\left\{ \sum_{i=1}^{n}\log\frac{Q(y_{i})}{P(y_{i})}\geq-n\zeta\right\} \leq\exp\left(-n\Big[\mathsf{H}^2(P\,\|\,Q)-\frac{\zeta}{2}\Big]\right),
\label{eq:prob-log-PQ-exp-Hel}
\end{equation}
%
where $\mathsf{H}^2(P\,\|\,Q)$ is the squared Hellinger distance between $P$ and $Q$ defined in \eqref{eq:defn-Hellinger-PQ}.
\end{lemma}
%
%
\begin{lemma}
\label{lem:sum-Bernoulli-CD}
%
Consider two sequences of independent random variables
%
\[
z_{i}\sim\mathsf{Bern}(p),\qquad w_{i}\sim\mathsf{Bern}(q),\qquad1\leq i\leq n,
\]
%
and suppose that $p>q$. For any $\xi\in \mathbb{R}$, it follows that
%
\begin{align*}
\mathbb{P}\left\{ \sum_{i=1}^{n}(z_{i}-w_i) \leq n\xi\right\}  & \leq\exp\left(-n\Big[\mathsf{H}_{p,q}^{2}-\frac{\xi}{2}\log\frac{p(1-q)}{q(1-p)}\Big]\right) ,
\end{align*}
%
where $\mathsf{H}_{p,q}^{2}  \coloneqq  \big(\sqrt{p}-\sqrt{q} \, \big)^{2}$.
%
%\begin{align}
%	.
%	\label{eq:defn-Hpq}
%\end{align}
% 
\end{lemma}
%
In what follows, we first establish Lemmas~\ref{lemma:LLR-Hellinger} and \ref{lem:sum-Bernoulli-CD}, and then return to prove Lemma~\ref{lemma:M-ustar-lower-bound-CD}. 



\paragraph{Proof of Lemma~\ref{lemma:LLR-Hellinger}.}
Apply the Chernoff bound to yield 
%
\begin{align}
 & \mathbb{P}\left\{ \sum_{i=1}^{n}\log\frac{Q(y_{i})}{P(y_{i})}\geq-n\zeta\right\} {\leq}\frac{\mathbb{E}_{y_{i}\sim P}\left[\exp\left(\frac{1}{2}\sum_{i=1}^{n}\log\frac{Q(y_{i})}{P(y_{i})}\right)\right]}{\exp(-n\zeta/2)} \notag\\
 & \qquad \qquad\overset{(\mathrm{i})}{=}
 %\frac{\prod_{i=1}^{n}\mathbb{E}_{y_{i}\sim P}\left[\exp\left(\frac{1}{2}\log\frac{Q(y_{i})}{P(y_{i})}\right)\right]}{\exp(-n\zeta/2)}
 %= 
 \frac{\left(\mathbb{E}_{y\sim P}\Big[\left(\frac{Q(y)}{P(y)}\right)^{1/2}\Big]\right)^{n}}{\exp(-n\zeta/2)},
\label{eq:prob-log-PQ-UB}
\end{align}
%
where 
%(i) results from the Chernoff bound, and (ii) 
(i) holds due to the i.i.d.~assumption of the $y_{i}$'s.  In addition,
%It is seen from the definition of $P$ and $Q$ that
%
\begin{align}
\mathbb{E}_{y\sim P}\left[\left(\frac{Q(y)}{P(y)}\right)^{1/2}\right] & =\sum_{y}P(y)\left(\frac{Q(y)}{P(y)}\right)^{1/2}=\sum_{y}\sqrt{P(y)Q(y)} \notag\\
 & =1-\frac{1}{2}\sum_{y}\left(P(y)+Q(y)-2\sqrt{P(y)Q(y)}\right) \notag\\
  & =1-\frac{1}{2}\sum_{y}\left( \sqrt{P(y)} - \sqrt{Q(y)} \right)^2 \notag\\
 & =1-\mathsf{H}^2(P\,\|\,Q)\leq  \exp \big(-\mathsf{H}^2(P\,\|\,Q) \big),
 \label{eq:P-Q-Hellinger-bound}
\end{align}
%
where the second line follows since $\sum_yP(y) = \sum_yQ(y)=1$, and  the last line uses the definition \eqref{eq:defn-Hellinger-PQ} and the elementary inequality $1-x\leq \exp(-x)$. 
Substituting \eqref{eq:P-Q-Hellinger-bound} into \eqref{eq:prob-log-PQ-UB} concludes the proof. 







\paragraph{Proof of Lemma~\ref{lem:sum-Bernoulli-CD}.}

Set  $y_i\coloneqq z_i - w_i$. The proof is built upon a mapping between $\sum_{i=1}^{n}y_{i}$ and a certain log-likelihood ratio. 
Specifically, let us introduce two distributions $P$ and $Q$ supported on $\{1,0,-1\}$:
%
\[
P(x)=\begin{cases}
p(1-q),\quad & \text{if }x=1,\\
q(1-p), & \text{if }x=-1,\\
pq+(1-p)(1-q), \quad & \text{if }x=0,
\end{cases} 
%\qquad \text{and}
\]
% 
%
\[
Q(x)=\begin{cases}
q(1-p),\quad & \text{if }x=1,\\
p(1-q), & \text{if }x=-1,\\
pq+(1-p)(1-q), \quad & \text{if }x=0.
\end{cases}
\]
%
Apparently, $P$ (resp.~$Q$) corresponds to the distribution of $y_i$ (resp.~$-y_i$).  A key observation is that
%
\begin{align*}
\sum_{i=1}^{n}\log\frac{Q(y_{i})}{P(y_{i})} & =\sum_{i=1}^{n}\left\{ \mathbbm1\{y_{i}=1\}\log\frac{q(1-p)}{p(1-q)}+\mathbbm1\{y_{i}=-1\}\log\frac{p(1-q)}{q(1-p)}\right\} \\
 & =\sum_{i=1}^{n}y_{i}\log\frac{q(1-p)}{p(1-q)} ,
\end{align*}
%
which relies on the fact that $y_i$ is supported on $\{1,0,-1\}$. 
Recognizing that $\log\frac{p(1-q)}{q(1-p)}>0$ holds as long as
$p>q$ (since $q(1-p)<p(1-q)$), we can further derive
%
\begin{align*}
\mathbb{P}\left\{ \sum_{i=1}^{n}y_{i}\leq n\xi\right\}  & =\mathbb{P}\left\{ \frac{1}{\log\frac{p(1-q)}{q(1-p)}}\sum_{i=1}^{n}\log\frac{Q(y_{i})}{P(y_{i})}\geq-n\xi\right\} \\
 & \leq\exp\left(-n\Big[\mathsf{H}^2(P\,\|\,Q)-\frac{\xi}{2}\log\frac{p(1-q)}{q(1-p)}\Big]\right),
\end{align*}
%
where the last inequality comes from Lemma~\ref{lemma:LLR-Hellinger}. 
From the constructions of $P$ and $Q$ and the definition \eqref{eq:defn-Hellinger-PQ} of $\mathsf{H}^2(P\,\|\,Q)$, it is easily seen that 
%
\begin{align*}
\mathsf{H}^{2}(P\,\|\,Q) & =\Big(\sqrt{p(1-q)}-\sqrt{q(1-p)}\,\Big)^{2}
	= \frac{\big(p-q\big)^{2}}{\big(\sqrt{p(1-q)}+\sqrt{q(1-p)}\,\big)^{2}}\\
 & \geq\frac{\big(p-q\big)^{2}}{\big(\sqrt{p}+\sqrt{q}\,\big)^{2}}=\big(\sqrt{p}-\sqrt{q}\big)^{2} ,
\end{align*}
%
thus concluding the proof. 






\paragraph{Proof of Lemma~\ref{lemma:M-ustar-lower-bound-CD}.}

Let us start by looking at the first entry of $\bm{M}\bm{u}^{\star}$. 
It is seen from the construction \eqref{eq:M-data-matrix-SBM} that
%
\begin{align}
\bm{M}_{1,\cdot}\bm{u}^{\star}=\bm{A}_{1,\cdot}\bm{u}^{\star}-\frac{p+q}{2}\big(\bm{1}^{\top}\bm{u}^{\star}\big)\bm{1}+pu_{1}^{\star}
\geq \bm{A}_{1,\cdot}\bm{u}^{\star} ,
%+ O\Big(\frac{1}{\sqrt{n}}\Big),
	\label{eq:connection-M1u-A1u-CD}
\end{align}
%
where we have used the fact that $\bm{1}^{\top}\bm{u}^{\star}=0$  and $u_1^{\star}>0$. 
The expression  $\bm{u}^{\star} =\frac{1}{\sqrt{n}}  {\small
\left[\begin{array}{c}
\bm{1}_{n/2}\\
-\bm{1}_{n/2}
\end{array}\right] }$ admits the following decomposition
%
\begin{align}
	\bm{A}_{1,\cdot}\bm{u}^{\star}
%=\frac{1}{\sqrt{n}}\sum_{i=2}^{n/2}\big(A_{1,i}-A_{1,i+n/2}\big)+\frac{1}{\sqrt{n}}A_{1,n/2+1}
	= \frac{1}{\sqrt{n}}\sum\nolimits_{i=1}^{n/2}\big(A_{1,i}-A_{1,i+n/2}\big) , 
	\label{eq:A1ustar-decompose-CD}
\end{align}
%
which can be controlled via Lemma~\ref{lem:sum-Bernoulli-CD}. 
%Here and throughout, we denote by $\mathsf{Bern}(p)$  the Bernoulli distribution with mean $p$.



Observe that $A_{1,i} \sim \mathsf{Bern}(p)$ for all $1< i\leq n/2$ and $A_{1,i} \sim \mathsf{Bern}(q)$ otherwise. Using the definitions of $z_i$ and $w_i$ in Lemma~\ref{lem:sum-Bernoulli-CD}, we obtain
%
\begin{align}
 & \mathbb{P}\left\{ \sum\nolimits_{i=1}^{n/2}\big(A_{1,i}-A_{1,i+n/2}\big)\leq\frac{n\zeta}{2}-1\right\}  
	\leq\mathbb{P}\left\{ \sum\nolimits_{i=1}^{n/2}\big(z_{i}-w_{i}\big)\leq\frac{n\zeta}{2}\right\}  \notag\\
	& \qquad \leq\exp\left(-\frac{n}{2}\Big[\mathsf{H}_{p,q}^{2}-\frac{\zeta}{2}\log\frac{p(1-q)}{q(1-p)}\Big]\right) \leq \frac{1}{n^{1+\delta}}
	\label{eq:A-sum-complicated-CD}
\end{align}
%
for some $\delta >0$, where the first inequality follows since $A_{1,1}=0\leq z_1+1$ (so that $A_{1,1}-A_{1,n/2+1}$ is stochastically dominated by $z_1 - w_1$), and the last inequality holds as long as
%
\begin{equation}
	\mathsf{H}_{p,q}^{2}-\frac{\zeta}{2}\log\frac{p(1-q)}{q(1-p)}
	\geq\frac{2(1+\delta)\log n}{n} ,
	\label{eq:H-pq-lower-bound-123}
\end{equation}
which we shall ensure at the end of the proof. Substituting \eqref{eq:A-sum-complicated-CD} into \eqref{eq:A1ustar-decompose-CD} and \eqref{eq:connection-M1u-A1u-CD} yields
\begin{align*}
\mathbb{P}\left\{ \bm{M}_{1,\cdot}\bm{u}^{\star}\leq \frac{n\zeta-2}{2\sqrt{n}} \right\}  
%& \leq\mathbb{P}\left\{ \bm{A}_{1,\cdot}\bm{u}^{\star}\leq\varepsilon_{0}-1/\sqrt{n}\right\} \\
 & \leq \mathbb{P}\left\{ \bm{A}_{1,\cdot}\bm{u}^{\star}\leq\frac{n\zeta-2}{2\sqrt{n}} \right\} \leq\frac{1}{n^{1+\delta}}.
\end{align*}
%
%which combined with \eqref{eq:connection-M1u-A1u-CD} leads to
% and the assumption ${1}/{\lambda^{\star}} = o(\varepsilon_0)$ leads to
%
%\[
%	\mathbb{P}\left\{ \bm{M}_{1,\cdot}\bm{u}^{\star} \leq  {\varepsilon_0}  \right\} \leq
%	\mathbb{P}\left\{ \bm{A}_{1,\cdot}\bm{u}^{\star} \leq  {\varepsilon_0}   \right\} 
%	\leq\frac{1}{n^{1+\delta}}.
%\]
%
Repeating the preceding analysis for $ \bm{M}_{l,\cdot}\bm{u}^{\star}$ with other  $l$'s  and taking the union bound, we see that with probability at least $1- n^{-\delta}$,
%
\begin{subequations}
	\label{eq:M-ustar-bound-CD}
\begin{align}
%\begin{split}
% & \frac{1}{\lambda^{\star}}\bm{A}_{l,\cdot}\bm{u}^{\star}\geq\frac{\sqrt{n}\zeta}{4\lambda^{\star}}\ \qquad\,\text{for all }l\leq\frac{n}{2}\\
% & \frac{1}{\lambda^{\star}}\bm{A}_{l,\cdot}\bm{u}^{\star}\leq-\frac{\sqrt{n}\zeta}{4\lambda^{\star}}\ \ \quad \text{for all }l>\frac{n}{2}
%\end{split}
	\bm{M}_{l,\cdot}\bm{u}^{\star} &\geq \frac{n\zeta-2}{2\sqrt{n}}
	%\frac{\sqrt{n}\zeta}{4\lambda^{\star}}
	\quad  &&\text{if } l\leq {n}/{2} \\
	%\quad
	%\text{and}\quad
	\bm{M}_{l,\cdot}\bm{u}^{\star}
	&\leq - \frac{n\zeta-2}{2\sqrt{n}}
	%-\frac{\sqrt{n}\zeta}{4\lambda^{\star}}
	\quad  &&\text{if } l> n/2
	%\big)
\end{align}
\end{subequations}
%
%\begin{equation}
%\begin{array}{cc}
%	& \frac{1}{\lambda^{\star}}\bm{M}_{l,\cdot}\bm{u}^{\star}  \geq\frac{\sqrt{n}\zeta-2}{2\lambda^{\star}},\ \quad~~\text{}l\leq\frac{n}{2}\\
%	& \frac{1}{\lambda^{\star}}\bm{M}_{l,\cdot}\bm{u}^{\star}  \leq-\frac{\sqrt{n}\zeta-2}{2\lambda^{\star}},\ \ ~~\text{}l>\frac{n}{2}
%\end{array}
%\end{equation}
%
hold simultaneously for all $1\leq l\leq n$. 



Finally, it remains to ensure satisfaction of \eqref{eq:H-pq-lower-bound-123}. As it turns out, if the condition \eqref{eq:H-pq-lower-bound-lemma}
holds, then it suffices to take $\delta\leq \varepsilon / 2$ and 
$\zeta=\frac{2\varepsilon\log n}{n\log\frac{p(1-q)}{q(1-p)}}$. This completes the proof.  






\section{Distributional theory and uncertainty quantification}
\label{sec:distribution-theory}


Thus far, we have demonstrated intriguing statistical performance of estimators developed based on spectral methods. 
As one can anticipate, the quality of a spectral estimator is largely affected by the imperfectness of data generating mechanisms (e.g., noise corruption, missing data). 
The uncertainty of the estimator due to these factors would inevitably influence any subsequent decision making based on it. 
Viewed in this light, it is recommended to accompany the estimator in hand with valid measures of uncertainty (or ``confidence''), 
in order to better inform decision makers.  


Take the low-rank matrix estimation problem in Section~\ref{sec:setup-general-theory} for instance:  
an important uncertainty quantification task can be posed as the construction of a valid confidence
interval---based on the spectral estimator---that is likely to cover an unseen entry of the matrix of interest $\bm{M}^{\star}$. 
More precisely, for any location $(i,j)$ and any target coverage level $1-\alpha \in (0,1)$ (e.g., 95\%), we aim to identify a short interval---denoted by $\mathsf{CI}_{i,j}^{1-\alpha}$---based on the spectral estimator such that
%
\begin{equation}
	\mathbb{P}\big( M_{i,j}^{\star} \in \mathsf{CI}_{i,j}^{1-\alpha} \big) \approx 1-\alpha,  
\end{equation}
%
which essentially augments a point estimate into an interval that is guaranteed to cover the unknown with the pre-specified target probability.  
Note that the problem of constructing a valid confidence interval falls within the realm of {\em statistical inference} in the statistics literature,
which constitutes an important step beyond statistical estimation. 
Accomplishing this task in high dimension often calls for a refined statistical reasoning toolbox that offers quantitative distributional characterizations of the estimator. 

%The $\ell_{\infty}$ and $\ell_{2,\infty}$ theory presented in this chapter becomes particularly effective for this purpose.    


%calls for a modern suite
%of statistical reasoning tools, with the aim of providing comprehensive uncertainty evaluations for the estimation
%algorithms in use.
%i
%


%quantitative characterization of 

%statistical validity 



\subsection{Entrywise distributional guarantees}
\label{sec:distribution-general}


As a natural starting point to build confidence intervals, 
we seek to develop comprehensive understanding about the distribution of the spectral estimator. 
In general, obtaining a non-asymptotic yet tractable distributional characterization of a nonconvex estimator like the spectral method 
could be remarkably challenging. 
Fortunately, the $\ell_{\infty}$ and $\ell_{2,\infty}$ perturbation theory introduced previously (e.g., Theorem~\ref{thm:UsgnH-Ustar-MUstar-general}) allows one to make progress for some important scenarios. 



Let us revisit the setting in Section~\ref{sec:setup-general-theory}, and consider the following estimator of the unknown low-rank matrix $\bm{M}^{\star}$:
%
\begin{align}
	\widehat{\bm{M}}=\big[\widehat{M}_{i,j} \big]_{1\leq i,j\leq n}=\bm{U}\bm{\Lambda}\bm{U}^{\top} ,
	\label{eq:estimator-M-hat-distribution}
\end{align}
%
obtained via the spectral method. The aim is to develop tractable distributional guarantees for each entry of $\widehat{\bm{M}}-\bm{M}^{\star}$. 


Towards this end, we first examine whether our previous results shed light on certain distributional properties of $\widehat{\bm{M}}-\bm{M}^{\star}$. 
Informally,  Theorem~\ref{thm:UsgnH-Ustar-MUstar-general} (in particular, \eqref{eq:UsgnH-MUstar-bound-theorem-general}) reveals that
%
\begin{align}
	\bm{U}\mathsf{sgn}(\bm{H}) \approx \bm{M} \bm{U}^{\star} \big(\bm{\Lambda}^{\star}\big)^{-1} 
	= \bm{U}^{\star} + \bm{E} \bm{U}^{\star} \big(\bm{\Lambda}^{\star}\big)^{-1}. 
	\label{eq:U-sgn-H-first-order-approx-discuss}
\end{align}
%
Assuming tightness of this first-order approximation, one further derives 
%
\begin{align}
\bm{U}\bm{\Lambda}\bm{U}^{\top}-\bm{M}^{\star} & \overset{(\mathrm{i})}{\approx}\bm{U}\mathsf{sgn}(\bm{H})\bm{\Lambda}^{\star}\big(\bm{U}\mathsf{sgn}(\bm{H})\big)^{\top}-\bm{U}^{\star}\bm{\Lambda}^{\star}\bm{U}^{\star\top} \notag\\
 & \overset{(\mathrm{ii})}{\approx}\big(\bm{U}\mathsf{sgn}(\bm{H})-\bm{U}^{\star}\big)\bm{\Lambda}^{\star}\bm{U}^{\star\top}+\bm{U}^{\star}\bm{\Lambda}^{\star}\big(\bm{U}\mathsf{sgn}(\bm{H})-\bm{U}^{\star}\big)^{\top} 
 \notag\\
 & \overset{(\mathrm{iii})}{\approx}\bm{E}\bm{U}^{\star}\big(\bm{\Lambda}^{\star}\big)^{-1}\bm{\Lambda}^{\star}\bm{U}^{\star\top}+\bm{U}^{\star}\bm{\Lambda}^{\star}\big(\bm{E}\bm{U}^{\star}\big(\bm{\Lambda}^{\star}\big)^{-1}\big)^{\top} \notag\\
 & =\bm{E}\bm{U}^{\star}\bm{U}^{\star\top}+\bm{U}^{\star}\bm{U}^{\star\top}\bm{E}, 
	\label{eq:approx-ULambdaU-Mstar-discuss}
\end{align}
%
where (i) holds as long as $\mathsf{sgn}(\bm{H})\bm{\Lambda}^{\star} \mathsf{sgn}(\bm{H})^{\top} \approx \bm{\Lambda}$ (which has already been illuminated in the analysis 
of Corollary~\ref{cor:entrywise-error-general} and will be solidified momentarily),  
(ii) is obtained by dropping the higher-order term $\big(\bm{U}\mathsf{sgn}(\bm{H})-\bm{U}^{\star}\big)\bm{\Lambda}^{\star}\big(\bm{U}\mathsf{sgn}(\bm{H})-\bm{U}^{\star}\big)^{\top}$, and (iii) relies upon the approximation \eqref{eq:U-sgn-H-first-order-approx-discuss}.  


Given that \eqref{eq:approx-ULambdaU-Mstar-discuss} is a linear map of the noise matrix $\bm{E}$, 
this essentially forms a first-order approximation of $\widehat{\bm{M}}$,
which in turn enables a tractable distributional theory for $\widehat{\bm{M}}$. 
Observe that each entry of the matrix in \eqref{eq:approx-ULambdaU-Mstar-discuss} 
is a weighted superposition of the independent zero-mean entries of $\bm{E}$.   
Equipped with this observation, some variant of the central limit theorem
suggests that each entry of $\widehat{\bm{M}} - \bm{M}^{\star}$ is approximately zero-mean Gaussian,
as formalized by the theorem below. For notational convenience, we shall define a projection matrix
	%
	\begin{equation}
	%\bm{P}\coloneqq\bm{U}^{\star}\bm{U}^{\star\top},
		\bm{P}^{\star}= \big[ P_{i,j}^{\star} \big]_{1\leq i,j\leq n} \coloneqq \bm{U}^{\star}\bm{U}^{\star\top},
		\label{eq:defn-P-UU-T}
	\end{equation}	
	%
and impose a lower bound requirement on the noise variance: 
%
\begin{equation}
	 \sigma_{\min}^2 \leq \sigma_{i,j}^2 \leq \sigma^2, \qquad 1\leq i,j\leq n. 
	 \label{eq:sigma-min-definition}
\end{equation}
%
\begin{theorem}
\label{thm:matrix-distribution-complete}
%
Suppose that the assumptions of Theorem~\ref{thm:UsgnH-Ustar-MUstar-general} hold. For any $1\leq i,j\leq n$, set
%
\begin{align}
	v_{i,j}^{\star}=
		\begin{cases}
			\sum_{l=1}^{n}\sigma_{i,l}^{2}P^{\star 2}_{l,j}+\sum_{l=1}^{n}P^{\star 2}_{i,l}\sigma_{l,j}^{2}+2\sigma_{i,j}^{2}P^{\star}_{i,i}P^{\star}_{j,j}, & \text{if }i\neq j,\\
			4\sum_{l=1}^{n}\sigma_{i,l}^{2}P^{\star 2}_{l,i}, & \text{if }i=j. 
		\end{cases}
	\label{eq:variance-formula-Mij-lem}
\end{align}
%
Assume that $\sigma/\sigma_{\min}=O(1)$, and that
%
\begin{equation}
	\frac{\big\|\bm{U}_{j,\cdot}^{\star}\big\|_{2}^{2}+\big\|\bm{U}_{i,\cdot}^{\star}\big\|_{2}^{2}}{\big\|\bm{U}^{\star}\big\|_{\mathrm{F}}^{2}}
	\gtrsim \frac{B^{2}\kappa^{2}\mu^{2}r^{2}\log^{2}n}{\sigma^{2}n^{2}}+\frac{\sigma^{2}\mu^{2}r\kappa^{4}\log^{3}n}{(\lambda_{r}^{\star})^{2}}.
	%\frac{B^{2}\mu^{2}r\log n}{n^{2}\sigma^{2}}+\frac{\sigma^{2}\mu^{2}r\kappa^{4}\log^{3}n}{(\lambda_{r}^{\star})^{2}}+\frac{\kappa^{2}\mu^{2}r^{2}\log^{2}n}{n^{2}}.
	\label{eq:Ui-Uj-condition-thm-distribution-1}
\end{equation}
%
Then the estimator \eqref{eq:estimator-M-hat-distribution} obeys
%
\[
	\sup_{z\in\mathbb{R}}\left|\mathbb{P}\left(\widehat{M}_{i,j}-M_{i,j}^{\star}\leq z\sqrt{v_{i,j}^{\star}}\,\right)-\Phi(z)\right|=o(1) ,
\]
where $\Phi(\cdot)$ represents the cumulative density function (CDF) of the standard Gaussian distribution.
% 
\end{theorem}
%
The proof of this theorem is postponed to Section~\ref{sec:proof-thm:matrix-distribution-complete}. 
In a nutshell, Theorem~\ref{thm:matrix-distribution-complete} tells us that $\widehat{\bm{M}}$ is a nearly unbiased estimator of the truth $\bm{M}^{\star}$, as long as the signal strength---as captured by $\|\bm{U}^{\star}_{i,\cdot}\|_2$ and $\|\bm{U}^{\star}_{j,\cdot}\|_2$ when estimating the $(i,j)$-th entry---is sufficiently large (cf.~\eqref{eq:Ui-Uj-condition-thm-distribution-1}).  
The resulting estimation error in each entry is well approximated by a zero-mean Gaussian random variable, 
whose variance can be determined in a tractable fashion.  
As can be easily verified, the variance $v_{i,j}^{\star}$ is precisely the variance of the $(i,j)$-th entry of $\bm{E}\bm{U}^{\star}\bm{U}^{\star\top}+\bm{U}^{\star}\bm{U}^{\star\top}\bm{E}$ (as singled out in \eqref{eq:approx-ULambdaU-Mstar-discuss}). 
%In stark contrast to classical large-sample asymptotic analysis \citep{van2000asymptotic}, 
The above distributional theory is non-asymptotic, 
which lends itself well to high-dimensional applications.  





\subsection{Inference and uncertainty quantification}
\label{sec:UQ-general}



The Gaussian approximation unveiled in Theorem~\ref{thm:matrix-distribution-complete},
which is dictated by a single parameter ${v}_{i,j}^{\star}$, 
paves the way for statistical inference and uncertainty quantification tailored to this model. 
In order to construct a valid confidence interval for each entry of $\bm{M}^{\star}$, 
everything boils down to identifying an estimator that approximates the variance parameter ${v}_{i,j}^{\star}$, ideally in a data-driven yet faithful manner. 

In view of the variance characterization \eqref{eq:variance-formula-Mij-lem}, 
computing ${v}_{i,j}^{\star}$ requires information about both the noise variances $\{\sigma_{i,j}^2\}_{1\leq i,j\leq n}$ and the projection matrix $\bm{P}^{\star}$ (cf.~\eqref{eq:defn-P-UU-T}). However, estimating the noise variances is in general statistically infeasible, given that we only have access to a single observation (i.e., $M_{i,j}$) related to each individual variance $\sigma_{i,j}^2$.  Fortunately, the variance  ${v}_{i,j}^{\star}$ involves only the summation or equivalently the average of these individual variances, whose stochastic errors will be averaged out.  This leads us to the following surrogate
%
\begin{equation}
	\widetilde{v}_{i,j}=\begin{cases}
		\sum_{l=1}^{n}{E}_{i,l}^{2}{P}_{l,j}^{\star 2}+\sum_{l=1}^{n}{P}_{i,l}^{\star 2}{E}_{l,j}^{2}+2{E}_{i,j}^{2}{P}^{\star}_{i,i}{P}_{j,j}^{\star}, & \text{if }i\neq j,\\
	4\sum_{l=1}^{n}{E}_{i,l}^{2}{P}_{l,i}^{\star 2}, & \text{if }i=j, 
\end{cases}
	\label{eq:v-ij-plugin-surrogate-123}
\end{equation}
%
which is clearly an unbiased estimator of ${v}_{i,j}^{\star}$. 
Given the statistical independence of $\{E_{i,j}\}_{i\geq j}$, 
we can expect to have $\widetilde{v}_{i,j}\approx v_{i,j}^{\star}$, owing to the concentration of measure.  



%we propose to replace $\sigma_{i,j}^2$ by the corresponding entry $E_{i,j}^2$, or more realistically, an estimator $\widehat{E}_{i,j}^2$ of $E_{i,j}^2$. 



However, the above surrogate $\widetilde{v}_{i,j}$ remains practically incomputable, 
due to the absence of knowledge about both $\bm{E}$ and $\bm{P}^{\star}$.   
To address this issue, we propose the following plug-in estimator:
%
\begin{equation}
	\widehat{v}_{i,j}=\begin{cases}
	\sum_{l=1}^{n}\widehat{E}_{i,l}^{2}\widehat{P}_{l,j}^{2}+\sum_{l=1}^{n}\widehat{P}_{i,l}^{2}\widehat{E}_{l,j}^{2}+2\widehat{E}_{i,j}^{2}\widehat{P}_{i,i}\widehat{P}_{j,j}, & \text{if }i\neq j,\\
	4\sum_{l=1}^{n}\widehat{E}_{i,l}^{2}\widehat{P}_{l,i}^{2}, & \text{if }i=j, 
\end{cases}
	\label{eq:v-ij-plugin-estimator}
\end{equation}
%
where $\widehat{\bm{E}}=[\widehat{E}_{i,j}]_{1\leq i,j\leq n}$ and $\widehat{\bm{P}}=[\widehat{P}_{i,j}]_{1\leq i,j\leq n}$ stand for estimators of $\bm{E}$ and $\bm{P}^{\star}$, respectively. 
In particular, we employ the following specific estimators of $\bm{E}$ and $\bm{P}^{\star}$, again adopting the plug-in strategy:
%
\begin{subequations}
\label{eq:estimators-E-and-P}
\begin{align}
\widehat{\bm{E}} & \coloneqq\bm{M}-\bm{U}\bm{\Lambda}\bm{U}^{\top},\label{eq:estimator-noise-matrix-E}\\
\widehat{\bm{P}} & \coloneqq\bm{U}\bm{U}^{\top},\label{eq:estimator-projection-P}
\end{align}
\end{subequations}
%
where $\bm{U}$ and $\bm{\Lambda}$ are, as usual, computed via eigendecomposition of $\bm{M}$. 
For a prescribed coverage level $1-\alpha$ (with $0<\alpha <1$), we construct the following confidence interval for the $(i,j)$-th entry of $\bm{M}^{\star}$, 
motivated by the Gaussian approximation in Theorem~\ref{thm:matrix-distribution-complete}: 
%
\begin{align}
	\label{eq:confidence-interval-construction-general}
	\mathsf{CI}_{i,j}^{1-\alpha} \coloneqq \Big[\widehat{M}_{i,j}\pm\Phi^{-1}(1-\alpha/2)\sqrt{\widehat{v}_{i,j}}\,\Big].
\end{align}
%
Here and throughout, for any $b>0$, we let $[a\pm b]$ abbreviate the interval $[a-b,a+b]$, and we use $\Phi^{-1}(\cdot)$ to represent the inverse CDF of the standard Gaussian distribution. 


As encouraging news, the above construction of entrywise confidence intervals is provably valid with high probability, as revealed by the following theorem. The proof is postponed to Section~\ref{sec:proof-thm:CI-general}. 
%
\begin{theorem}
	\label{thm:CI-general}
	Consider the settings and assumptions in Section~\ref{sec:setup-general-theory},  
	%the assumptions of Theorem~\ref{thm:UsgnH-Ustar-MUstar-general} hold.
	and suppose that $\sigma / \sigma_{\min} = O(1)$,  $\kappa^{4}\mu^{2}r^{2}\log n\leq n$ and $\sigma \kappa\sqrt{n\log n}\lesssim |\lambda_{r}^{\star}|$. Consider any $1\leq i,j\leq n$, and assume that
	%
	\begin{align}
		\frac{\big\|\bm{U}_{j,\cdot}^{\star}\big\|_{2}^{2}+\big\|\bm{U}_{i,\cdot}^{\star}\big\|_{2}^{2}}{\big\|\bm{U}^{\star}\big\|_{\mathrm{F}}^{2}}\gtrsim\frac{B\kappa^{2}\mu^{2}r^{2}\log^{2}n}{\sigma n^{3/2}}+\frac{\sigma\mu^{2}r\kappa^{3}\log^{3}n}{|\lambda_{r}^{\star}|\sqrt{n}}.  
		\label{eq:Uj-Ui-lower-bound-thm}
	\end{align}
	%
	For any fixed coverage level $1-\alpha \in (0,1)$, 
	the confidence interval $\mathsf{CI}_{i,j}^{1-\alpha}$ constructed in \eqref{eq:confidence-interval-construction-general} obeys
	%
	\begin{align}
		%\left| 
		\mathbb{P} \Big( M_{i,j}^{\star} \in \mathsf{CI}_{i,j}^{1-\alpha}  
		\Big) 
		= 1-\alpha + o(1).
		% \right| =o(1).
	\end{align}
	%
\end{theorem}
%
%It is noteworthy that: 
Theorem~\ref{thm:CI-general} confirms that the confidence interval proposed above meets the prescribed coverage requirement,
provided that the associated signal strength is not too low (see \eqref{eq:Uj-Ui-lower-bound-thm}). 
In addition to its statistical validity, the proposed procedure enjoys several features that make it practically appealing: 
%
\begin{itemize}
	\item {\em Adaptive to unknown noise levels and distributions.} The above inference procedure is fully data-driven, which does not require prior knowledge about the noise levels or noise distributions. As alluded to previously, it is in general impossible to estimate the noise variance in each entry, and hence a data-driven yet valid approach is of critical value. 

	\item {\em Adaptive to heteroskedastic noise.} Our statistical guarantees hold without relying on homogeneity of noise components. In other words, this inference procedure automatically accommodates heteroskedastic noise, a scenario where the variance of the noise components might vary across different locations.   
\end{itemize}


Careful readers might remark that Theorem~\ref{thm:CI-general} is concerned with statistical inference for a single entry. 
Interestingly, the distributional theory presented in Section~\ref{sec:distribution-general} (see also Lemma~\ref{thm:matrix-distribution} in the proof of Theorem~\ref{thm:matrix-distribution-complete}) might also be instrumental in pursuing simultaneous inference, namely, the problem of constructing a confidence region that simultaneously accounts for more than one unknown entries. We omit such an extension for the sake of conciseness. 


%
%\end{proof}

%%
%\begin{align*}
% & \big\|\mathsf{sgn}(\bm{H})\bm{\Lambda}^{\star}\mathsf{sgn}(\bm{H})^{\top}-\bm{\Lambda}\big\|\\
% & \qquad\leq\big\|\mathsf{sgn}(\bm{H})\bm{\Lambda}^{\star}\mathsf{sgn}(\bm{H})^{\top}-\bm{H}\bm{\Lambda}^{\star}\bm{H}^{\top}\big\|+\big\|\bm{H}\bm{\Lambda}^{\star}\bm{H}^{\top}-\bm{\Lambda}\big\|\\
% & \qquad\lesssim\frac{\kappa\sigma^{2}n}{\lambda_{r}^{\star}}+\sigma\sqrt{r\log n}.
%\end{align*}



%
\section{Matrix completion}\label{sec:mc-l2}

A pressing challenge often encountered in data science applications is estimation and learning in the face of \emph{missing data}.
To elucidate how to tackle this challenge via spectral methods, we delve into the renowned matrix completion problem in this section,
followed by another application called tensor completion in Section~\ref{sec:tensor-completion}.


Imagine that one observes a small subset of the entries in a large unknown matrix
and seeks to fill in all missing entries. An archetypal example is collaborative filtering, where one aims to predict the users' preferences on a collection of products based on partially revealed user-product ratings.   See Figure~\ref{fig:completion} for an illustration.  The problem, often referred to as matrix completion, is apparently ill-posed in general, as there are (much) fewer measurements than the unknowns.

%This is indeed a factor model in Section~\ref{sec:PCA} with missing data.  

Fortunately, if the matrix of interest exhibits certain low-dimensional structure, then reliable recovery becomes feasible.
A commonly encountered example of this kind concerns the case when the target matrix enjoys a low-rank structure. Again, take collaborative filtering for example:  the user-product rating matrix  might be well explained by a relatively small number of latent factors connecting users' preferences with products' attributes, thus resulting in an approximately low-rank matrix.  Motivated by its fundamental importance,  recent years have witnessed a flurry of research activity in studying low-rank matrix completion \citep{candes2009exact, keshavan2010matrix, gross2011recovering};
see \citet{chen2018harnessing,davenport2016overview} for overviews of recent developments.
In the sequel, we present a simple yet effective approach enabled by the spectral method, originally proposed in~\citet{achlioptas2007fast,keshavan2010matrix}.

\begin{figure}
\begin{center}
\includegraphics[width=0.35\textwidth]{figures/matrix_completion}
\end{center}
\caption{Illustration of matrix completion, where each ``$\checkmark$'' stands for an observed entry and each ``$?$'' represents a missing entry.} \label{fig:completion}
\end{figure}


\subsection{Problem formulation and assumptions}
\label{sec:problem-formulation-MC}

Suppose that we are interested in estimating an $n_{1}\times n_{2}$  rank-$r$ matrix  $\bm{M}^{\star} = [{M}^{\star}_{i,j}]_{1\leq i\leq n_1, 1\leq j\leq n_2}$.
Without loss of generality, we assume
$$n_{1}\leq n_{2}.$$  Denote by  $\bm{M}^{\star}=\bm{U}^{\star}\bm{\Sigma}^{\star}\bm{V}^{\star\top}$ the SVD of $\bm{M}^{\star}$, where the columns of $\bm{U}^{\star}\in \mathbb{R}^{n_1\times r}$ (resp.~$\bm{V}^{\star}\in \mathbb{R}^{n_2\times r}$) are the left (resp.~right) singular vectors of $\bm{M}^{\star}$, and $\bm{\Sigma}^{\star}$ is a diagonal matrix whose diagonal entries are the singular values of $\bm{M}^{\star}$. We define the condition number of the matrix $\bm{M}^{\star}$ to be $\kappa \coloneqq \sigma_1(\bM^\star)/\sigma_r(\bM^\star)$.


To capture the presence of missing data, we introduce an index subset $\Omega\subseteq [n_1]\times [n_2]$,
such that each entry ${M}^{\star}_{i,j}$ is observed if and only if $(i,j)\in \Omega$.
The goal is to reconstruct the singular subspaces $\bm{U}^{\star}$ and $\bm{V}^{\star}$, as well as the  full matrix $\bm{M}^{\star}$, based on entries observed over the sampling set $\Omega$.



\paragraph{Random sampling.}
Apparently, not all sampling patterns admit reliable estimation. For instance, if $\Omega$ contains only entries in the top half of the matrix,
then there is in general no hope to predict the bottom half of the matrix. In order to allow for meaningful matrix completion,
this monograph focuses on a natural random observation model commonly adopted in the literature, as formulated below.
%
\begin{assumption}[Random sampling]
\label{assump:mc-bernoulli}
Each entry of $\bm{M}^{\star}$ is observed independently with probability $0<p<1$, namely,
each $(i,j)\in [ n_{1}] \times [n_{2}]$ is included in $\Omega$ independently with probability $p$.
\end{assumption}
%
Under this model, we shall view the expected number of observed entries---namely, $pn_1n_2$---as the sample size. In truth, as long as $p$ is not overly small,  the number of observed entries is expected to concentrate  around its mean $pn_1n_2$.



\paragraph{Incoherence conditions.}
Caution needs to be exercised, however,  that the random sampling model alone does not guarantee effective recovery of an arbitrary low-rank matrix $\bm{M}^{\star}$.
Consider, for example, the following rank-1 matrix $\bm{M}^{\star}$ containing a single nonzero entry:
%
\[
%{\small
\left[\begin{array}{cccc}
1 & 0 & \cdots & 0\\
0 & 0 & \cdots & 0\\
\vdots & \vdots & \ddots & \vdots\\
0 & 0 & \cdots & 0
\end{array}\right] .
%}
\]
%
If $p=o(1)$, then with probability $1-p=1-o(1)$, the sampling pattern will fail to include the nonzero entry $M_{1,1}^{\star}$,
thus ruling out the possibility of faithful matrix recovery.
Consequently,  one needs to make sure that the sampling pattern does not suppress too much useful information.
Towards this end, the pioneering work \citet{candes2009exact, candes2010NearOptimalMC} singled out an incoherence parameter that plays a vital role.
%
\begin{definition}
\label{assump:mc-incoherence}
The incoherence parameter $\mu$ of the matrix $\bm{M}^{\star}\in\mathbb{R}^{n_{1}\times n_{2}}$ is defined as
%
\begin{align*}
	\mu\coloneqq \max\left\{  \frac{ n_1 \|\bm{U}^{\star}\|_{2,\infty}^2 }{r} ,  \frac{ n_2 \|\bm{V}^{\star}\|_{2,\infty}^2 }{r} \right\} .
	%\\
	%\|\bm{V}^{\star}\|_{2,\infty} & \leq\sqrt{\mu/n_{2}}\,\|\bm{V}^{\star}\|_{\mathrm{F}}=\sqrt{\mu r/n_{2}}.
\end{align*}
%
\end{definition}
%
\begin{remark}
\label{remark:range-mu-MC}
%
Recognizing the following basic relation
%
\[
	\frac{r}{n_1}=\frac{1}{n_1} \|\bm{U}^{\star}\|_{\mathrm{F}}^2 \leq \|\bm{U}^{\star}\|_{2,\infty}^2 \leq \|\bm{U}^{\star}\|^2 = 1
\]
%
and an analogous one for $\bm{V}^{\star}$, we have $1\leq \mu \leq \max\{n_1,n_2\}/r= n_2/r$.
%
\end{remark}


In words, a small $\mu$ indicates that the energy of the singular vectors is spread out across different elements, namely,  the singular subspace of $\bm{M}^{\star}$ is not too ``aligned'' with any of the standard basis vectors,
thus ensuring that entrywise observations provide somewhat equalized information about the full spectrum of $\bm{M}^{\star}$.
The following lemma summarizes a few immediate consequences of this definition, with the proof deferred to Section~\ref{sec:proof-auxiliary-lemmas-MC-L2}.
%
\begin{lemma}\label{claim:mu-incoherence}
Assume that  $\bm{M}^{\star}\in\mathbb{R}^{n_{1}\times n_{2}}$ is $\mu$-incoherent. Then the following relations hold
%
\begin{subequations}
\label{subeq:mu-incoherence}
\begin{align}
	\|\bm{M}^{\star}\|_{2,\infty} & \leq\sqrt{\mu r/n_{1}}\, \big\| \bm{M}^{\star} \big\|; \quad
	\|\bm{M}^{\star\top}\|_{2,\infty}  \leq\sqrt{\mu r/n_{2}}\, \big\| \bm{M}^{\star} \big\|;\label{eq:mc-row-norm-upper}\\
	& \qquad\quad \|\bm{M}^{\star}\|_{\infty}  \leq\mu r  \big\| \bm{M}^{\star} \big\| /\sqrt{n_{1}n_{2}} . \label{eq:mc-entry-upper}
\end{align}
\end{subequations}
%
\end{lemma}



\paragraph{Additional notation.}

We find it convenient to introduce a Euclidean projection operator $\mathcal{P}_{\Omega}:\mathbb{R}^{n_{1}\times n_{2}}\mapsto\mathbb{R}^{n_{1}\times n_{2}}$ such that
%
\begin{equation}\label{defn:Pomega}
\big[\mathcal{P}_{\Omega}(\bm{A}) \big]_{i,j}=\begin{cases}
A_{i,j},\quad & \text{if }(i,j)\in\Omega \\
0, & \text{else}
\end{cases}
\end{equation}
%
for any matrix $\bm{A} =[A_{i,j}]\in\mathbb{R}^{n_{1}\times n_{2}}$.
With this notation in place, matrix completion amounts to recovering $\bm{M}^{\star}$ on the basis of $\mathcal{P}_{\Omega}(\bm{M}^{\star})$.




\subsection{Algorithm}\label{sec:mc-alg}

To apply the spectral method, the first step is to form a reasonable approximation $\bm{M}$ of the unknown matrix $\bm{M}^{\star}$.
By virtue of the random sampling model (cf.~Assumption~\ref{assump:mc-bernoulli}), a candidate approximation can be obtained from the observed data matrix via inverse probability weighting:
%
\begin{align}
	\bm{M}   \coloneqq p^{-1}\mathcal{P}_{\Omega}(\bm{M}^{\star}).
	\label{eq:rescaled-M-matrix-completion-l2}
\end{align}
%
The rationale is that $\bm{M}$ forms an unbiased estimate of the ground truth, namely, $$\mathbb{E}[\bm{M}]=\bm{M}^{\star},$$ where
the expectation is taken over the randomness in $\Omega$.


As a result, the proposed spectral method proceeds by computing the rank-$r$ SVD $\bm{U}\bm{\Sigma}\bm{V}^{\top}$ of the matrix $\bm{M}$ constructed in \eqref{eq:rescaled-M-matrix-completion-l2}, and employing $\bm{U}\in \mathbb{R}^{n_1\times r}$, $\bm{V}\in \mathbb{R}^{n_2\times r}$  and $\bm{U}\bm{\Sigma}\bm{V}^{\top}$ as estimates of $\bm{U}^{\star}$, $\bm{V}^{\star}$ and $\bm{M}^{\star}$, respectively.






\subsection{Performance guarantees}\label{sec:mc-l2-theory}

As before, whether the subspace $\bm{U}$ (resp.~$\bm{V}$) is close
to $\bm{U}^{\star}$ (resp.~$\bm{V}^{\star}$) relies crucially on
the size of the perturbation $\|\bm{M}-\bm{M}^{\star}\|$.
Therefore, we begin by developing an upper bound on this quantity; the proof is based on the matrix Bernstein inequality and is postponed to Section~\ref{sec:proof-auxiliary-lemmas-MC-L2}.

\begin{lemma}
\label{lemma:mc-noise-bound}
Consider the settings in Section~\ref{sec:problem-formulation-MC}. Suppose that $n_{2}p\geq C\mu r\log n_{2}$ for some
 constant $C>0$. Then with probability at least
	$1-O(n_{2}^{-10})$, the matrix $\bm{M}$ constructed in \eqref{eq:rescaled-M-matrix-completion-l2} obeys
%
\[
	\big\| \bm{M} -\bm{M}^{\star} \big\|
	\lesssim\sqrt{\frac{\mu r\log n_{2}}{n_{1}p}}\, \big\| \bm{M}^{\star} \big\|.
\]
%
\end{lemma}
%
With this perturbation bound in place, we are equipped to apply Wedin's $\sin\bm{\Theta}$ theorem to obtain the following results.  The condition on the sample size in Theorem~\ref{thm:mc-l2-subspace} is stronger than that in Lemma~\ref{lemma:mc-noise-bound}, as we need to control the eigengap in the following theorem.
%
\begin{theorem}
\label{thm:mc-l2-subspace}
Consider the settings in Section~\ref{sec:problem-formulation-MC}.
Suppose that $n_{1}p\geq C_{1}\kappa^{2}\mu r\log n_{2}$ for some sufficiently
large constant $C_{1}>0$. Then with probability exceeding $1-O(n_2^{-10})$,
%
\begin{align*}
\max\Big\{ \mathsf{dist}\left(\bm{U},\bm{U}^{\star}\right),\mathsf{dist}\left(\bm{V},\bm{V}^{\star}\right)\Big\}  & \lesssim\kappa\sqrt{\frac{\mu r \log n_{2}}{n_{1}p}} .
%;\\
%\max\Big\{ \mathsf{dist}_{\mathrm{F}}\left(\bm{U},\bm{U}^{\star}\right),\mathsf{dist}_{\mathrm{F}}\left(\bm{V},\bm{V}^{\star}\right)\Big\}  & \lesssim\kappa\sqrt{\frac{\mu r^{2} \log n_{2}}{n_{1}p}}.
\end{align*}
%
\end{theorem}
%
\begin{proof}
%	
As a direct consequence of Lemma \ref{lemma:mc-noise-bound}, one has
%
\[
	\big\| \bm{M} -\bm{M}^{\star} \big\|
	\lesssim\sqrt{\frac{\mu r\log n_{2}}{n_{1}p}} \,\big\| \bm{M}^{\star}\big\| \le \Big( 1 - \frac{1}{\sqrt{2}} \Big) \sigma_{r}(\bm{M}^{\star}),
\]
%
provided that $n_{1}p\geq C_{1}\kappa^{2}\mu r\log n_{2}$ for some  large enough constant $C_{1}>0$.
Apply Wedin's theorem (cf.~\eqref{eq:wedin-simpler-friendly}) and Lemma \ref{lemma:mc-noise-bound} to obtain
%
\begin{align*}
 & \max\Big\{ \mathsf{dist}\left(\bm{U},\bm{U}^{\star}\right),\mathsf{dist}\left(\bm{V},\bm{V}^{\star}\right) \Big\}
	\leq  \frac{2\big\|\bm{M} -\bm{M}^{\star}\big\|}{\sigma_{r}(\bm{M}^{\star})}
	%\\ % & \qquad\qquad
	%\lesssim\frac{\sqrt{\frac{\mu r\log n_{2}}{n_{1}p}}\,\big\|\bm{M}^{\star}\big\|}{\sigma_{r}(\bm{M}^{\star})}\asymp
	\lesssim \kappa\sqrt{\frac{\mu r\log n_{2}}{n_{1}p}}
\end{align*}
%
as claimed.
%The proof for the $\mathsf{dist}_{\mathrm{F}}(\cdot,\cdot)$ bounds follows from the same argument and is hence omitted.
%
\end{proof}


As an important implication of Theorem~\ref{thm:mc-l2-subspace}, once the sample size exceeds
$$pn_1n_2 \gg \kappa^{2}\mu r n_2  \log n_{2}, $$
then the spectral estimate achieves consistent estimation in the sense that
%
\[
	\max\Big\{ \mathsf{dist}\left(\bm{U},\bm{U}^{\star}\right),\mathsf{dist}\left(\bm{V},\bm{V}^{\star}\right)\Big\} = o(1).
\]
%
Given that $pn_1n_2\gtrsim \mu n_2 r \log n_2$ is an information-theoretic sampling requirement for reliable matrix completion when $r=o(n_1/\log n_2)$ \citep{candes2010NearOptimalMC},
Theorem~\ref{thm:mc-l2-subspace} confirms the near optimality of spectral methods---in terms of the scaling with $n_1$, $n_2$ and $p$---when it comes to consistent subspace estimation.



Before moving forward to the proof, we further characterize the statistical accuracy of $\bm{U}\bm{\Sigma}\bm{V}^{\top}$ in estimating the unknown matrix $\bm{M}^{\star}$. Accomplishing this only requires Lemma~\ref{lemma:mc-noise-bound}, without any need of the singular subspace perturbation theory. This result will also come in handy when we turn to discussing entrywise estimation accuracy in Chapter~\ref{cha:Linf-theory}.
%
\begin{theorem}
\label{thm:mc-l2-matrix}
Consider the settings in Section~\ref{sec:problem-formulation-MC}.
Suppose that $n_{2}p\geq C\mu r\log n_{2}$ for some sufficiently large constant
$C>0$. Then with probability at least $1-O(n_{2}^{-10})$, one has
%
\begin{align*}
%\|\bm{U}\bm{\Sigma}\bm{V}^{\top}-\bm{M}^{\star}\| & \lesssim\sqrt{\frac{\mu r}{n_{1}p}\log n_{2}}\,\sigma_{1}(\bm{M}^{\star});\\
\|\bm{U}\bm{\Sigma}\bm{V}^{\top}-\bm{M}^{\star}\|_{\mathrm{F}} & \lesssim\sqrt{\frac{\mu r^{2}\log n_{2}}{n_{1}p}}\, \big\| \bm{M}^{\star} \big\|.
\end{align*}
%
\end{theorem}
%
\begin{proof}
First, note
%
\begin{align*}
\|\bm{U}\bm{\Sigma}\bm{V}^{\top}-\bm{M}^{\star}\| &\leq\|\bm{U}\bm{\Sigma}\bm{V}^{\top}- \bm{M} \|+\|\bm{M} -\bm{M}^{\star}\| \leq2\|\bm{M}-\bm{M}^{\star}\|,
\end{align*}
%
where the first inequality comes from the triangle inequality, and the second inequality follows from the fact that $\bm{U}\bm{\Sigma}\bm{V}^{\top}$ is the best rank-$r$ approximation to $\bm{M}$,
i.e.,
$$\|\bm{U}\bm{\Sigma}\bm{V}^{\top}-\bm{M}\|=\min_{\bm{Z}: \mathsf{rank}(\bm{Z})\leq r}\|\bm{Z}-\bm{M}\| \leq \|\bm{M}-\bm{M}^{\star}\|. $$
Additionally, it is observed that $\bm{U}\bm{\Sigma}\bm{V}^{\top}-\bm{M}^{\star}$ has rank at most $2r$, which implies
%
\[
	\|\bm{U}\bm{\Sigma}\bm{V}^{\top}-\bm{M}^{\star}\|_{\mathrm{F}}\leq\sqrt{2r}\|\bm{U}\bm{\Sigma}\bm{V}^{\top}-\bm{M}^{\star}\| \leq 2\sqrt{2r} \big\|\bm{M} - \bm{M}^{\star} \big\|.
\]
%
This combined with Lemma~\ref{lemma:mc-noise-bound} immediately concludes the proof. \end{proof}

 


\subsection{Proof of auxiliary lemmas}
\label{sec:proof-auxiliary-lemmas-MC-L2}

\paragraph{Proof of Lemma~\ref{claim:mu-incoherence}.}
First of all, the  $\|\cdot\|_{2,\infty}$ norm of $\bm{M}^{\star}$
can be upper bounded by
%
\begin{equation*}
\|\bm{M}^{\star}\|_{2,\infty}=\|\bm{U}^{\star}\bm{\Sigma}^{\star}\bm{V}^{\star\top}\|_{2,\infty}\leq\|\bm{U}^{\star}\|_{2,\infty}\|\bm{\Sigma}^{\star}\| \, \|\bm{V}^{\star}\|\leq\sqrt{\frac{\mu r}{n_{1}}} \big\| \bm{M}^{\star} \big\|.
\end{equation*}
%
Here, the first inequality arises from the elementary bounds
$\|\bm{A}\bm{B}\|_{2,\infty}\leq\|\bm{A}\|_{2,\infty}\|\bm{B}\|$
and $\|\bm{A}\bm{B}\|\leq\|\bm{A}\|\, \|\bm{B}\|$, whereas the last relation
uses Definition~\ref{assump:mc-incoherence}, the orthonormality
of $\bm{V}^{\star}$, and identifies $\|\bm{\Sigma}^{\star}\|$ with
$\| \bm{M}^{\star} \|$. The  bound on $\|\bm{M}^{\star\top}\|_{2,\infty}$ can be derived analogously and is omitted for brevity.



In addition, the matrix $\bm{M}^{\star}$ is elementwise bounded by
%
\begin{equation*}
\|\bm{M}^{\star}\|_{\infty}=\|\bm{U}^{\star}\bm{\Sigma}^{\star}\bm{V}^{\star\top}\|_{\infty}\leq\|\bm{U}^{\star}\|_{2,\infty}\|\bm{V}^{\star}\|_{2,\infty} \|\bm{\Sigma}^{\star}\|
	\leq\frac{\mu r}{\sqrt{n_{1}n_{2}}} \big\| \bm{M}^{\star} \big\|.
\end{equation*}
%
Here, the first inequality follows from the fact $\|\bm{A}\bm{B}^{\top}\|_{\infty}\leq\|\bm{A}\|_{2,\infty}\|\bm{B}\|_{2,\infty}$
and the aforementioned one $\|\bm{A}\bm{B}\|_{2,\infty}\leq\|\bm{A}\|_{2,\infty}\|\bm{B}\|$,
while the last inequality again relies on Definition~\ref{assump:mc-incoherence}.
%\end{proof}


\paragraph{Proof of Lemma~\ref{lemma:mc-noise-bound}.} Note that the matrix $\bm{E}\coloneqq p^{-1}\mathcal{P}_{\Omega}(\bm{M}^{\star})-\bm{M}^{\star}$
can be expressed as the sum of $n_{1}n_{2}$ i.i.d.~random matrices
%
\[
  \frac{1}{p}  \mathcal{P}_{\Omega}(\bm{M}^{\star})-\bm{M}^{\star}=\sum_{i=1}^{n_{1}}\sum_{j=1}^{n_{2}}\underbrace{\big(p^{-1} \delta_{i,j} -1 \big)M_{i,j}^{\star}\bm{e}_{i}\bm{e}_{j}^{\top}}_{\eqqcolon\bm{X}_{i,j}}.
\]
%
Here, $\delta_{i,j}$ (which indicates whether the $(i,j)$-th entry is observed) follows an independent Bernoulli distribution
with parameter~$p$, and $\bm{e}_{i}$ stands for the $i$-th standard
basis vector of appropriate dimensions. It is easily seen that for each $(i,j)$,
%
\[
	\mathbb{E}[\bm{X}_{i,j}]=\bm{0}\quad\text{and}\quad\|\bm{X}_{i,j}\|\leq  \frac{1}{p} \|\bm{M}^{\star}\|_{\infty}\leq\frac{\mu r}{p\sqrt{n_{1}n_{2}}} \big\| \bm{M}^{\star} \big\|,
\]
%
where the last relation results from the entrywise upper bound (\ref{eq:mc-entry-upper})
on $\bm{M}^{\star}$. In order to apply the matrix Bernstein inequality (cf.~Corollary~\ref{thm:matrix-Bernstein-friendly}),
we need to control the variance statistic
%
\[
	v\coloneqq\max \Big\{ \Big\Vert \sum\nolimits _{i,j}\mathbb{E}\left[\bm{X}_{i,j}\bm{X}_{i,j}^{\top}\right]\Big\Vert ,\Big\Vert \sum\nolimits _{i,j}\mathbb{E}\left[\bm{X}_{i,j}^{\top}\bm{X}_{i,j}\right] \Big\Vert \Big\} .
\]
%
Regarding the first variance term, we have
%
\begin{align*}
\sum\nolimits _{i,j}\mathbb{E}\left[\bm{X}_{i,j}\bm{X}_{i,j}^{\top}\right]
	& =\sum\nolimits _{i,j}\mathbb{E}\left[ \big( p^{-1} \delta_{i,j} -1 \big)^{2}(M_{i,j}^{\star})^{2}\bm{e}_{i}\bm{e}_{j}^{\top}\bm{e}_{j}\bm{e}_{i}^{\top}\right]\\
 & =\frac{1-p}{p}\sum\nolimits _{i,j}\big(M_{i,j}^{\star}\big)^{2}\bm{e}_{i}\bm{e}_{i}^{\top}=\frac{1-p}{p}\sum_{i=1}^{n_{1}}\|\bm{M}_{i,\cdot}^{\star}\|_{2}^{2}\bm{e}_{i}\bm{e}_{i}^{\top} \\
	& \preceq \frac{1-p}{p}  \|\bm{M}^{\star}\|_{2,\infty}^{2} \bm{I}_{n_1} \preceq \frac{\mu r}{n_1p}  \|\bm{M}^{\star}\|^{2} \,\bm{I}_{n_1} .
\end{align*}
%
Here, the first identity arises from the definition of $\bm{X}_{i,j}$,
the second one calculates the variance of Bernoulli random variables,
and the last line relies on the upper bound~(\ref{eq:mc-row-norm-upper}).
Similarly, the second term in the variance statistic enjoys the following characterization:
%
\[
	\sum\nolimits _{i,j}\mathbb{E}\left[\bm{X}_{i,j}^{\top}\bm{X}_{i,j}\right] \preceq \frac{\mu r}{n_2p} \|\bm{M}^{\star}\|^{2} \,\bm{I}_{n_2}.
\]
%
Taking the above  relations together and recalling that $n_1\leq n_2$ give
%
\begin{align*}
v & %\leq\frac{1}{p}\max\left\{ \|\bm{M}^{\star}\|_{2,\infty}^{2},\|\bm{M}^{\star\top}\|_{2,\infty}^{2}\right\}
	% \leq\frac{1}{p}\left\{ \frac{\mu r}{n_{1}} \big\| \bm{M}^{\star} \big\|^2,  \frac{\mu r}{n_{2}} \big\| \bm{M}^{\star} \big\|^2 \right\} \\
	\leq\frac{\mu r}{n_{1}p} \big\| \bm{M}^{\star} \big\|^2.
\end{align*}
% 


With the above bounds in place, invoking matrix Bernstein (see Corollary~\ref{thm:matrix-Bernstein-friendly}) reveals that: with probability at least $1-O(n_2^{-10})$,
%
\begin{align*}
	\| \bm{E}\| & \lesssim \sqrt{\frac{\mu r \| \bm{M}^{\star} \|^2 \log n_{2}}{n_{1}p} }+\frac{\mu r \| \bm{M}^{\star} \| \log n_{2}}{p\sqrt{n_{1}n_{2}}}
	\asymp \sqrt{\frac{\mu r \| \bm{M}^{\star} \|^2 \log n_{2}}{n_{1}p} } ,
\end{align*}
%
where the last inequality is valid as long as  $n_{2}p \gtrsim \mu r\log n_{2}$.







\section{Application: Confidence intervals for matrix completion}
\label{sec:CI-matrix-completion}

As an illustration of the applicability of the inference procedure described in Section~\ref{sec:UQ-general}, 
we develop concrete consequences of Theorem~\ref{thm:CI-general} in application to noisy matrix completion---an extension of the formulation
in Section~\ref{sec:mc-l2} to noisy settings. 


\paragraph{Model: noisy matrix completion.} 

Suppose that we are asked to reconstruct a symmetric rank-$r$ matrix
 $\bm{M}^{\star}=[M_{i,j}^{\star}]_{1\leq k,l \leq n}\in \mathbb{R}^{n\times n}$  
 with eigendecomposition $\bm{M}^{\star}=\bm{U}^{\star}\bm{\Lambda}^{\star}\bm{U}^{\star\top}$.
 We only get to acquire noisy observations of a subset of the entries of $\bm{M}^{\star}$; more precisely, 
there exists a sampling set $\Omega \subseteq [n]\times [n]$ such that we observe
%
\begin{align}
	M_{k,l}^{\star} + \eta_{k,l} \qquad & \text{if } (k,l)\in \Omega. 
	% M_{i,j} = \begin{cases} , \qquad & \text{if } (i,j)\in \Omega, \\ 0, &\text{else.}  \end{cases}
\end{align}
%
Here, $\{\eta_{k,l}\mid k\geq l\}$ denotes independent Gaussian noise obeying  
%
\begin{align}
	\eta_{k,l} = \eta_{l,k} \overset{\mathrm{i.i.d.}}{\sim} \mathcal{N}(0, \sigma_{\eta}^2), \qquad k\geq l. 
\end{align}
%
As before, we focus on the random sampling model such that each location $(k,l)$ with $k\geq l$ is included in the sampling set $\Omega$ independently with probability $p$. Further,
 assume that $\bm{M}^{\star}$ has eigenvalues obeying \eqref{eq:assumption-lambda-r-positive-setup}, condition number $\kappa$ (cf.~\eqref{eq:defn-kappa}), and  incoherence parameter $\mu$ (cf.~\eqref{eq:defn-Ustar-incoherence}). 
Can we build a confidence interval for each entry $M_{i,j}^{\star}$, on the basis of the output of the spectral method?  






\paragraph{Computing entrywise confidence intervals.} 
%
In order to apply the inference procedure in Section~\ref{sec:UQ-general}, it suffices to determine the data matrix $\bm{M}=[M_{i,j}]_{1\leq i,j\leq n}$, which can be selected as usual. Specifically, a possible inference procedure proceeds as follows:
%
\begin{itemize}
	\item  Set $\bm{M}$ such that for any $1\leq i,j \leq n$, 
%
\begin{align}
	M_{i,j} = \begin{cases} \frac{1}{p} \big( M_{i,j}^{\star} + \eta_{i,j} \big), \qquad & \text{if } (i,j)\in \Omega, \\ 0, &\text{else}, \end{cases}
\end{align}
%
which clearly obeys $\mathbb{E}[\bm{M}]=\bm{M}^{\star}$. 

	\item Compute the estimate $\widehat{\bm{M}}$ (cf.~\eqref{eq:estimator-M-hat-distribution}) via the spectral method.

	\item For a given coverage level $1-\alpha$ and a given pair $(i,j)$, 
		construct the confidence interval $\mathsf{CI}_{i,j}^{1-\alpha}$ according to \eqref{eq:confidence-interval-construction-general}, 
		with auxiliary parameters provided in \eqref{eq:v-ij-plugin-estimator} and \eqref{eq:estimators-E-and-P}. 

\end{itemize}
 


\paragraph{Performance guarantees and implications.}
When specialized to noisy matrix completion, our inference theory in Theorem~\ref{thm:CI-general} leads to the following statistical guarantees. 
%
\begin{theorem}
	\label{thm:CI-noisy-matrix-completion}
	Consider the noisy matrix completion setting in this section. 
	Suppose that $\kappa^{4}\mu^{2}r^{2}\log n\leq n$, 
	%$\sigma \kappa\sqrt{n\log n}\lesssim |\lambda_{r}^{\star}|$ and 
	$\frac{\max_{k,l}|M_{k,l}^{\star}|}{\min_{k,l}|M_{k,l}^{\star}|}=O(1)$, 
	%
	\begin{equation}
		np\gtrsim \kappa^{4}r\log n\quad\text{and}\quad\sigma_{\eta}\kappa\sqrt{\frac{n\log n}{p}}\lesssim|\lambda_{r}^{\star}|.  
		\label{eq:Mstar-noise-requirement-noisy-MC}
	\end{equation}
	%	
	Consider any $1\leq i,j\leq n$, and assume that
	%
\begin{align}
	\frac{\big\|\bm{U}_{j,\cdot}^{\star}\big\|_{2}^{2}+\big\|\bm{U}_{i,\cdot}^{\star}\big\|_{2}^{2}}{\big\|\bm{U}^{\star}\big\|_{\mathrm{F}}^{2}} & \gtrsim\frac{\mu^{2}r^{2}\kappa^{4}\log^{3}n}{n\sqrt{np}}+\frac{\sigma_{\eta}\mu^{2}r\kappa^{3}\log^{3}n}{|\lambda_{r}^{\star}|\sqrt{np}}.
	\label{eq:Ui-Uj-lower-bound-thm-noisy-MC-146}
\end{align}
%
	For any fixed coverage level $1-\alpha \in (0,1)$, 
	the confidence interval $\mathsf{CI}_{i,j}^{1-\alpha}$ constructed in \eqref{eq:confidence-interval-construction-general} obeys
	%
	\begin{align}
		\mathbb{P} \Big( M_{i,j}^{\star} \in \mathsf{CI}_{i,j}^{1-\alpha}  
		\Big) 
		= 1-\alpha + o(1).
	\end{align}
	%
\end{theorem}
%


In order to help interpret the applicable range of Theorem~\ref{thm:CI-noisy-matrix-completion}, 
let us focus on the simple scenario with $\kappa, \mu, r\asymp 1$ to simplify discussion. 

\begin{itemize}

	\item First of all, Condition~\eqref{eq:Mstar-noise-requirement-noisy-MC} can be simplified as
	%
	\[
		np  \gtrsim  \log n\quad\text{and}\quad \sigma_{\eta}\sqrt{\frac{n\log n}{p}}\lesssim |\lambda_r^{\star}|. 
	\]
	%
	The first condition on the sampling size coincides with the fundamental requirement even if the goal is merely to enable reliable estimation \citep{candes2010NearOptimalMC}, 
	whereas the second condition on the signal-to-noise ratio is also necessary---up to some log factor---to ensure an estimation quality better than that of a random guess \citep[Theorem~3.3]{cai2019subspace}.  


	\item Next, we move on to interpret the other condition \eqref{eq:Ui-Uj-lower-bound-thm-noisy-MC-146} imposed in our theory, which simplifies to
	%
	\[
		\frac{\big\|\bm{U}_{j,\cdot}^{\star}\big\|_{2}^{2}+\big\|\bm{U}_{i,\cdot}^{\star}\big\|_{2}^{2}}{\big\|\bm{U}^{\star}\big\|_{\mathrm{F}}^{2}}
		\gtrsim\frac{\log^{3}n}{n\sqrt{np}}+\frac{\sigma_{\eta}\log^{3}n}{|\lambda_{r}^{\star}|\sqrt{np}}.
	\]
	%
	Let us consider the most challenging case where $np \gtrsim  \mathrm{poly}\log n$ and $\sigma_{\eta}\sqrt{\frac{n\mathrm{poly}\log n}{p}}\lesssim |\lambda_r^{\star}|$ (for some sufficiently large poly-log factor). 
	In such a case, the above condition only requires 
	%
	\[
		\frac{\big\|\bm{U}_{j,\cdot}^{\star}\big\|_{2}^{2}+\big\|\bm{U}_{i,\cdot}^{\star}\big\|_{2}^{2}}{\big\|\bm{U}^{\star}\big\|_{\mathrm{F}}^{2}}\gtrsim\frac{1}{n\mathrm{poly}\log n},
	\]
	%
	indicating that the associated signal power $\|\bm{U}_{j,\cdot}^{\star}\|_{2}^{2}+\|\bm{U}_{i,\cdot}^{\star}\|_{2}^{2}$ 
		is allowed to be much smaller than the average signal power across all rows (which can be captured by 
		$\|\bm{U}^{\star}\|_{\mathrm{F}}^{2}/n$). 

\end{itemize}
%
In a nutshell, the validity of our inference procedure is ensured for broad settings. 
Additionally, we have conducted a series of numerical experiments to examine the entrywise distributions of $\bm{M}$. As illustrated in Figure~\ref{fig:MC-inference-numerics}, 
the normalized estimation error $(\widehat{v}_{i,j})^{-1/2} (\widehat{M}_{i,j} - M_{i,j}^{\star})$ is close in distribution to a standard Gaussian random variable, 
which corroborates our theory on the confidence interval construction. 



Before concluding, we would like to remark that: 
while the distributional theory for spectral methods allows for valid construction of confidence intervals for an unseen entry,  
it is oftentimes not among the most effective statistical inference procedures that one can put forward. 
There exist other alternatives that are provably more efficient, including but not limited to inference procedures based on convex relaxation and nonconvex optimization \citep{chen2019inference,xia2021statistical}, and the ones based on more refined spectral methods \citep{yan2021inference,chernozhukov2021inference}. 


\begin{figure}

	\begin{tabular}{cc}
		\includegraphics[width=0.45\textwidth]{figures/mc_inference_histogram_plot.pdf} & \includegraphics[width=0.45\textwidth]{figures/mc_inference_QQ_plot.pdf} \tabularnewline
		(a)  & (b) \tabularnewline
	\end{tabular}
	
	\caption{Entrywise numerical distribution for noisy matrix completion. We generate $\bm{M}^{\star}=\bm{U}^{\star}\bm{U}^{\star\top}$ with $\bm{U}^{\star}$ being a random orthonormal matrix, and analyze the behavior of $Z_{1,2}= (\widehat{v}_{1,2})^{-0.5} ( \widehat{M}_{1,2}-M_{1,2}^{\star} )$, 
		where $\widehat{\bm{M}}$ is defined in \eqref{eq:estimator-M-hat-distribution} and $\widehat{v}_{i,j}$ is defined in \eqref{eq:v-ij-plugin-estimator}.
		The results are reported for 500 Monte Carlo trials when $p=0.3$, $n=1000$, $\sigma_{\eta} = 10^{-4}$, and $r=3$.  
		 (a) Histogram of the empirical distribution of $Z_{1,2}$; (b)  Q-Q (quantile-quantile) plot of $Z_{1,2}$ vs.~the standard normal distribution. \label{fig:MC-inference-numerics}}  
\end{figure}



%
\section{Matrix completion}\label{sec:mc-l2}

A pressing challenge often encountered in data science applications is estimation and learning in the face of \emph{missing data}.
To elucidate how to tackle this challenge via spectral methods, we delve into the renowned matrix completion problem in this section,
followed by another application called tensor completion in Section~\ref{sec:tensor-completion}.


Imagine that one observes a small subset of the entries in a large unknown matrix
and seeks to fill in all missing entries. An archetypal example is collaborative filtering, where one aims to predict the users' preferences on a collection of products based on partially revealed user-product ratings.   See Figure~\ref{fig:completion} for an illustration.  The problem, often referred to as matrix completion, is apparently ill-posed in general, as there are (much) fewer measurements than the unknowns.

%This is indeed a factor model in Section~\ref{sec:PCA} with missing data.  

Fortunately, if the matrix of interest exhibits certain low-dimensional structure, then reliable recovery becomes feasible.
A commonly encountered example of this kind concerns the case when the target matrix enjoys a low-rank structure. Again, take collaborative filtering for example:  the user-product rating matrix  might be well explained by a relatively small number of latent factors connecting users' preferences with products' attributes, thus resulting in an approximately low-rank matrix.  Motivated by its fundamental importance,  recent years have witnessed a flurry of research activity in studying low-rank matrix completion \citep{candes2009exact, keshavan2010matrix, gross2011recovering};
see \citet{chen2018harnessing,davenport2016overview} for overviews of recent developments.
In the sequel, we present a simple yet effective approach enabled by the spectral method, originally proposed in~\citet{achlioptas2007fast,keshavan2010matrix}.

\begin{figure}
\begin{center}
\includegraphics[width=0.35\textwidth]{figures/matrix_completion}
\end{center}
\caption{Illustration of matrix completion, where each ``$\checkmark$'' stands for an observed entry and each ``$?$'' represents a missing entry.} \label{fig:completion}
\end{figure}


\subsection{Problem formulation and assumptions}
\label{sec:problem-formulation-MC}

Suppose that we are interested in estimating an $n_{1}\times n_{2}$  rank-$r$ matrix  $\bm{M}^{\star} = [{M}^{\star}_{i,j}]_{1\leq i\leq n_1, 1\leq j\leq n_2}$.
Without loss of generality, we assume
$$n_{1}\leq n_{2}.$$  Denote by  $\bm{M}^{\star}=\bm{U}^{\star}\bm{\Sigma}^{\star}\bm{V}^{\star\top}$ the SVD of $\bm{M}^{\star}$, where the columns of $\bm{U}^{\star}\in \mathbb{R}^{n_1\times r}$ (resp.~$\bm{V}^{\star}\in \mathbb{R}^{n_2\times r}$) are the left (resp.~right) singular vectors of $\bm{M}^{\star}$, and $\bm{\Sigma}^{\star}$ is a diagonal matrix whose diagonal entries are the singular values of $\bm{M}^{\star}$. We define the condition number of the matrix $\bm{M}^{\star}$ to be $\kappa \coloneqq \sigma_1(\bM^\star)/\sigma_r(\bM^\star)$.


To capture the presence of missing data, we introduce an index subset $\Omega\subseteq [n_1]\times [n_2]$,
such that each entry ${M}^{\star}_{i,j}$ is observed if and only if $(i,j)\in \Omega$.
The goal is to reconstruct the singular subspaces $\bm{U}^{\star}$ and $\bm{V}^{\star}$, as well as the  full matrix $\bm{M}^{\star}$, based on entries observed over the sampling set $\Omega$.



\paragraph{Random sampling.}
Apparently, not all sampling patterns admit reliable estimation. For instance, if $\Omega$ contains only entries in the top half of the matrix,
then there is in general no hope to predict the bottom half of the matrix. In order to allow for meaningful matrix completion,
this monograph focuses on a natural random observation model commonly adopted in the literature, as formulated below.
%
\begin{assumption}[Random sampling]
\label{assump:mc-bernoulli}
Each entry of $\bm{M}^{\star}$ is observed independently with probability $0<p<1$, namely,
each $(i,j)\in [ n_{1}] \times [n_{2}]$ is included in $\Omega$ independently with probability $p$.
\end{assumption}
%
Under this model, we shall view the expected number of observed entries---namely, $pn_1n_2$---as the sample size. In truth, as long as $p$ is not overly small,  the number of observed entries is expected to concentrate  around its mean $pn_1n_2$.



\paragraph{Incoherence conditions.}
Caution needs to be exercised, however,  that the random sampling model alone does not guarantee effective recovery of an arbitrary low-rank matrix $\bm{M}^{\star}$.
Consider, for example, the following rank-1 matrix $\bm{M}^{\star}$ containing a single nonzero entry:
%
\[
%{\small
\left[\begin{array}{cccc}
1 & 0 & \cdots & 0\\
0 & 0 & \cdots & 0\\
\vdots & \vdots & \ddots & \vdots\\
0 & 0 & \cdots & 0
\end{array}\right] .
%}
\]
%
If $p=o(1)$, then with probability $1-p=1-o(1)$, the sampling pattern will fail to include the nonzero entry $M_{1,1}^{\star}$,
thus ruling out the possibility of faithful matrix recovery.
Consequently,  one needs to make sure that the sampling pattern does not suppress too much useful information.
Towards this end, the pioneering work \citet{candes2009exact, candes2010NearOptimalMC} singled out an incoherence parameter that plays a vital role.
%
\begin{definition}
\label{assump:mc-incoherence}
The incoherence parameter $\mu$ of the matrix $\bm{M}^{\star}\in\mathbb{R}^{n_{1}\times n_{2}}$ is defined as
%
\begin{align*}
	\mu\coloneqq \max\left\{  \frac{ n_1 \|\bm{U}^{\star}\|_{2,\infty}^2 }{r} ,  \frac{ n_2 \|\bm{V}^{\star}\|_{2,\infty}^2 }{r} \right\} .
	%\\
	%\|\bm{V}^{\star}\|_{2,\infty} & \leq\sqrt{\mu/n_{2}}\,\|\bm{V}^{\star}\|_{\mathrm{F}}=\sqrt{\mu r/n_{2}}.
\end{align*}
%
\end{definition}
%
\begin{remark}
\label{remark:range-mu-MC}
%
Recognizing the following basic relation
%
\[
	\frac{r}{n_1}=\frac{1}{n_1} \|\bm{U}^{\star}\|_{\mathrm{F}}^2 \leq \|\bm{U}^{\star}\|_{2,\infty}^2 \leq \|\bm{U}^{\star}\|^2 = 1
\]
%
and an analogous one for $\bm{V}^{\star}$, we have $1\leq \mu \leq \max\{n_1,n_2\}/r= n_2/r$.
%
\end{remark}


In words, a small $\mu$ indicates that the energy of the singular vectors is spread out across different elements, namely,  the singular subspace of $\bm{M}^{\star}$ is not too ``aligned'' with any of the standard basis vectors,
thus ensuring that entrywise observations provide somewhat equalized information about the full spectrum of $\bm{M}^{\star}$.
The following lemma summarizes a few immediate consequences of this definition, with the proof deferred to Section~\ref{sec:proof-auxiliary-lemmas-MC-L2}.
%
\begin{lemma}\label{claim:mu-incoherence}
Assume that  $\bm{M}^{\star}\in\mathbb{R}^{n_{1}\times n_{2}}$ is $\mu$-incoherent. Then the following relations hold
%
\begin{subequations}
\label{subeq:mu-incoherence}
\begin{align}
	\|\bm{M}^{\star}\|_{2,\infty} & \leq\sqrt{\mu r/n_{1}}\, \big\| \bm{M}^{\star} \big\|; \quad
	\|\bm{M}^{\star\top}\|_{2,\infty}  \leq\sqrt{\mu r/n_{2}}\, \big\| \bm{M}^{\star} \big\|;\label{eq:mc-row-norm-upper}\\
	& \qquad\quad \|\bm{M}^{\star}\|_{\infty}  \leq\mu r  \big\| \bm{M}^{\star} \big\| /\sqrt{n_{1}n_{2}} . \label{eq:mc-entry-upper}
\end{align}
\end{subequations}
%
\end{lemma}



\paragraph{Additional notation.}

We find it convenient to introduce a Euclidean projection operator $\mathcal{P}_{\Omega}:\mathbb{R}^{n_{1}\times n_{2}}\mapsto\mathbb{R}^{n_{1}\times n_{2}}$ such that
%
\begin{equation}\label{defn:Pomega}
\big[\mathcal{P}_{\Omega}(\bm{A}) \big]_{i,j}=\begin{cases}
A_{i,j},\quad & \text{if }(i,j)\in\Omega \\
0, & \text{else}
\end{cases}
\end{equation}
%
for any matrix $\bm{A} =[A_{i,j}]\in\mathbb{R}^{n_{1}\times n_{2}}$.
With this notation in place, matrix completion amounts to recovering $\bm{M}^{\star}$ on the basis of $\mathcal{P}_{\Omega}(\bm{M}^{\star})$.




\subsection{Algorithm}\label{sec:mc-alg}

To apply the spectral method, the first step is to form a reasonable approximation $\bm{M}$ of the unknown matrix $\bm{M}^{\star}$.
By virtue of the random sampling model (cf.~Assumption~\ref{assump:mc-bernoulli}), a candidate approximation can be obtained from the observed data matrix via inverse probability weighting:
%
\begin{align}
	\bm{M}   \coloneqq p^{-1}\mathcal{P}_{\Omega}(\bm{M}^{\star}).
	\label{eq:rescaled-M-matrix-completion-l2}
\end{align}
%
The rationale is that $\bm{M}$ forms an unbiased estimate of the ground truth, namely, $$\mathbb{E}[\bm{M}]=\bm{M}^{\star},$$ where
the expectation is taken over the randomness in $\Omega$.


As a result, the proposed spectral method proceeds by computing the rank-$r$ SVD $\bm{U}\bm{\Sigma}\bm{V}^{\top}$ of the matrix $\bm{M}$ constructed in \eqref{eq:rescaled-M-matrix-completion-l2}, and employing $\bm{U}\in \mathbb{R}^{n_1\times r}$, $\bm{V}\in \mathbb{R}^{n_2\times r}$  and $\bm{U}\bm{\Sigma}\bm{V}^{\top}$ as estimates of $\bm{U}^{\star}$, $\bm{V}^{\star}$ and $\bm{M}^{\star}$, respectively.






\subsection{Performance guarantees}\label{sec:mc-l2-theory}

As before, whether the subspace $\bm{U}$ (resp.~$\bm{V}$) is close
to $\bm{U}^{\star}$ (resp.~$\bm{V}^{\star}$) relies crucially on
the size of the perturbation $\|\bm{M}-\bm{M}^{\star}\|$.
Therefore, we begin by developing an upper bound on this quantity; the proof is based on the matrix Bernstein inequality and is postponed to Section~\ref{sec:proof-auxiliary-lemmas-MC-L2}.

\begin{lemma}
\label{lemma:mc-noise-bound}
Consider the settings in Section~\ref{sec:problem-formulation-MC}. Suppose that $n_{2}p\geq C\mu r\log n_{2}$ for some
 constant $C>0$. Then with probability at least
	$1-O(n_{2}^{-10})$, the matrix $\bm{M}$ constructed in \eqref{eq:rescaled-M-matrix-completion-l2} obeys
%
\[
	\big\| \bm{M} -\bm{M}^{\star} \big\|
	\lesssim\sqrt{\frac{\mu r\log n_{2}}{n_{1}p}}\, \big\| \bm{M}^{\star} \big\|.
\]
%
\end{lemma}
%
With this perturbation bound in place, we are equipped to apply Wedin's $\sin\bm{\Theta}$ theorem to obtain the following results.  The condition on the sample size in Theorem~\ref{thm:mc-l2-subspace} is stronger than that in Lemma~\ref{lemma:mc-noise-bound}, as we need to control the eigengap in the following theorem.
%
\begin{theorem}
\label{thm:mc-l2-subspace}
Consider the settings in Section~\ref{sec:problem-formulation-MC}.
Suppose that $n_{1}p\geq C_{1}\kappa^{2}\mu r\log n_{2}$ for some sufficiently
large constant $C_{1}>0$. Then with probability exceeding $1-O(n_2^{-10})$,
%
\begin{align*}
\max\Big\{ \mathsf{dist}\left(\bm{U},\bm{U}^{\star}\right),\mathsf{dist}\left(\bm{V},\bm{V}^{\star}\right)\Big\}  & \lesssim\kappa\sqrt{\frac{\mu r \log n_{2}}{n_{1}p}} .
%;\\
%\max\Big\{ \mathsf{dist}_{\mathrm{F}}\left(\bm{U},\bm{U}^{\star}\right),\mathsf{dist}_{\mathrm{F}}\left(\bm{V},\bm{V}^{\star}\right)\Big\}  & \lesssim\kappa\sqrt{\frac{\mu r^{2} \log n_{2}}{n_{1}p}}.
\end{align*}
%
\end{theorem}
%
\begin{proof}
%	
As a direct consequence of Lemma \ref{lemma:mc-noise-bound}, one has
%
\[
	\big\| \bm{M} -\bm{M}^{\star} \big\|
	\lesssim\sqrt{\frac{\mu r\log n_{2}}{n_{1}p}} \,\big\| \bm{M}^{\star}\big\| \le \Big( 1 - \frac{1}{\sqrt{2}} \Big) \sigma_{r}(\bm{M}^{\star}),
\]
%
provided that $n_{1}p\geq C_{1}\kappa^{2}\mu r\log n_{2}$ for some  large enough constant $C_{1}>0$.
Apply Wedin's theorem (cf.~\eqref{eq:wedin-simpler-friendly}) and Lemma \ref{lemma:mc-noise-bound} to obtain
%
\begin{align*}
 & \max\Big\{ \mathsf{dist}\left(\bm{U},\bm{U}^{\star}\right),\mathsf{dist}\left(\bm{V},\bm{V}^{\star}\right) \Big\}
	\leq  \frac{2\big\|\bm{M} -\bm{M}^{\star}\big\|}{\sigma_{r}(\bm{M}^{\star})}
	%\\ % & \qquad\qquad
	%\lesssim\frac{\sqrt{\frac{\mu r\log n_{2}}{n_{1}p}}\,\big\|\bm{M}^{\star}\big\|}{\sigma_{r}(\bm{M}^{\star})}\asymp
	\lesssim \kappa\sqrt{\frac{\mu r\log n_{2}}{n_{1}p}}
\end{align*}
%
as claimed.
%The proof for the $\mathsf{dist}_{\mathrm{F}}(\cdot,\cdot)$ bounds follows from the same argument and is hence omitted.
%
\end{proof}


As an important implication of Theorem~\ref{thm:mc-l2-subspace}, once the sample size exceeds
$$pn_1n_2 \gg \kappa^{2}\mu r n_2  \log n_{2}, $$
then the spectral estimate achieves consistent estimation in the sense that
%
\[
	\max\Big\{ \mathsf{dist}\left(\bm{U},\bm{U}^{\star}\right),\mathsf{dist}\left(\bm{V},\bm{V}^{\star}\right)\Big\} = o(1).
\]
%
Given that $pn_1n_2\gtrsim \mu n_2 r \log n_2$ is an information-theoretic sampling requirement for reliable matrix completion when $r=o(n_1/\log n_2)$ \citep{candes2010NearOptimalMC},
Theorem~\ref{thm:mc-l2-subspace} confirms the near optimality of spectral methods---in terms of the scaling with $n_1$, $n_2$ and $p$---when it comes to consistent subspace estimation.



Before moving forward to the proof, we further characterize the statistical accuracy of $\bm{U}\bm{\Sigma}\bm{V}^{\top}$ in estimating the unknown matrix $\bm{M}^{\star}$. Accomplishing this only requires Lemma~\ref{lemma:mc-noise-bound}, without any need of the singular subspace perturbation theory. This result will also come in handy when we turn to discussing entrywise estimation accuracy in Chapter~\ref{cha:Linf-theory}.
%
\begin{theorem}
\label{thm:mc-l2-matrix}
Consider the settings in Section~\ref{sec:problem-formulation-MC}.
Suppose that $n_{2}p\geq C\mu r\log n_{2}$ for some sufficiently large constant
$C>0$. Then with probability at least $1-O(n_{2}^{-10})$, one has
%
\begin{align*}
%\|\bm{U}\bm{\Sigma}\bm{V}^{\top}-\bm{M}^{\star}\| & \lesssim\sqrt{\frac{\mu r}{n_{1}p}\log n_{2}}\,\sigma_{1}(\bm{M}^{\star});\\
\|\bm{U}\bm{\Sigma}\bm{V}^{\top}-\bm{M}^{\star}\|_{\mathrm{F}} & \lesssim\sqrt{\frac{\mu r^{2}\log n_{2}}{n_{1}p}}\, \big\| \bm{M}^{\star} \big\|.
\end{align*}
%
\end{theorem}
%
\begin{proof}
First, note
%
\begin{align*}
\|\bm{U}\bm{\Sigma}\bm{V}^{\top}-\bm{M}^{\star}\| &\leq\|\bm{U}\bm{\Sigma}\bm{V}^{\top}- \bm{M} \|+\|\bm{M} -\bm{M}^{\star}\| \leq2\|\bm{M}-\bm{M}^{\star}\|,
\end{align*}
%
where the first inequality comes from the triangle inequality, and the second inequality follows from the fact that $\bm{U}\bm{\Sigma}\bm{V}^{\top}$ is the best rank-$r$ approximation to $\bm{M}$,
i.e.,
$$\|\bm{U}\bm{\Sigma}\bm{V}^{\top}-\bm{M}\|=\min_{\bm{Z}: \mathsf{rank}(\bm{Z})\leq r}\|\bm{Z}-\bm{M}\| \leq \|\bm{M}-\bm{M}^{\star}\|. $$
Additionally, it is observed that $\bm{U}\bm{\Sigma}\bm{V}^{\top}-\bm{M}^{\star}$ has rank at most $2r$, which implies
%
\[
	\|\bm{U}\bm{\Sigma}\bm{V}^{\top}-\bm{M}^{\star}\|_{\mathrm{F}}\leq\sqrt{2r}\|\bm{U}\bm{\Sigma}\bm{V}^{\top}-\bm{M}^{\star}\| \leq 2\sqrt{2r} \big\|\bm{M} - \bm{M}^{\star} \big\|.
\]
%
This combined with Lemma~\ref{lemma:mc-noise-bound} immediately concludes the proof. \end{proof}

 


\subsection{Proof of auxiliary lemmas}
\label{sec:proof-auxiliary-lemmas-MC-L2}

\paragraph{Proof of Lemma~\ref{claim:mu-incoherence}.}
First of all, the  $\|\cdot\|_{2,\infty}$ norm of $\bm{M}^{\star}$
can be upper bounded by
%
\begin{equation*}
\|\bm{M}^{\star}\|_{2,\infty}=\|\bm{U}^{\star}\bm{\Sigma}^{\star}\bm{V}^{\star\top}\|_{2,\infty}\leq\|\bm{U}^{\star}\|_{2,\infty}\|\bm{\Sigma}^{\star}\| \, \|\bm{V}^{\star}\|\leq\sqrt{\frac{\mu r}{n_{1}}} \big\| \bm{M}^{\star} \big\|.
\end{equation*}
%
Here, the first inequality arises from the elementary bounds
$\|\bm{A}\bm{B}\|_{2,\infty}\leq\|\bm{A}\|_{2,\infty}\|\bm{B}\|$
and $\|\bm{A}\bm{B}\|\leq\|\bm{A}\|\, \|\bm{B}\|$, whereas the last relation
uses Definition~\ref{assump:mc-incoherence}, the orthonormality
of $\bm{V}^{\star}$, and identifies $\|\bm{\Sigma}^{\star}\|$ with
$\| \bm{M}^{\star} \|$. The  bound on $\|\bm{M}^{\star\top}\|_{2,\infty}$ can be derived analogously and is omitted for brevity.



In addition, the matrix $\bm{M}^{\star}$ is elementwise bounded by
%
\begin{equation*}
\|\bm{M}^{\star}\|_{\infty}=\|\bm{U}^{\star}\bm{\Sigma}^{\star}\bm{V}^{\star\top}\|_{\infty}\leq\|\bm{U}^{\star}\|_{2,\infty}\|\bm{V}^{\star}\|_{2,\infty} \|\bm{\Sigma}^{\star}\|
	\leq\frac{\mu r}{\sqrt{n_{1}n_{2}}} \big\| \bm{M}^{\star} \big\|.
\end{equation*}
%
Here, the first inequality follows from the fact $\|\bm{A}\bm{B}^{\top}\|_{\infty}\leq\|\bm{A}\|_{2,\infty}\|\bm{B}\|_{2,\infty}$
and the aforementioned one $\|\bm{A}\bm{B}\|_{2,\infty}\leq\|\bm{A}\|_{2,\infty}\|\bm{B}\|$,
while the last inequality again relies on Definition~\ref{assump:mc-incoherence}.
%\end{proof}


\paragraph{Proof of Lemma~\ref{lemma:mc-noise-bound}.} Note that the matrix $\bm{E}\coloneqq p^{-1}\mathcal{P}_{\Omega}(\bm{M}^{\star})-\bm{M}^{\star}$
can be expressed as the sum of $n_{1}n_{2}$ i.i.d.~random matrices
%
\[
  \frac{1}{p}  \mathcal{P}_{\Omega}(\bm{M}^{\star})-\bm{M}^{\star}=\sum_{i=1}^{n_{1}}\sum_{j=1}^{n_{2}}\underbrace{\big(p^{-1} \delta_{i,j} -1 \big)M_{i,j}^{\star}\bm{e}_{i}\bm{e}_{j}^{\top}}_{\eqqcolon\bm{X}_{i,j}}.
\]
%
Here, $\delta_{i,j}$ (which indicates whether the $(i,j)$-th entry is observed) follows an independent Bernoulli distribution
with parameter~$p$, and $\bm{e}_{i}$ stands for the $i$-th standard
basis vector of appropriate dimensions. It is easily seen that for each $(i,j)$,
%
\[
	\mathbb{E}[\bm{X}_{i,j}]=\bm{0}\quad\text{and}\quad\|\bm{X}_{i,j}\|\leq  \frac{1}{p} \|\bm{M}^{\star}\|_{\infty}\leq\frac{\mu r}{p\sqrt{n_{1}n_{2}}} \big\| \bm{M}^{\star} \big\|,
\]
%
where the last relation results from the entrywise upper bound (\ref{eq:mc-entry-upper})
on $\bm{M}^{\star}$. In order to apply the matrix Bernstein inequality (cf.~Corollary~\ref{thm:matrix-Bernstein-friendly}),
we need to control the variance statistic
%
\[
	v\coloneqq\max \Big\{ \Big\Vert \sum\nolimits _{i,j}\mathbb{E}\left[\bm{X}_{i,j}\bm{X}_{i,j}^{\top}\right]\Big\Vert ,\Big\Vert \sum\nolimits _{i,j}\mathbb{E}\left[\bm{X}_{i,j}^{\top}\bm{X}_{i,j}\right] \Big\Vert \Big\} .
\]
%
Regarding the first variance term, we have
%
\begin{align*}
\sum\nolimits _{i,j}\mathbb{E}\left[\bm{X}_{i,j}\bm{X}_{i,j}^{\top}\right]
	& =\sum\nolimits _{i,j}\mathbb{E}\left[ \big( p^{-1} \delta_{i,j} -1 \big)^{2}(M_{i,j}^{\star})^{2}\bm{e}_{i}\bm{e}_{j}^{\top}\bm{e}_{j}\bm{e}_{i}^{\top}\right]\\
 & =\frac{1-p}{p}\sum\nolimits _{i,j}\big(M_{i,j}^{\star}\big)^{2}\bm{e}_{i}\bm{e}_{i}^{\top}=\frac{1-p}{p}\sum_{i=1}^{n_{1}}\|\bm{M}_{i,\cdot}^{\star}\|_{2}^{2}\bm{e}_{i}\bm{e}_{i}^{\top} \\
	& \preceq \frac{1-p}{p}  \|\bm{M}^{\star}\|_{2,\infty}^{2} \bm{I}_{n_1} \preceq \frac{\mu r}{n_1p}  \|\bm{M}^{\star}\|^{2} \,\bm{I}_{n_1} .
\end{align*}
%
Here, the first identity arises from the definition of $\bm{X}_{i,j}$,
the second one calculates the variance of Bernoulli random variables,
and the last line relies on the upper bound~(\ref{eq:mc-row-norm-upper}).
Similarly, the second term in the variance statistic enjoys the following characterization:
%
\[
	\sum\nolimits _{i,j}\mathbb{E}\left[\bm{X}_{i,j}^{\top}\bm{X}_{i,j}\right] \preceq \frac{\mu r}{n_2p} \|\bm{M}^{\star}\|^{2} \,\bm{I}_{n_2}.
\]
%
Taking the above  relations together and recalling that $n_1\leq n_2$ give
%
\begin{align*}
v & %\leq\frac{1}{p}\max\left\{ \|\bm{M}^{\star}\|_{2,\infty}^{2},\|\bm{M}^{\star\top}\|_{2,\infty}^{2}\right\}
	% \leq\frac{1}{p}\left\{ \frac{\mu r}{n_{1}} \big\| \bm{M}^{\star} \big\|^2,  \frac{\mu r}{n_{2}} \big\| \bm{M}^{\star} \big\|^2 \right\} \\
	\leq\frac{\mu r}{n_{1}p} \big\| \bm{M}^{\star} \big\|^2.
\end{align*}
% 


With the above bounds in place, invoking matrix Bernstein (see Corollary~\ref{thm:matrix-Bernstein-friendly}) reveals that: with probability at least $1-O(n_2^{-10})$,
%
\begin{align*}
	\| \bm{E}\| & \lesssim \sqrt{\frac{\mu r \| \bm{M}^{\star} \|^2 \log n_{2}}{n_{1}p} }+\frac{\mu r \| \bm{M}^{\star} \| \log n_{2}}{p\sqrt{n_{1}n_{2}}}
	\asymp \sqrt{\frac{\mu r \| \bm{M}^{\star} \|^2 \log n_{2}}{n_{1}p} } ,
\end{align*}
%
where the last inequality is valid as long as  $n_{2}p \gtrsim \mu r\log n_{2}$.






\paragraph{Proof of Theorem~\ref{thm:CI-noisy-matrix-completion}.} 

Given that $\eta_{i,j}$ is a Gaussian random variable and hence possibly
unbounded, we find it convenient to introduce a truncated version
as follows
%
\[
	\widetilde{\eta}_{i,j}=\eta_{i,j}\mathbbm{1}\big\{|\eta_{i,j}|\leq 5\sigma_{\eta}\sqrt{\log n}\big\},\qquad1\leq i,j\leq n.
\]
%
and 
%
\[
	\widetilde{M}_{i,j} = \begin{cases} \frac{1}{p} \big( M_{i,j}^{\star} + \widetilde{\eta}_{i,j} \big), \qquad & \text{if } (i,j)\in \Omega, \\ 0, &\text{else}. \end{cases}
\]
%
Repeating the analysis in Section~\ref{sec:proof-eqn-noise-matrix-E-bound-denoising}, we can show that 
%
\[
	\mathbb{P}\big\{ \bm{M} = \widetilde{\bm{M}} \big\}
	= \mathbb{P} \big\{ \eta_{i,j} = \widetilde{\eta}_{i,j}, \forall i,j\in [n] \big\} 
	\geq 1- n^{-10},
\]
%
meaning that  $\bm{M}$ and $\widetilde{\bm{M}}$ are equivalent with high probability. 
As a result, we shall concentrate on validating the confidence interval computed based on $\widetilde{\bm{M}}$ in the subsequent analysis. %, unless otherwise noted 
Before proceeding, we record several key properties about  $\widetilde{\eta}_{i,j}$ as follows:
%
\begin{align}
	\mathbb{E}[ \widetilde{\eta}_{i,j}] = 0, ~~~ \mathbb{E}[ \widetilde{\eta}^2_{i,j}]=(1-o(1)) \sigma_{\eta}^2, 
	 ~~~ \big|\widetilde{\eta}_{i,j}\big|\leq  5\sigma_{\eta}\sqrt{\log n}.  
	 \label{eq:eta-property-CI-MC}
\end{align}
%
 


The proof follows by invoking Theorem~\ref{thm:CI-general}, as long as the conditions required therein are satisfied. 
To begin with, the associated variance parameters are given by
%
\begin{align*}
\sigma_{i,j}^{2} & \coloneqq\mathbb{E}\left[\big(\widetilde{M}_{i,j}-M_{i,j}^{\star}\big)^{2}\right]\\
 & =p\mathbb{E}\left[\left(\frac{1-p}{p}M_{i,j}^{\star}+\frac{1}{p}\widetilde{\eta}_{i,j}\right)^{2}\right]+(1-p)\big(M_{i,j}^{\star}\big)^{2}\\
 & = 
\frac{\left(1-p\right)^{2}}{p} \big(M_{i,j}^{\star}\big)^{2}+\frac{1}{p}\mathbb{E}\left[\widetilde{\eta}_{i,j}^{2}\right]+(1-p)\big(M_{i,j}^{\star}\big)^{2} \\
 & =\frac{1-p}{p}\big(M_{i,j}^{\star}\big)^{2}+\frac{1-o(1)}{p}\sigma_{\eta}^{2}, 
\end{align*}
%
thus leading to
%
\begin{align*}
\sigma_{\min}^{2} & \coloneqq \min_{i,j}\sigma_{i,j}^{2}=\frac{1-p}{p}\min_{i,j}\big(M_{i,j}^{\star}\big)^{2}+\frac{1-o(1)}{p}\sigma_{\eta}^{2},\\
\sigma^{2} & \coloneqq \max_{i,j}\sigma_{i,j}^{2}=\frac{1-p}{p} \|\bm{M}^{\star}\|_{\infty}^2 +\frac{1-o(1)}{p}\sigma_{\eta}^{2}.
\end{align*}
%
Apparently, $\sigma^2/\sigma_{\min}^{2}=O(1)$ holds true under the assumptions of Theorem~\ref{thm:CI-noisy-matrix-completion}. 
In addition, the random variables $\{\widetilde{M}_{i,j}-M_{i,j}^{\star}\}$ are all bounded obeying
%
\[
\big|\widetilde{M}_{i,j}-M_{i,j}^{\star}\big|\leq\frac{1-p}{p}\big|M_{i,j}^{\star}\big|+\frac{1}{p}\big|\widetilde{\eta}_{i,j}\big|
\leq\frac{(1-p)\|\bm{M}^{\star}\|_{\infty}+\sigma_{\eta}\sqrt{5\log n}}{p}\eqqcolon B.
\]
%
These bounds readily imply that
%
\begin{equation}
	\frac{B}{\sigma}\asymp\frac{~\frac{(1-p)\|\bm{M}^{\star}\|_{\infty}+\sigma_{\eta}\sqrt{5\log n}}{p}~}{\sqrt{\frac{1-p}{p}}\|\bm{M}^{\star}\|_{\infty}+\frac{1}{\sqrt{p}}\sigma_{\eta}}\lesssim\sqrt{\frac{\log n}{p}}. 
	\label{eq:B-sigma-ratio-noisy-MC}
\end{equation}
%



Moving to the condition $\sigma \kappa\sqrt{n\log n}\lesssim |\lambda_{r}^{\star}|$ in Theorem~\ref{thm:CI-noisy-matrix-completion}, it can be guaranteed if
	%
	\begin{equation*}
		\|\bm{M}^{\star}\|_{\infty}\kappa\sqrt{\frac{(1-p)n\log n}{p}}
		\lesssim|\lambda_{r}^{\star}| 
		\quad \text{and} \quad
		\sigma_{\eta}\kappa\sqrt{\frac{n\log n}{p}}\lesssim|\lambda_{r}^{\star}|.  
		\label{eq:Mstar-noise-requirement-noisy-MC-135}
	\end{equation*}
	%
Given the assumption  $\frac{\max_{k,l}|M_{k,l}^{\star}|}{\min_{k,l}|M_{k,l}^{\star}|}=O(1)$, one has 
%
\begin{equation}
	\label{eq:Mstar-inf-F-relation-noisy-MC}
	\|\bm{M}^{\star}\|_{\infty} \asymp \frac{1}{n} \|\bm{M}^{\star}\|_{\mathrm{F}} 
	\leq \frac{\sqrt{r}}{n} \|\bm{M}^{\star}\| 
	= \frac{\kappa \sqrt{r}}{n} |\lambda_{r}^{\star}|. 
\end{equation}
%
As a consequence, the condition $\sigma \kappa\sqrt{n\log n}\lesssim |\lambda_{r}^{\star}|$ can be ensured under Condition~\eqref{eq:Mstar-noise-requirement-noisy-MC}. 
%


It remains to certify Condition~\eqref{eq:Uj-Ui-lower-bound-thm}. 
By virtue of the above calculations of $\sigma$ and $B$ as well as the property \eqref{eq:B-sigma-ratio-noisy-MC}, it is easily seen that Condition~\eqref{eq:Uj-Ui-lower-bound-thm} is valid as long as the following holds:
%
\begin{align*}
\frac{\big\|\bm{U}_{j,\cdot}^{\star}\big\|_{2}^{2}+\big\|\bm{U}_{i,\cdot}^{\star}\big\|_{2}^{2}}{\big\|\bm{U}^{\star}\big\|_{\mathrm{F}}^{2}} & \gtrsim\frac{\kappa^{2}\mu^{2}r^{2}\log^{5/2}n}{n\sqrt{np}}+\frac{\sqrt{1-p}\,\|\bm{M}^{\star}\|_{\infty}\mu^{2}r\kappa^{3}\log^{3}n}{|\lambda_{r}^{\star}|\sqrt{np}}\\
 & \qquad+\frac{\sigma_{\eta}\mu^{2}r\kappa^{3}\log^{3}n}{|\lambda_{r}^{\star}|\sqrt{np}}.
\end{align*}
%
Taking this together with the relation \eqref{eq:Mstar-inf-F-relation-noisy-MC}, 
we can demonstrate straightforwardly that Condition~\eqref{eq:Uj-Ui-lower-bound-thm} is guaranteed to hold as long as Condition~\eqref{eq:Ui-Uj-lower-bound-thm-noisy-MC-146} is satisfied. 
This completes the proof. 




%\paragraph{Statistical guarantees.} 


%\begin{subappendices}














%%%%%%%%%%%%%%%%%%%%%%%%%%%%%%%%%%%%%%%%%%%%%%%%%%%%%%%%%%%%%%%%%%%%%%



\section{Appendix A: Proof of Theorem~\ref{thm:UsgnH-Ustar-MUstar-general}}
\label{sec:proof-thm:UsgnH-Ustar-MUstar-general}

%In the rest of this chapter, we establish the results Theorem~\ref{thm:UsgnH-Ustar-MUstar-general} and Corollary~\ref{cor:entrywise-error-general}.
To simplify notation, we assume throughout the proof that $\lambda_r^{\star} > 0$, namely,
%
\begin{align}
	|\lambda_1^{\star}| \geq  \cdots \geq | \lambda_{r-1}^{\star} | \geq \lambda_r^{\star} > 0.
	\label{eq:assumption-lambda-r-positive}
\end{align}
%Now we turn to the proof of Theorem~\ref{thm:UsgnH-Ustar-MUstar-general}.
The challenge of the proof arises due to the complicated statistical dependency between $\bm{M}$ and $\bm{U}$, and the leave-one-out analysis paves a plausible path to decouple the dependency.






\subsection{Construction of leave-one-out auxiliary estimates}
\label{sec:construction-LOO-general}

As elucidated in the rank-1 matrix denoising example in Section~\ref{sec:rank-1-denoising-LOO}, the key to enabling fine-grained analysis is to seek assistance from  a collection of leave-one-out estimates.  Akin to Section~\ref{sec:LOO-estimates-denoising}, for each $1\leq l\leq n$, we construct  two auxiliary matrices  $\bm{M}^{(l)}$ and $\bm{E}^{(l)}=\big[E^{(l)}_{i,j} \big]_{1\leq i,j\leq n}$ as follows:
%
\begin{align}
\bm{M}^{(l)}\coloneqq\bm{M}^{\star}+\bm{E}^{(l)},
\quad
E_{i,j}^{(l)}  \coloneqq
\begin{cases}
E_{i,j}, & \text{if }i\neq l\text{ and }j\neq l,\\
0, & \text{else},
\end{cases}
\label{eq:Ml-construction-general}
\end{align}
%
which are generated by simply discarding all random noise incurred in the $l$-th column/row of the data matrix.
In addition, let $\lambda_1^{(l)},\cdots,\lambda_n^{(l)}$ be the eigenvalues of $\bm{M}^{(l)}$ sorted by
%
\begin{align}
	\big|\lambda_1^{(l)}\big| \geq \big|\lambda_2^{(l)}\big| \geq \cdots \geq \big| \lambda_{n}^{(l)} \big| ,
\end{align}
%
and denote by $\bm{u}_{i}^{(l)}$ the eigenvector of $\bm{M}^{(l)}$  associated with $\lambda_i^{(l)}$. The leave-one-out spectral estimates $\bm{U}^{(l)}$ and $\bm{\Lambda}^{(l)}$ are, therefore, given by
%
\begin{align}
	\bm{U}^{(l)}\coloneqq\big[\bm{u}_{1}^{(l)},\cdots,\bm{u}_{r}^{(l)}\big]\in\mathbb{R}^{n\times r};
	\quad
	\bm{\Lambda}^{(l)}\coloneqq\mathsf{diag}\big(\big[\lambda_{1}^{(l)},\cdots,\lambda_{r}^{(l)}\big]\big).	
\end{align}
%

We emphasize again that the main advantage of introducing the leave-one-out estimate $\bm{U}^{(l)}$ stems from its statistical independence from the $l$-th row of $\bm{M}$, which substantially simplifies the analysis for the $l$-th row of the estimate.
In principle, our analysis employs the leave-one-out estimates
to help decouple delicate statistical dependency  in a row-by-row fashion. Another crucial aspect of the analysis lies in the exploitation of the proximity of all these auxiliary estimates, a feature that is enabled by the ``stability'' of the spectral method.





\subsection{Preliminary facts}
\label{sec:preliminary-facts-Linf-general}

%Before embarking on  the $\ell_{\infty}/\ell_{2,\infty}$ analysis, we gather a couple of useful facts immediately available from the $\ell_2$
%perturbation theory and the matrix tail bounds. In what follows, $\bm{\Theta}$ (resp.~$\bm{\Theta}^{(l)}$) denotes a diagonal matrix whose diagonal entries are the principal angles between $\bm{U}$ (resp.~$\bm{U}^{(l)}$) and $\bm{U}^{\star}$.
%%
%\begin{lemma}
%\label{lem:preliminary-result-general-iid}
%Consider the setting in Section~\ref{sec:setup-general-theory-independent}.
%There is some constant $c_2>0$ such that with probability at least $1-O(n^{-7})$,
%%
%\begin{subequations}
%\label{eq:L2-perturbation-summary-general}
%%
%\begin{align}
%	\|\bm{E} \|  &\leq c_2\sigma \sqrt{n}  \quad  & \big\|\bm{E}^{(l)} \big\| &\leq c_2 \sigma \sqrt{n}  &&  \label{eq:noise-size-general} \\
%  \mathsf{dist}( \bm{U} , \bm{U}^{\star} )  &\leq \frac{2c_2\sigma\sqrt{n}}{\lambda_r^{\star}}  ~~
% & \mathsf{dist}\big( \bm{U}^{(l)} , \bm{U}^{\star} \big)  &\leq \frac{2c_2\sigma\sqrt{n}}{\lambda_r^{\star}} && \label{eq:dist-U-Ustar-general} \\
%	  \big\|\sin \bm{\Theta} \big\|  &\leq \frac{c_2\sigma\sqrt{2n}}{\lambda_r^{\star}}  ~~
%	& \big\|\sin \bm{\Theta}^{(l)} \big\|  &\leq \frac{c_2\sigma\sqrt{2n}}{\lambda_r^{\star}} && \label{eq:distp-U-Ustar-general} \\
%	\max_{1\leq j\leq r} | \lambda_j  |  &\geq \lambda_r^{\star} - c_2 \sigma \sqrt{n} &  \max_{1\leq j\leq r} | \lambda^{(l)}_j  |  &\geq \lambda_r^{\star} - c_2\sigma \sqrt{n} &&  \label{eq:lambda-size-large-general}  \\
%	\max_{j: j> r} | \lambda_j   |  &\leq c_2 \sigma \sqrt{n} & \max_{j: j> r}| \lambda_j^{(l)} |  &\leq c_2\sigma \sqrt{n}  \label{eq:lambda-size-small-general}
%\end{align}
%%
%%
%hold simultaneously for all $1\leq l\leq n$, provided that $c_2\sigma \sqrt{n}\leq (1-1/\sqrt{2})\lambda_r^{\star}$.
%Moreover, for any fixed matrix $\bm{A}\in \mathbb{R}^{n\times d}$ with $d\leq n$, one has
%%
%\begin{align}
%	\|\bm{E} \bm{A} \|_{2,\infty} \leq
%	4\sigma\sqrt{\log n}\,\|\bm{A}\|_{\mathrm{F}}+(6B\log n)\|\bm{A}\|_{2,\infty}
%	\label{eq:E-concentration-inf2-general}
%\end{align}
%\end{subequations}
%%
%with probability exceeding $1-2n^{-5}$.
%\end{lemma}
%%
%\begin{remark}
%In view of \eqref{eq:E-concentration-inf2-general}, with probability at least $1-2n^{-5}$,
%%
%\begin{align}
%	& \|\bm{E}\bm{U}^{\star}\|_{2,\infty} \leq 4\sigma\sqrt{\log n}\,\|\bm{U}^{\star}\|_{\mathrm{F}}+(6B\log n)\|\bm{U}^{\star}\|_{2,\infty}
%		\nonumber \\
% 	&\qquad = 4\sigma\sqrt{r\log n}+6B\sqrt{\frac{\mu r\log^{2}n}{n}}
%	\leq (4+6c_{\mathsf{b}} ) \sigma \sqrt{r\log n},
%	\label{eq:EU-2-inf-bound-general}
%\end{align}
%%
%which relies on the definition \eqref{eq:defn-Ustar-incoherence} and \eqref{eq:assumption-B-sigma} in Assumption~\ref{assumption-noise-general}.  As a result,
%%
%\begin{align}
%\big\|\bm{M}\bm{U}^{\star}\big\|_{2,\infty} & \leq\big\|\bm{M}^{\star}\bm{U}^{\star}\big\|_{2,\infty}+\big\|\bm{E}\bm{U}^{\star}\big\|_{2,\infty} \notag\\
% % \leq\big\|\bm{U}^{\star}\big\|_{2,\infty}\|\bm{\Lambda}^{\star}\|+\big\|\bm{E}\bm{U}^{\star}\big\|_{2,\infty} \notag\\
% & \leq \sqrt{\frac{\mu r}{n}}\big|\lambda_{1}^{\star}\big|+(4+6c_{\mathsf{b}} )\sigma\sqrt{r\log n},
%	\label{eq:MU-2-inf-bound-general}
%\end{align}
%%
%which follows from the fact $\bm{M}^{\star}\bm{U}^{\star}=\bm{U}^{\star}\bm{\Lambda}^{\star}$ (so that $\|\bm{M}^{\star}\bm{U}^{\star}\|_{2,\infty}\leq \big\|\bm{U}^{\star}\big\|_{2,\infty}\|\bm{\Lambda}^{\star}\| = \sqrt{\mu r/n} \,|\lambda_1^{\star}|$).
%\end{remark}
%
%
%Furthermore, we find it helpful to introduce the following  matrices
%%
%\begin{equation}
%	\bm{H} \coloneqq \bm{U}^{\top}\bm{U}^{\star}
%	\qquad\text{and}\qquad
%	\bm{H}^{(l)} \coloneqq \bm{U}^{(l)\top}\bm{U}^{\star},
%	\label{eq:defn-H-Hl-general}
%\end{equation}
%%
%which turn out to be close to being orthonormal.
%%
%Our results concerning $\bm{H}$ and $\bm{H}^{(l)}$ are formalized as follows.
%%
%\begin{lemma}
%\label{lem:H-property-summary-general}
%Instate the assumptions of  Lemma~\ref{lem:preliminary-result-general-iid}. With probability at least $1-O(n^{-7})$, the following holds:
%%
%\begin{subequations}  \label{eq:H-property-summary-general}
%%
%\begin{align}
%	\|\bm{H}^{-1}\|  &\leq 2  \quad  & \big\|\big(\bm{H}^{(l)} \big) ^{-1} \big\| &\leq 2  &&  \label{eq:H-inv-norm-bound-general} \\
%  \big\|\bm{H}-\mathsf{sgn}(\bm{H})\big\|  &\leq  \frac{2c_{2}^{2}n \sigma^{2}}{(\lambda_{r}^{\star})^{2}}  ~~
%	&      \big\|\bm{H}^{(l)}-\mathsf{sgn}(\bm{H}^{(l)})\big\|  &  \leq  \frac{2c_{2}^{2}n\sigma^{2} }{(\lambda_{r}^{\star})^{2}}   &&
%	\label{eq:H-sgnH-dist-general}
%\end{align}
%%
%\end{subequations}
%%
%simultaneously for all $1\leq l\leq n$ and, in addition,
%%
%\begin{align}
%  \|\bm{U}\bm{H}-\bm{U}^{\star}\|_{\mathrm{F}}  \leq\frac{2c_{2}\sigma \sqrt{rn}}{\lambda_{r}^{\star}} .
%\end{align}
%%
%
%
%\end{lemma}
%
%\begin{remark}\label{remark:AH-norm-bound}
%As a consequence of \eqref{eq:H-property-summary-general} and Proposition~\ref{prop:unitary_norm_relation},  one has
%%
%\begin{subequations}  \label{eq:A-2inf-H-Hl-general}
%\begin{align}
%	\vertiii{\bm{A}} &= \vertiiibig{ \bm{A}\bm{H} \bm{H} ^{-1}}
%	\leq \vertiii{ \bm{A}\bm{H}} \, \big\| \bm{H} ^{-1}\big\|
%	\leq 2 \vertiii{ \bm{A}\bm{H} }
%	\label{eq:A-2inf-H-general} \\
%	\vertiii{\bm{A}}
%	& \leq \vertiiibig{ \bm{A}\bm{H}^{(l)}} \, \big\| \big( \bm{H}^{(l)}  \big) ^{-1}\big\|
%  	\leq 2 \vertiiibig{ \bm{A}\bm{H}^{(l)} }
%	\label{eq:A-2inf-Hl-general}
%\end{align}
%\end{subequations}
%%
%for any matrix $\bm{A}$ and any unitarily invariant norm $\vertiii{\cdot}$ (e.g., the Frobenius norm  and the $\ell_{2,\infty}$ norm $\|\cdot\|_{2,\infty}$).
%\end{remark}
%
%
%
%The proofs for the above results  are postponed to Section~\ref{sec:proof-auxiliary-lemmas-general}.
%
%
Before embarking on the $\ell_{\infty}$ and $\ell_{2,\infty}$ analyses,
we gather a couple of useful facts, whose proofs are postponed to
Section~\ref{sec:proof-auxiliary-lemmas-general}. In what follows,
$\bm{\Theta}$ (resp.~$\bm{\Theta}^{(l)}$) denotes a diagonal matrix
whose diagonal entries are the principal angles between $\bm{U}$
(resp.~$\bm{U}^{(l)}$) and $\bm{U}^{\star}$. In addition, we find
it helpful to introduce the following matrices
\begin{equation}
\bm{H}\coloneqq\bm{U}^{\top}\bm{U}^{\star}\qquad\text{and}\qquad\bm{H}^{(l)}\coloneqq\bm{U}^{(l)\top}\bm{U}^{\star},\label{eq:defn-H-Hl-general}
\end{equation}
which turn out to be close to being orthonormal.

The first set of results follows from the statistical nature of the
perturbation matrix $\bm{E}$ (cf.~Assumption~\ref{assumption-noise-general}),
which is immediately available from the matrix tail bounds. 
%
\begin{lemma}\label{lemma:general-noise-bound}
Consider the setting in Section~\ref{sec:setup-general-theory-independent}.
There is some constant $c_{2}>0$ such that with probability at least
$1-O(n^{-7})$,
\begin{equation}
\max_{l}\big\|\bm{E}^{(l)}\big\|\leq\|\bm{E}\|\leq c_{2}\sigma\sqrt{n}.\label{eq:noise-size-general}
\end{equation}
Moreover, for any fixed matrix $\bm{A}\in\mathbb{R}^{n\times d}$
with $d\leq n$, one has
\begin{align}
\|\bm{E}\bm{A}\|_{2,\infty}\leq4\sigma\sqrt{\log n}\,\|\bm{A}\|_{\mathrm{F}}+(6B\log n)\|\bm{A}\|_{2,\infty}.\label{eq:E-concentration-inf2-general}
\end{align}
 \end{lemma}
 %
%with probability at least $1-2n^{-5}$.

\begin{remark} In view of \eqref{eq:E-concentration-inf2-general},
with probability at least $1-2n^{-5}$,
\begin{align}
 & \|\bm{E}\bm{U}^{\star}\|_{2,\infty}\leq4\sigma\sqrt{\log n}\,\|\bm{U}^{\star}\|_{\mathrm{F}}+(6B\log n)\|\bm{U}^{\star}\|_{2,\infty}\nonumber \\
 & \qquad=4\sigma\sqrt{r\log n}+6B\sqrt{\frac{\mu r\log^{2}n}{n}}
 = (4+6c_{\mathsf{b}})\sigma\sqrt{r\log n},\label{eq:EU-2-inf-bound-general}
\end{align}
which relies on the definition \eqref{eq:defn-Ustar-incoherence}
and the definition of $c_{\mathsf{b}}$ in \eqref{eq:assumption-B-sigma}.
As a result,
\begin{align}
\big\|\bm{M}\bm{U}^{\star}\big\|_{2,\infty} & \leq\big\|\bm{M}^{\star}\bm{U}^{\star}\big\|_{2,\infty}+\big\|\bm{E}\bm{U}^{\star}\big\|_{2,\infty}\nonumber \\
&\leq \sqrt{\frac{\mu r}{n}}|\lambda_1^{\star}| + (4+6c_{\mathsf{b}})\sigma\sqrt{r\log n}, \label{eq:MU-2-inf-bound-general}
\end{align}
which follows from the fact $\bm{M}^{\star}\bm{U}^{\star}=\bm{U}^{\star}\bm{\Lambda}^{\star}$
(so that $\|\bm{M}^{\star}\bm{U}^{\star}\|_{2,\infty}\leq\big\|\bm{U}^{\star}\big\|_{2,\infty}\|\bm{\Lambda}^{\star}\|=\sqrt{\mu r/n}\,|\lambda_{1}^{\star}|$).
\end{remark}

With the size of the perturbations (i.e., $\|\bm{E}^{(l)}\|$ and
$\|\bm{E}\|$) under control, the $\ell_{2}$ perturbation theory
established in Chapter~\ref{cha:matrix-perturbation} leads to the following set of conclusions.
\begin{lemma}\label{lem:preliminary-result-general-iid} Suppose
that $c_{2}\sigma\sqrt{n}\leq(1-1/\sqrt{2})\lambda_{r}^{\star}$,
where $c_{2}$ is the same constant as in Lemma~\ref{lemma:general-noise-bound}.
Then with probability at least $1-O(n^{-7})$, one has \begin{subequations}
\label{eq:L2-perturbation-summary-general}
\begin{align}
\mathsf{dist}(\bm{U},\bm{U}^{\star}) & \leq\frac{2c_{2}\sigma\sqrt{n}}{\lambda_{r}^{\star}},~~ & \mathsf{dist}\big(\bm{U}^{(l)},\bm{U}^{\star}\big) & \leq\frac{2c_{2}\sigma\sqrt{n}}{\lambda_{r}^{\star}}, \label{eq:dist-U-Ustar-general}\\
\big\|\sin\bm{\Theta}\big\| & \leq\frac{c_{2}\sigma\sqrt{2n}}{\lambda_{r}^{\star}}, ~~ & \big\|\sin\bm{\Theta}^{(l)}\big\| & \leq\frac{c_{2}\sigma\sqrt{2n}}{\lambda_{r}^{\star}}, \label{eq:distp-U-Ustar-general}\\
\max_{1\leq j\leq r}|\lambda_{j}| & \geq\lambda_{r}^{\star}-c_{2}\sigma\sqrt{n}, & \max_{1\leq j\leq r}|\lambda_{j}^{(l)}| & \geq\lambda_{r}^{\star}-c_{2}\sigma\sqrt{n}, \label{eq:lambda-size-large-general}\\
\max_{j:j>r}|\lambda_{j}| & \leq c_{2}\sigma\sqrt{n}, & \max_{j:j>r}|\lambda_{j}^{(l)}| & \leq c_{2}\sigma\sqrt{n}\label{eq:lambda-size-small-general}
\end{align}
hold simultaneously for all $1\leq l\leq n$. In addition,
\begin{align}
\|\bm{U}\bm{H}-\bm{U}^{\star}\|_{\mathrm{F}}\leq\frac{2c_{2}\sigma\sqrt{rn}}{\lambda_{r}^{\star}}.\label{eq:dist-under-H-general}
\end{align}
 \end{subequations} %
with probability exceeding $1-2n^{-5}$. \end{lemma}


\begin{remark}\label{remark:perturbation-size-eigen-gap}
Lemmas~\ref{lemma:general-noise-bound} and \ref{lem:preliminary-result-general-iid} allow us to bound the eigengap and perturbation size as follows
%
\begin{subequations}\label{eq:lambda-rl-lambda-r1l-M-Ml-LB-step}
\begin{align}
	\big| \lambda_{r}^{(l)} \big| - \big|\lambda_{r+1}^{(l)}\big|
	& \geq \lambda_{r}^{\star}-  c_2 \sigma \sqrt{n} - c_2 \sigma \sqrt{n}
	\geq \lambda_{r}^{\star}/2,
	\label{eq:lambda-rl-lambda-r1l-LB-step}  \\
	\|\bm{M}-\bm{M}^{(l)}\| & \leq\|\bm{M}-\bm{M}^{\star}\|+\|\bm{M}^{\star}-\bm{M}^{(l)}\|=\|\bm{E}\|+\|\bm{E}^{(l)}\|  \notag\\
	& \leq2c_{2}\sigma\sqrt{n}\leq(1-1/\sqrt{2}) \big( \big| \lambda_{r}^{(l)} \big| - \big|\lambda_{r+1}^{(l)}\big| \big),
	\label{eq:M-Ml-size-step}
\end{align}
%
\end{subequations}
%
which are valid as long as $20 c_2 \sigma \sqrt{n} \leq \lambda_r^{\star}$. These will prove useful when bounding the approximation error of $\bm{U}$ using $\bm{U}^{(l)}$.
\end{remark}

Another collection of results is concerned with $\bm{H}$ and $\bm{H}^{(l)}$.
%
\begin{lemma} \label{lem:H-property-summary-general} Suppose that the
assumptions of Lemma~\ref{lem:preliminary-result-general-iid} hold. With
probability at least $1-O(n^{-7})$, \begin{subequations}
\label{eq:H-property-summary-general}
\begin{align}
\|\bm{H}^{-1}\| & \leq2,\quad & \big\|\big(\bm{H}^{(l)}\big)^{-1}\big\| & \leq2,\label{eq:H-inv-norm-bound-general}\\
\big\|\bm{H}-\mathsf{sgn}(\bm{H})\big\| & \leq\frac{2c_{2}^{2}n\sigma^{2}}{(\lambda_{r}^{\star})^{2}},~~ & \big\|\bm{H}^{(l)}-\mathsf{sgn}(\bm{H}^{(l)})\big\| & \leq\frac{2c_{2}^{2}n\sigma^{2}}{(\lambda_{r}^{\star})^{2}}\label{eq:H-sgnH-dist-general}
\end{align}
\end{subequations} hold simultaneously for all $1\leq l\leq n$.

\end{lemma}

\begin{remark}\label{remark:AH-norm-bound} As a consequence of \eqref{eq:H-property-summary-general}
and Proposition~\ref{prop:unitary_norm_relation}, one has \begin{subequations}
\label{eq:A-2inf-H-Hl-general}
\begin{align}
\vertiii{\bm{A}} & =\vertiiibig{\bm{A}\bm{H}\bm{H}^{-1}}\leq\vertiii{\bm{A}\bm{H}}\,\big\|\bm{H}^{-1}\big\|\leq2\vertiii{\bm{A}\bm{H}}, \label{eq:A-2inf-H-general}\\
\vertiii{\bm{A}} & \leq\vertiiibig{\bm{A}\bm{H}^{(l)}}\,\big\|\big(\bm{H}^{(l)}\big)^{-1}\big\|\leq2\vertiiibig{\bm{A}\bm{H}^{(l)}}\label{eq:A-2inf-Hl-general}
\end{align}
\end{subequations} for any matrix $\bm{A}$. Here, $\vertiii{\bm{A}}$ could either be the Frobenius norm or the $\ell_{2,\infty}$
norm $\|\cdot\|_{2,\infty}$. \end{remark}




\subsection{Leave-one-out analysis}

Now we move on to the main part of the analysis, which is further decomposed into four steps. 

\subsubsection{Step 1: decomposing the $\ell_{2,\infty}$ estimation error of $\bm{U}$}




By virtue of the proximity of  $\bm{H}$ and $\mathsf{sgn}(\bm{H})$ unveiled in Lemma~\ref{lem:H-property-summary-general}, we are allowed to employ $\bm{U}\bm{H}$ as a surrogate for $\bm{U}\mathsf{sgn}(\bm{H})$, which is more convenient to work with. As it turns out, the discrepancy between $\bm{U}\bm{H}$ and the first-order approximation $\bm{M}\bm{U}^{\star}(\bm{\Lambda}^{\star})^{-1}$, and the discrepancy between $\bm{U}\bm{H}$ and the truth,  can be bounded by three important terms, as asserted below. The proof is built upon elementary  algebra and basic $\ell_2$ perturbation bounds in Section~\ref{sec:preliminary-facts-Linf-general}, and is deferred to Section~\ref{sec:proof-auxiliary-lemmas-general}.




\begin{lemma}
\label{lem:UH-MUstar-diff-decompose}
%
Suppose that $2c_2\sigma\sqrt{n}\leq \lambda_r^{\star}$ for some sufficiently large constant $c_2>0$. Then with probability at least $1-O(n^{-7})$, one has
%
\begin{subequations}
	\label{eq:UH-MUstar-Ustar-diff-2inf-general}
	\begin{align}
		\big\|\bm{U}\bm{H}-\bm{M}\bm{U}^{\star}\big(\bm{\Lambda}^{\star}\big)^{-1}\big\|_{2,\infty} & \leq\mathcal{E}_{1}+\mathcal{E}_{2},
		\label{eq:UH-MUstarLambdastar-diff-2inf-general}\\
		\big\|\bm{U}\bm{H}-\bm{U}^{\star}\big\|_{2,\infty} & \leq \mathcal{E}_{1}+\mathcal{E}_{2}+ \mathcal{E}_3,
		\label{eq:UH-Ustar-diff-2inf-general}
	\end{align}
\end{subequations}
%
where
%
%\begin{subequations}
%\begin{align*}
%	\label{eq:UB-UH-MULambda-UB-lemma}
%
%\begin{array}{cc}
$\mathcal{E}_{1}\coloneqq\frac{2\|\bm{M}(\bm{U}\bm{H}-\bm{U}^{\star})\|_{2,\infty}}{\lambda_{r}^{\star}}$, $\mathcal{E}_{2}\coloneqq\frac{4\|\bm{M}\bm{U}^{\star}\|_{2,\infty}\|\bm{E}\|}{(\lambda_{r}^{\star})^{2}}$, and $\mathcal{E}_{3}\coloneqq\frac{\|\bm{E}\bm{U}^{\star}\|_{2,\infty}}{\lambda_{r}^{\star}}$.
%\end{array}
%
%\end{align*}
%\end{subequations}
%
\end{lemma}
%




Lemma~\ref{lem:UH-MUstar-diff-decompose} leaves us with three important terms to deal with. The term $\mathcal{E}_1$ is most complicated as it involves the product of two random matrices, whereas $\mathcal{E}_2$ and $\mathcal{E}_3$ can be controlled straightforwardly through our preliminary facts in Section~\ref{sec:preliminary-facts-Linf-general}.   Specifically,   the term $\mathcal{E}_2$ can be bounded by combining \eqref{eq:MU-2-inf-bound-general} with the bound \eqref{eq:noise-size-general} on $\bm{E}$ to obtain
%
% \begin{align*}
% \big\|\bm{M}\bm{U}^{\star}\big\|_{2,\infty} & \leq\big\|\bm{M}^{\star}\bm{U}^{\star}\big\|_{2,\infty}+\big\|\bm{E}\bm{U}^{\star}\big\|_{2,\infty}
%   \leq\big\|\bm{U}^{\star}\big\|_{2,\infty}\|\bm{\Lambda}^{\star}\|+\big\|\bm{E}\bm{U}^{\star}\big\|_{2,\infty}\\
%  & \leq \sqrt{\frac{\mu r}{n}}\big|\lambda_{1}^{\star}\big|+(4+6c_{1})\sigma\sqrt{r\log n},
% \end{align*}
%
% where the first identity makes use of the fact $\bm{M}^{\star}\bm{U}^{\star}=\bm{U}^{\star}\bm{\Lambda}^{\star}$ (so that $\|\bm{M}^{\star}\bm{U}^{\star}\|_{2,\infty}\leq \big\|\bm{U}^{\star}\big\|_{2,\infty}\|\bm{\Lambda}^{\star}\|$), and the last line relies on \eqref{eq:EU-2-inf-bound-general}. This combined with
% the bound \eqref{eq:L2-perturbation-summary-general} on $\|\bm{E}\|$ yields
%
\begin{align}
	% \frac{\big\|\bm{M}\bm{U}^{\star}\big\|_{2,\infty}\|\bm{E}\|}{(\lambda_{r}^{\star})^{2}}
	\mathcal{E}_2 & \leq \frac{4c_{2}\kappa\sigma\sqrt{\mu r}}{\lambda_{r}^{\star}}+\frac{4c_{2}(4+6c_{\mathsf{b}})\sigma^{2}\sqrt{rn\log n}}{(\lambda_{r}^{\star})^{2}}.
	\label{eq:MUstar-E-norm-lambda2-bound}
\end{align}
%
Regarding the term $\mathcal{E}_3$, the inequality \eqref{eq:EU-2-inf-bound-general} readily gives
%
\begin{align}
	%\big\|\bm{E}\bm{U}^{\star}\big(\bm{\Lambda}^{\star}\big)^{-1}\big\|_{2,\infty}\leq\frac{\big\|\bm{E}\bm{U}^{\star}\big\|_{2,\infty}}{\lambda_{r}^{\star}}
	\mathcal{E}_3 \leq\frac{(4+6c_{\mathsf{b}} )\sigma\sqrt{r\log n}}{\lambda_{r}^{\star}}.
	\label{eq:E3-bound-general}
\end{align}
%




Turning to controlling the remaining term $\mathcal{E}_{1}$, a closer inspection, however, reveals substantial challenges, due to the complicated statistical dependency between $\bm{M}$ and $\bm{U}$. To further complicate matters, the term $\mathcal{E}_1$---as we shall demonstrate momentarily---depend on some intrinsic properties of interest about $\bm{U}$ (e.g., $\|\bm{U}\bm{H}-\bm{U}^{\star}\|_{2,\infty}$),  which might lead to circular reasoning if not handled properly. In order to circumvent this issue, we intend to establish the following relation
%
\begin{equation}
	\mathcal{E}_{1}\leq\mathcal{E}_{1,1}+\rho_{1}\big\|\bm{U}\bm{H}-\bm{U}^{\star}\big\|_{2,\infty}
	\label{eq:E1-UB-E11-rho1}
\end{equation}
%
for some quantity $\mathcal{E}_{1,1}>0$ that does not involve $\|\bm{U}\bm{H}-\bm{U}^{\star} \|_{2,\infty}$  as well as some contraction factor $0<\rho_{1}\leq1/2$. Assuming the relation~\eqref{eq:E1-UB-E11-rho1} holds for the moment, we have the following useful claim (the proof is straightforward and again postponed to Section~\ref{sec:proof-auxiliary-lemmas-general}).
%
\begin{lemma}
\label{lem:UH-sequence-bounds-E123}
If Conditions \eqref{eq:UH-MUstar-Ustar-diff-2inf-general} and \eqref{eq:E1-UB-E11-rho1} hold with $0<\rho_{1}\leq1/2$, then we have
%
\begin{subequations}\label{eq:UH-sequence-bounds-E123}
\begin{align}
\big\|\bm{U}\bm{H}-\bm{U}^{\star}\big\|_{2,\infty} & \leq2\big(\mathcal{E}_{1,1}+\mathcal{E}_{2}+\mathcal{E}_{3}\big), \label{eq:UH-Ustar-E11-E2-E3}\\
\big\|\bm{U}\bm{H}-\bm{M}\bm{U}^{\star}\big(\bm{\Lambda}^{\star}\big)^{-1}\big\|_{2,\infty} & \leq2\big(\mathcal{E}_{1,1}+\mathcal{E}_{2}+\rho_{1}\mathcal{E}_{3}\big),
\label{eq:UH-MUstar-Lambdastar-E11-E2-E3}
\end{align}
\begin{align}
	& \big\|\bm{U}\mathsf{sgn}(\bm{H})-\bm{U}^{\star}\big\|_{2,\infty}  \leq4\big(\mathcal{E}_{1,1}+\mathcal{E}_{2}+\mathcal{E}_{3}\big)+\frac{4c_{2}^{2}\sigma^{2}\sqrt{\mu rn}}{(\lambda_{r}^{\star})^{2}}, \label{eq:UsgnH-Ustar-Lambdastar-E11-E2-E3} \\
	& \big\|\bm{U}\mathsf{sgn}(\bm{H})-\bm{M}\bm{U}^{\star}\big(\bm{\Lambda}^{\star}\big)^{-1}\big\|_{2,\infty} \nonumber \\
& \qquad \leq3\mathcal{E}_{1,1}+3\mathcal{E}_{2}+\Big(2\rho_{1}+\frac{8c_{2}^{2}\sigma^{2}n}{(\lambda_{r}^{\star})^{2}}\Big)\mathcal{E}_{3}
 +\frac{4c_{2}^{2}\sigma^{2}\sqrt{\mu rn}}{(\lambda_{r}^{\star})^{2}}.\label{eq:UsgnH-MUstar-Lambdastar-E11-E2-E3}
\end{align}
\end{subequations}
%
\end{lemma}
%
\begin{remark}
When $\mathcal{E}_3$ is the dominant term, the  bound   \eqref{eq:UH-MUstar-Lambdastar-E11-E2-E3}  might be stronger than \eqref{eq:UH-Ustar-E11-E2-E3} if $\rho_1$ is  small. % This subtle improvement proves beneficial for some applications like community detection.
\end{remark}


% Let us take a closer inspection of the three terms $\mathcal{E}_{1}$, $\mathcal{E}_{2}$ and $\mathcal{E}_3$.




%Suppose for the moment that there exist some quantities $\mathcal{E}_{1,1}>0$ and $0<\rho_{1}\leq1/2$ obeying








With this lemma in mind, everything boils down to (i) establishing the relation \eqref{eq:E1-UB-E11-rho1} and (ii) deriving a tight bound on $\mathcal{E}_{1,1}$, which form the main content of the rest of the proof. In light of the triangle inequality
%
\begin{align*}
	\big\|\bm{M}\big(\bm{U}\bm{H}-\bm{U}^{\star}\big)\big\|_{2,\infty}
	\leq
	\big\|\bm{E} \big(\bm{U}\bm{H}-\bm{U}^{\star}\big)\big\|_{2,\infty}
	+ \big\|\bm{M}^{\star}\big(\bm{U}\bm{H}-\bm{U}^{\star}\big)\big\|_{2,\infty},
\end{align*}
%
we dedicate the next two steps to bounding $\|\bm{E} (\bm{U}\bm{H}-\bm{U}^{\star} ) \|_{2,\infty}$ and $\|\bm{M}^{\star} (\bm{U}\bm{H}-\bm{U}^{\star}) \|_{2,\infty}$ respectively.





\subsubsection{Step 2: bounding $\|\bm{E} (\bm{U}\bm{H}-\bm{U}^{\star}) \|_{2,\infty}$ via leave-one-out analysis}

To obtain tight row-wise control of $\bm{E}\big(\bm{U}\bm{H}-\bm{U}^{\star}\big)$, one needs to carefully decouple the statistical dependency between $\bm{E}$ and $\bm{U}$,  which is where the leave-one-out idea comes into play.



\paragraph{Step 2.1: a convenient decomposition.}

We start by invoking the triangle inequality to decompose the target quantity as follows
%
\begin{align}
 & \big\|\bm{E}\big(\bm{U}\bm{H}-\bm{U}^{\star}\big)\big\|_{2,\infty}=\max_{l}\big\|\bm{E}_{l,\cdot}\big(\bm{U}\bm{H}-\bm{U}^{\star}\big)\big\|_{2}\nonumber \\
 & \leq\max_{l}\Big\{\big\|\bm{E}_{l,\cdot}\big(\bm{U}^{(l)}\bm{H}^{(l)}-\bm{U}^{\star}\big)\big\|_{2}+\big\|\bm{E}_{l,\cdot}\big(\bm{U}\bm{H}-\bm{U}^{(l)}\bm{H}^{(l)}\big)\big\|_{2}\Big\}\nonumber \\
 & \leq\max_{l}\Big\{\big\|\bm{E}_{l,\cdot}\big(\bm{U}^{(l)}\bm{H}^{(l)}-\bm{U}^{\star}\big)\big\|_{2}+\|\bm{E}\| \big\|\bm{U}\bm{H}-\bm{U}^{(l)}\bm{H}^{(l)}\big\|_{\mathrm{F}}\Big\}.
	\label{eq:decomposition-E-UH-Ustar}
\end{align}
%
In words, when controlling the $l$-th row of $\bm{E}\big(\bm{U}\bm{H}-\bm{U}^{\star}\big)$, we attempt to employ $\bm{U}^{(l)}\bm{H}^{(l)}$ as a surrogate of $\bm{U}\bm{H}$.  The benefits to be harvested from this decomposition are:
%
\begin{itemize}
	\item The statistical independence between $\bm{E}_{l,\cdot}$ and $\bm{U}^{(l)}\bm{H}^{(l)}$ allows for convenient upper bounds on $\big\|\bm{E}_{l,\cdot}\big(\bm{U}^{(l)}\bm{H}^{(l)}-\bm{U}^{\star}\big)\big\|_{2}$;
	
	\item  $\bm{U}\bm{H}$ and  $\bm{U}^{(l)}\bm{H}^{(l)}$ are expected to be exceedingly close, so that the discrepancy incurred by replacing $\bm{U}\bm{H}$ with $\bm{U}^{(l)}\bm{H}^{(l)}$ is negligible.
\end{itemize}
%
In what follows, we flesh out the proof details.

%With regards to the second term in \eqref{eq:MUh-Ustar-UB12},





%Our starting point is the following decomposition
%%
%\begin{align}
%	& \big\|\bm{M}\big(\bm{U}\bm{H}-\bm{U}^{\star}\big) \big\|_{2,\infty}  \notag\\
%  %\leq\big\|\bm{M}\big(\bm{U}\bm{H}-\bm{U}^{\star}\big)\big\|_{2,\infty}\big\|\big(\bm{\Lambda}^{\star}\big)^{-1}\big\|\\
% %& \qquad
%	%=\frac{\big\|\bm{M}\big(\bm{U}\bm{H}-\bm{U}^{\star}\big)\big\|_{2,\infty}}{\lambda_{r}^{\star}}
%	& \quad \leq   \big\|\bm{M}^{\star}\big(\bm{U}\bm{H}-\bm{U}^{\star}\big)\big\|_{2,\infty} + \big\|\bm{E}\big(\bm{U}\bm{H}-\bm{U}^{\star}\big)\big\|_{2,\infty}.
%	\label{eq:MUh-Ustar-UB12}
%\end{align}
%%


%In what follows, we shall take care of each of these two terms separately.
%




\paragraph{Step 2.2: the proximity of $\bm{U}\bm{H}$ and $\bm{U}^{(l)}\bm{H}^{(l)}$.}
Given that
%
\begin{align}
	\big\|\bm{U}\bm{H}-\bm{U}^{(l)}\bm{H}^{(l)}\big\|_{\mathrm{F}}
	& =\big\|\bm{U}\bm{U}^{\top}\bm{U}^{\star}-\bm{U}^{(l)}\bm{U}^{(l)\top}\bm{U}^{\star}\big\|_{\mathrm{F}} \notag\\
 	& \leq\big\|\bm{U}\bm{U}^{\top}-\bm{U}^{(l)}\bm{U}^{(l)\top}\big\|_{\mathrm{F}} \big\|\bm{U}^{\star}\big\| \notag\\
 	& =\big\|\bm{U}\bm{U}^{\top}-\bm{U}^{(l)}\bm{U}^{(l)\top}\big\|_{\mathrm{F}},
	\label{eq:UU-UlUl-UB01}
\end{align}
%
it boils down to bounding $\|\bm{U}\bm{U}^{\top}-\bm{U}^{(l)}\bm{U}^{(l)\top}\|_{\mathrm{F}}$.
Under simple conditions on the eigengap and the perturbation size (see \eqref{eq:lambda-rl-lambda-r1l-M-Ml-LB-step}  in Remark~\ref{remark:perturbation-size-eigen-gap}),  the Davis-Kahan theorem (cf.~Corollary~\ref{cor:davis-kahan-conclusion-corollary}) yields
%
\begin{align}
	\big\|\bm{U}\bm{U}^{\top}-\bm{U}^{(l)}\bm{U}^{(l)\top}\big\|_{\mathrm{F}}
	& \leq\frac{2\big\|\big(\bm{M}-\bm{M}^{(l)}\big)\bm{U}^{(l)}\big\|_{\mathrm{F}} }{ \big| \lambda_{r}^{(l)} \big| - \big|\lambda_{r+1}^{(l)} \big|} \nonumber\\
	& \leq\frac{4\big\|\big(\bm{M}-\bm{M}^{(l)}\big)\bm{U}^{(l)}\big\|_{\mathrm{F}} }{\lambda_{r}^{\star}}.
	\label{eq:UU-UlUl-UB1}
\end{align}
%



It remains to develop an upper bound on $\| (\bm{M}-\bm{M}^{(l)} )\bm{U}^{(l)} \|$.
The way we construct $\bm{M}^{(l)}$ (see Section~\ref{sec:construction-LOO-general}) allows us to express
%
\[
	\big(\bm{M}-\bm{M}^{(l)}\big)\bm{U}^{(l)}=\bm{e}_{l}\bm{E}_{l,\cdot}\bm{U}^{(l)}+  \big(\bm{E}_{\cdot,l} - E_{l,l} \bm{e}_l \big)  \bm{e}_{l}^{\top}\bm{U}^{(l)}.
\]
%
This together with the triangle inequality and the fact \eqref{eq:A-2inf-H-Hl-general} gives
%
\begin{align*}
 & \big\|\big(\bm{M}-\bm{M}^{(l)}\big)\bm{U}^{(l)}\big\|_{\mathrm{F}}\leq\big\|\bm{E}_{l,\cdot}\bm{U}^{(l)}\big\|_{2}+\big\|\bm{E}_{\cdot,l} - E_{l,l} \bm{e}_l \big\|_{2}\,\big\|\bm{U}^{(l)}\big\|_{2,\infty}\\
%	& \quad\leq\big\|\bm{E}_{l,\cdot}\bm{U}^{(l)}\big\|_{2}+\big\|\bm{E}\big\|\cdot\big\|\bm{U}^{(l)} \bm{H}^{(l)}  \big\|_{2,\infty} \big\| \big(\bm{H}^{(l)}\big)^{-1} \big\| \\
 & \quad\leq\big\|\bm{E}_{l,\cdot}\bm{U}^{(l)}\big\|_{2}+2\big\|\bm{E}\big\|\, \big\|\bm{U}^{(l)}\bm{H}^{(l)}\big\|_{2,\infty}\\
 & \quad\leq\big\|\bm{E}_{l,\cdot}\bm{U}^{(l)}\big\|_{2}+2\big\|\bm{E}\big\|\, \big\|\bm{U}\bm{H}\big\|_{2,\infty}+2\big\|\bm{E}\big\|\,\big\|\bm{U}\bm{H}-\bm{U}^{(l)}\bm{H}^{(l)}\big\|_{\mathrm{F}}.
\end{align*}
%
%where the last line arises from the triangle inequality.
%
Substitution into \eqref{eq:UU-UlUl-UB01} and \eqref{eq:UU-UlUl-UB1} gives
%
\begin{align*}
 & \big\|\bm{U}\bm{H}-\bm{U}^{(l)}\bm{H}^{(l)}\big\|_{\mathrm{F}} \leq \big\|\bm{U}\bm{U}^{\top}-\bm{U}^{(l)}\bm{U}^{(l)\top}\big\|_{\mathrm{F}} \\
 & \quad\leq\frac{4\big\|\bm{E}_{l,\cdot}\bm{U}^{(l)}\big\|_{2}+8\big\|\bm{E}\big\|\,\big\|\bm{U}\bm{H}\big\|_{2,\infty}+8\big\|\bm{E}\big\|\, \big\|\bm{U}\bm{H}-\bm{U}^{(l)}\bm{H}^{(l)}\big\|_{\mathrm{F}}}{\lambda_{r}^{\star}}.
\end{align*}
%
As long as $\|\bm{E}\|/\lambda_r^{\star}\leq 1/16$,
%(which is guaranteed if $16c_2\sigma\sqrt{n}\leq \lambda_r^{\star}$),
 one can further rearrange terms to obtain
%
\begin{align}
\big\|\bm{U}\bm{H}-\bm{U}^{(l)}\bm{H}^{(l)}\big\|_{\mathrm{F}} & \leq\frac{8\big\|\bm{E}_{l,\cdot}\bm{U}^{(l)}\big\|_{2}+16\big\|\bm{E}\big\|\, \big\|\bm{U}\bm{H}\big\|_{2,\infty}}{\lambda_{r}^{\star}}.
	\label{eq:UH-UlHl-diff-UB10}
\end{align}
%
In addition, the fact  \eqref{eq:A-2inf-H-Hl-general} combined with the triangle inequality yields
%
\[
	\tfrac{1}{2}\big\|\bm{E}_{l,\cdot}\bm{U}^{(l)}\big\|_{2}\leq \big\|\bm{E}_{l,\cdot}\bm{U}^{(l)}\bm{H}^{(l)}\big\|_{2}
	\leq \big\|\bm{E}_{l,\cdot}\big(\bm{U}^{(l)}\bm{H}^{(l)}-\bm{U}^{\star}\big)\big\|_{2}+ \big\|\bm{E}_{l,\cdot}\bm{U}^{\star}\big\|_{2},
\]
%
which taken collectively with \eqref{eq:UH-UlHl-diff-UB10} reveals that
%\yxc{check whether to use $UH$ or $U$}
%
\begin{align}
 & \big\|\bm{U}\bm{H}-\bm{U}^{(l)}\bm{H}^{(l)}\big\|_{\mathrm{F}}
	\leq\frac{16\big\|\bm{E}_{l,\cdot}\big(\bm{U}^{(l)}\bm{H}^{(l)}-\bm{U}^{\star}\big)\big\|_{2}}{\lambda_{r}^{\star}} \nonumber \\
	&  \quad + \frac{ 16 \Big\{ \big\|\bm{E}_{l,\cdot}\bm{U}^{\star}\big\|_{2}  + \big\|\bm{E}\big\| \, \big\|\bm{U}\bm{H}-\bm{U}^{\star}\big\|_{2,\infty} +  \big\|\bm{E}\big\|\, \big\|\bm{U}^{\star}\big\|_{2,\infty} \Big\} }{\lambda_{r}^{\star}}.
	\label{eq:UH-UlHl-diff-UB123}
\end{align}
%










\paragraph{Step 2.3: bounding $\|\bm{E}_{l,\cdot} (\bm{U}^{(l)}\bm{H}^{(l)}-\bm{U}^{\star} ) \|_{2}$.}

%Another term that appears in  \eqref{eq:decomposition-E-UH-Ustar}
%
%is $\|\bm{E}_{l,\cdot} (\bm{U}^{(l)}\bm{H}^{(l)}-\bm{U}^{\star} ) \|_{2}$, which we look into in this step.
Recognizing that $\bm{E}_{l,\cdot}$ is statistically independent of $\bm{U}^{(l)}$ (since $\bm{U}^{(l)}$ is computed without using $\bm{E}_{l,\cdot}$),
we invoke Lemma~\ref{lemma:general-noise-bound} (more precisely, we use the proof of this lemma) to demonstrate that
with probability exceeding $1-2n^{-6}$,
%
\begin{align}
 & \big\|\bm{E}_{l,\cdot}\big(\bm{U}^{(l)}\bm{H}^{(l)}-\bm{U}^{\star}\big)\big\|_{2} \notag\\
	%\leq4\sqrt{v\log n}+6L\log n\nonumber \\
 & \leq4\sigma\sqrt{\log n} \, \|\bm{U}^{(l)}\bm{H}^{(l)}-\bm{U}^{\star}\|_{\mathrm{F}}
	+ (6B \log n) \|\bm{U}^{(l)}\bm{H}^{(l)}-\bm{U}^{\star}\|_{2,\infty} \nonumber \\
	& \leq4\sigma \sqrt{\log n} \, \|\bm{U}\bm{H}-\bm{U}^{\star}\|_{\mathrm{F}} + (6B \log n) \|\bm{U}\bm{H}-\bm{U}^{\star}\|_{2,\infty} \nonumber \\
	& \qquad+ ( 10B\log n) \|\bm{U}\bm{H}-\bm{U}^{(l)}\bm{H}^{(l)}\|_{\mathrm{F}}
 \label{eq:El-UH-Ustar-UB-Bern1}
\end{align}
%
holds simultaneously for all $1\leq l\leq n$, where the last line results from the triangle inequality and the fact $4\sigma\sqrt{\log n}+6B\log n\leq 10B\log n$.


% using Assumption~\ref{eq:assumption-B-sigma}


\paragraph{Step 2.4: combining the above bounds.}

The careful reader would immediately remark that the inequalities \eqref{eq:UH-UlHl-diff-UB123} and \eqref{eq:El-UH-Ustar-UB-Bern1} are convoluted, both of which  involve the terms  $\|\bm{E}_{l,\cdot} (\bm{U}^{(l)}\bm{H}^{(l)}-\bm{U}^{\star} ) \|_{2}$ and $\|\bm{U}\bm{H}-\bm{U}^{(l)}\bm{H}^{(l)}\|_{\mathrm{F}}$. Fortunately, one can substitute \eqref{eq:El-UH-Ustar-UB-Bern1} into \eqref{eq:UH-UlHl-diff-UB123} to produce a cleaner bound. By doing so and exploiting the condition $320B\log n \leq \lambda_r^{\star}$, we  rearrange terms to reach
%
\begin{align*}
	& \big\|\bm{U}\bm{H}-\bm{U}^{(l)}\bm{H}^{(l)}\big\|_{\mathrm{F}}
	\leq \frac{32\big\|\bm{E}_{l,\cdot}\bm{U}^{\star}\big\|_{2}+32\big\|\bm{E}\big\|\, \big\|\bm{U}^{\star}\big\|_{2,\infty}}{\lambda_{r}^{\star}}\nonumber \\
	& +\frac{128\sigma\sqrt{\log n}\big\|\bm{U}\bm{H}-\bm{U}^{\star}\big\|_{\mathrm{F}}+(32c_2\sigma \sqrt{n} +192B\log n)\big\|\bm{U}\bm{H}-\bm{U}^{\star}\big\|_{2,\infty}}{\lambda_{r}^{\star}},
	%\label{eq:UH-UlHl-diff-UB456}
\end{align*}
%
where we have also used the upper bound on $\|\bm{E}\|$ derived in \eqref{eq:L2-perturbation-summary-general}.
Meanwhile, plugging the above inequality into \eqref{eq:El-UH-Ustar-UB-Bern1} yields
%
\begin{align}
	& \big\|\bm{E}_{l,\cdot}\big(\bm{U}^{(l)}\bm{H}^{(l)}-\bm{U}^{\star}\big)\big\|_{2}
	\leq \frac{320B\log n}{\lambda_{r}^{\star}}\big(\big\|\bm{E}_{l,\cdot}\bm{U}^{\star}\big\|_{2}+\big\|\bm{E}\big\|\,\big\|\bm{U}^{\star}\big\|_{2,\infty}\big) \nonumber\\
	&\quad + 5\sigma\sqrt{\log n}\,\|\bm{U}\bm{H}-\bm{U}^{\star}\|_{\mathrm{F}}+(7B\log n)\|\bm{U}\bm{H}-\bm{U}^{\star}\|_{2,\infty}  ,
	\label{eq:El-UH-Ustar-UB-Bern123}
\end{align}
%
provided that $\max \{ \sigma \sqrt{n}, B\log n\} \leq c_3\lambda_r^{\star}$ for some  small constant $c_3$.



Substituting the preceding two bounds into \eqref{eq:decomposition-E-UH-Ustar} and combining terms reveal the existence of some constant $c_4>0$ such that
%
\begin{align}
	&  \big\|\bm{E}\big(\bm{U}\bm{H}-\bm{U}^{\star}\big)\big\|_{2,\infty} \notag\\
	& \qquad \leq \alpha_0 + \alpha_1 + c_4\big( \sigma  \sqrt{n} + B\log n \big)\|\bm{U}\bm{H}-\bm{U}^{\star}\|_{2,\infty}
	\label{eq:El-UH-Ustar-UB-Bern123}
\end{align}
%
provided that $\max \{ \sigma \sqrt{n}, B\log n\} \leq c_3\lambda_r^{\star}$ for some constant $c_3>0$ small enough, where
\begin{align}
\begin{array}{ll}
 & \alpha_{0}\coloneqq\frac{32c_{2}\sigma\sqrt{n}+320B\log n}{\lambda_{r}^{\star}}\big(\big\|\bm{E}\bm{U}^{\star}\big\|_{2,\infty}+\big\|\bm{E}\big\|\,\big\|\bm{U}^{\star}\big\|_{2,\infty}\big) , \\
 & \alpha_{1}\coloneqq6\sigma\sqrt{\log n}\,\|\bm{U}\bm{H}-\bm{U}^{\star}\|_{\mathrm{F}} .
\end{array}	
\label{eq:defn-alpha-0-1-general}
\end{align}
%


Before continuing,  note that we are already well-equipped to bound the above two quantities. First, $\alpha_0$ can be bounded by
%
\begin{align}
\alpha_{0} & \leq \frac{32c_{2}\sigma\sqrt{n}+320B\log n}{\lambda_{r}^{\star}}\nonumber \\
 & \qquad\cdot\Big(4\sigma\sqrt{r\log n}+6B\sqrt{\frac{\mu r}{n}}\log n+c_{2}\sigma\sqrt{\mu r}\Big),
\label{eq:defn-alpha-0-general}
\end{align}
%
where we have used the bounds concerning $\bm{E}$ from Lemma~\ref{lemma:general-noise-bound} as well as the definition \eqref{eq:defn-Ustar-incoherence}.
Regarding $\alpha_1$, it is seen from (\ref{eq:dist-under-H-general}) that
%
\begin{align}
\alpha_{1}
  \leq\frac{12c_{2}\sigma^{2}\sqrt{rn\log n}}{\lambda_{r}^{\star}}.
	\label{eq:alpha1-bound-1234}
\end{align}


%Substituting the inequalities \eqref{eq:UH-UlHl-diff-UB123} and \eqref{eq:El-UH-Ustar-UB-Bern1} into \eqref{eq:decomposition-E-UH-Ustar} with a little algebra yields
%%
%\begin{align}
% & \big\|\bm{E}(\bm{U}\bm{H}-\bm{U}^{\star})\big\|_{2,\infty}\nonumber \\
% & \quad\leq4\sigma\sqrt{\log n}\|\bm{U}\bm{H}-\bm{U}^{\star}\|_{\mathrm{F}}+(6B\log n)\|\bm{U}\bm{H}-\bm{U}^{\star}\|_{2,\infty}\nonumber \\
% & \qquad+\frac{4\big(\|\bm{E}\|+10B\log n\big)\max_{l}\Big\{\big\|\bm{E}_{l,\cdot}\bm{U}^{(l)}\big\|_{2}+\big\|\bm{E}\big\|\cdot\big\|\bm{U}^{(l)}\big\|_{2,\infty}\Big\}}{\lambda_{r}^{\star}}
%	\label{eq:E-UH-Ustar-UB-234}
%\end{align}
%%


\subsubsection{Step 3: bounding $\|\bm{M}^{\star} (\bm{U}\bm{H}-\bm{U}^{\star} ) \|_{2,\infty}$}


We make the key observation that
%
\begin{align}
\big\|\bm{M}^{\star}\big(\bm{U}\bm{H}-\bm{U}^{\star}\big)\big\|_{2,\infty} & =\big\|\bm{U}^{\star}\bm{\Lambda}^{\star}\bm{U}^{\star\top}\big(\bm{U}\bm{H}-\bm{U}^{\star}\big)\big\|_{2,\infty} \notag\\
 & \leq\big\|\bm{U}^{\star}\big\|_{2,\infty} \|\bm{\Lambda}^{\star}\|\, \big\|\bm{U}^{\star\top}\big(\bm{U}\bm{H}-\bm{U}^{\star}\big)\big\| \notag\\
	& =  \sqrt{\frac{\mu r}{n}} \|\bm{\Lambda}^{\star}\|\, \big\|\bm{U}\bm{U}^{\top} - \bm{U}^{\star} \bm{U}^{\star\top} \big\|^{2}.
	\label{eq:Mstar-UH-Ustar-2-UB-123}
\end{align}
%
Here, the last identity holds true due to the following observation
%
\begin{align*}
\big\|\bm{U}^{\star\top}\big(\bm{U}\bm{H}-\bm{U}^{\star}\big)\big\| & =\big\|\bm{U}^{\star\top}\bm{U}\bm{U}^{\top}\bm{U}^{\star}-\bm{U}^{\star\top}\bm{U}^{\star}\big\|=\big\|\bm{U}^{\star\top}\bm{U}\bm{U}^{\top}\bm{U}^{\star}-\bm{I}\big\|\\
 & =\big\|\bm{Y}(\cos^{2}\bm{\Theta})\bm{Y}^{\top}-\bm{I}\big\|=\big\| \cos^{2}  \bm{\Theta} - \bm{I} \big\|\\
	& =\big\|\sin^{2}\bm{\Theta}\big\|=\big\|\sin\bm{\Theta}\big\|^2,
\end{align*}
%
where we denote by $\bm{X}(\cos \bm{\Theta}) \bm{Y}^{\top}$ the SVD of $\bm{U}^{\top}\bm{U}^{\star}$, with $\bm{X}$ and $\bm{Y}$ being orthonormal matrices and $\bm{\Theta}$ the diagonal matrix consisting of the principal angles between $\bm{U}$ and $\bm{U}^{\star}$ (see the definition in~\eqref{defn:principal-angles}). The above bounds combined with \eqref{eq:distp-U-Ustar-general} indicate that
%
\begin{align}
\big\|\bm{M}^{\star}\big(\bm{U}\bm{H}-\bm{U}^{\star}\big)\big\|_{2,\infty} & \leq
	%\big|\lambda_{1}^{\star}\big|\sqrt{\frac{\mu r}{n}}\,\frac{4c_{2}^{2}\sigma^{2}n}{\big(\lambda_{r}^{\star}\big)^{2}}
	\big|\lambda_{1}^{\star}\big| \sqrt{\frac{\mu r}{n}} \, \big\|\sin\bm{\Theta}\big\|^2
	% = \big|\lambda_{1}^{\star}\big| \sqrt{\frac{\mu r}{n}} \, \big\| \sin \bm{\Theta} \big\|^{2}
	\notag\\
	& \leq \frac{4c_{2}^{2}\kappa\sigma^{2}\sqrt{\mu rn}}{\lambda_{r}^{\star}}
 \eqqcolon \alpha_2.
	\label{eq:Mstar-UH-Ustar-2inf}
\end{align}
%



\subsubsection{Step 4: putting all pieces together}

Combining the bounds \eqref{eq:El-UH-Ustar-UB-Bern123} and \eqref{eq:Mstar-UH-Ustar-2inf} in Steps 2-3 and using the definition of $\mathcal{E}_1$ (see Lemma~\ref{lem:UH-MUstar-diff-decompose}) give
%
\begin{align}
\mathcal{E}_{1} & \leq\frac{2\big\|\bm{E}\big(\bm{U}\bm{H}-\bm{U}^{\star}\big)\big\|_{2,\infty}+2\big\|\bm{M}^{\star}\big(\bm{U}\bm{H}-\bm{U}^{\star}\big)\big\|_{2,\infty}}{\lambda_{r}^{\star}} \nonumber \\
	& \leq\mathcal{E}_{1,1}+\rho_{1}\|\bm{U}\bm{H}-\bm{U}^{\star}\|_{2,\infty},
	\label{eq:E1-UB-general-1}
\end{align}
%
where
%
\begin{equation}
	\mathcal{E}_{1,1}\coloneqq\frac{2\big(\alpha_{0}+\alpha_{1}+\alpha_{2}\big)}{\lambda_{r}^{\star}}
	\quad\text{and}\quad
	\rho_{1}\coloneqq\frac{2c_{4}(\sigma\sqrt{n}+B\log n)}{\lambda_{r}^{\star}}.
	\label{eq:defn-E1-rho1}
\end{equation}
%
This matches precisely the relation hypothesized in \eqref{eq:E1-UB-E11-rho1}. In particular, one has $0<\rho\leq 1/2$
as long as $4c_4(\sigma \sqrt{n}+ B\log n) \leq \lambda_r^{\star}$, which holds whenever $\sigma \sqrt{n \log n}\leq c_{\sigma} \lambda_r^{\star}$ for a sufficiently small $c_{\sigma}>0$ in view of our assumption on $B$ (cf. \eqref{eq:assumption-B-sigma}).


Recall that $\mathcal{E}_{1,1}$, $\mathcal{E}_{2}$, $\mathcal{E}_{3}$, $\alpha_0$, $\alpha_1$, $\alpha_2$ and $\rho_1$ have been controlled in \eqref{eq:defn-E1-rho1}, \eqref{eq:MUstar-E-norm-lambda2-bound}, \eqref{eq:E3-bound-general}, \eqref{eq:defn-alpha-0-general}, \eqref{eq:alpha1-bound-1234}, \eqref{eq:Mstar-UH-Ustar-2inf} and \eqref{eq:defn-E1-rho1}, respectively.
In addition, recall our assumption \eqref{eq:assumption-B-sigma} and suppose that $\sigma{\sqrt{n}}\leq c_{\sigma}\lambda_{r}^{\star}$ for some sufficiently small constant $c_{\sigma}>0$.
With these bounds and assumptions in mind, invoking Lemma~\ref{lem:UH-sequence-bounds-E123} and combining terms immediately conclude the proof.



\begin{remark} We shall also make note of an immediate consequence of the above argument as follows
%
\begin{align}
%\big\|\bm{U}\bm{H}-\bm{M}\bm{U}^{\star}\big(\bm{\Lambda}^{\star}\big)^{-1}\big\|_{2,\infty} & \lesssim\frac{\sigma\kappa\sqrt{\mu r}}{\lambda_{r}^{\star}}+\frac{c_{\mathsf{b}}\sigma^{2}\sqrt{rn}\log n}{\big(\lambda_{r}^{\star}\big)^{2}},
	\big\|\bm{U}\bm{H}-\bm{U}^{\star}\big\|_{2,\infty}&\lesssim\frac{\sigma\kappa\sqrt{\mu r}+\sigma\sqrt{r\log n}}{\lambda_{r}^{\star}},
	\label{eq:UH-MUstar-Lambda-star-2inf}
\end{align}
%
which will prove useful for deriving other important results.
\end{remark}




%and recall the bound \eqref{eq:UH-Ustar-diff-2inf-general}.
%
%
%
%%By rearranging terms, we are left with
%%%
%%\begin{align}
%%\|\bm{U}\bm{H}-\bm{U}^{\star}\|_{2,\infty}
%%	& \leq
%%	\frac{4\big(\alpha_{0}+\alpha_{1} + \alpha_2 \big)}{\lambda_{r}^{\star}}+2\mathcal{E}_{2}+2\mathcal{E}_{3}.
%%	\label{eq:UH-Ustar-general-UB2}
%%\end{align}
%%%
%%Substituting this relation into \eqref{eq:E1-UB-general-1}, we arrive at
%%%
%%\begin{align*}
%%\mathcal{E}_{1} & \leq\frac{2\big(\alpha_{0}+\alpha_{1}+\alpha_{2})+2c_{4}(\sigma\sqrt{n}+B\log n)\Big(\frac{4(\alpha_{0}+\alpha_{1}+\alpha_{2})}{\lambda_{r}^{\star}}+2\mathcal{E}_{2}+3\mathcal{E}_{3}\Big)}{\lambda_{r}^{\star}}\\
%% & \leq\frac{4\big(\alpha_{0}+\alpha_{1}+\alpha_{2})+c_{4}(\sigma\sqrt{n}+B\log n)\Big(4\mathcal{E}_{2}+6\mathcal{E}_{3}\Big)}{\lambda_{r}^{\star}},
%%\end{align*}
%%%
%%provided that $4c_4(\sigma\sqrt{n}+B\log n) \leq \lambda_r^{\star}$.
%
%
%
%Returning to  \eqref{eq:UH-MUstarLambdastar-diff-2inf-general}, we can combine the above bound on $\mathcal{E}_1$ to deduce that
%%
%\begin{align*}
% & \big\|\bm{U}\bm{H}-\bm{M}\bm{U}^{\star}\big(\bm{\Lambda}^{\star}\big)^{-1}\big\|_{2,\infty}\leq\mathcal{E}_{1}+\mathcal{E}_{2}\\
% & \qquad\leq\frac{4\big(\alpha_{0}+\alpha_{1}+\alpha_{2})+6c_{4}(\sigma\sqrt{n}+B\log n)\mathcal{E}_{3}}{\lambda_{r}^{\star}}+2\mathcal{E}_{2},
%\end{align*}
%%
%provided that $4c_4(\sigma\sqrt{n}+B\log n) \leq \lambda_r^{\star}$
%
%
%
%Furthermore,
%%
%\begin{align*}
%\big\|\bm{U}\bm{H}-\bm{U}\mathsf{sgn}(\bm{H})\big\|_{2,\infty} & \leq\big\|\bm{U}\big\|_{2,\infty}\big\|\bm{H}-\mathsf{sgn}(\bm{H})\big\|\leq\frac{4c_{2}^{2}\sigma^{2}n}{(\lambda_{r}^{\star})^{2}}\big\|\bm{U}\bm{H}\big\|_{2,\infty}\\
% & \leq\frac{4c_{2}^{2}\sigma^{2}n}{(\lambda_{r}^{\star})^{2}}\left\{ \big\|\bm{U}\bm{H}-\bm{U}^{\star}\big\|_{2,\infty}+\big\|\bm{U}^{\star}\big\|_{2,\infty}\right\} ,
%\end{align*}
%%
%where the second inequality follows from \eqref{eq:H-property-summary-general} and \eqref{eq:A-2inf-H-Hl-general}. Consequently,
%%
%\begin{align*}
% & \big\|\bm{U}\mathsf{sgn}(\bm{H})-\bm{M}\bm{U}^{\star}\big(\bm{\Lambda}^{\star}\big)^{-1}\big\|_{2,\infty}\\
% & \quad\leq\big\|\bm{U}\bm{H}-\bm{U}\mathsf{sgn}(\bm{H})\big\|_{2,\infty}+\big\|\bm{U}\bm{H}-\bm{M}\bm{U}^{\star}\big(\bm{\Lambda}^{\star}\big)^{-1}\big\|_{2,\infty}\\
% & \quad\leq\frac{4c_{2}^{2}\sigma^{2}n}{(\lambda_{r}^{\star})^{2}}\left\{ \big\|\bm{U}\bm{H}-\bm{U}^{\star}\big\|_{2,\infty}+\sqrt{\frac{\mu r}{n}}\right\} +\big\|\bm{U}\bm{H}-\bm{M}\bm{U}^{\star}\big(\bm{\Lambda}^{\star}\big)^{-1}\big\|_{2,\infty}\\
% & \quad\leq\frac{6\big(\alpha_{0}+\alpha_{1}+\alpha_{2})+8c_{4}(\sigma\sqrt{n}+B\log n)\mathcal{E}_{3}}{\lambda_{r}^{\star}}+3\mathcal{E}_{2}+\frac{4c_{2}^{2}\sigma^{2}\sqrt{\mu nr}}{(\lambda_{r}^{\star})^{2}} .
%\end{align*}
%%
%\yxc{maybe sent earlier}
%


%combine the triangle inequality and the relation \eqref{eq:UH-Ustar-general-UB2} to yield
%%
%\begin{align*}
% & \big\|\bm{U}\bm{H}-\bm{M}\bm{U}^{\star}\big(\bm{\Lambda}^{\star}\big)^{-1}\big\|_{2,\infty}\leq\big\|\bm{U}\bm{H}-\bm{M}^{\star}\bm{U}^{\star}\big(\bm{\Lambda}^{\star}\big)^{-1}+\bm{E}\bm{U}^{\star}\big(\bm{\Lambda}^{\star}\big)^{-1}\big\|_{2,\infty}\\
% & \quad\quad\leq\big\|\bm{U}\bm{H}-\bm{M}^{\star}\bm{U}^{\star}\big(\bm{\Lambda}^{\star}\big)^{-1}\big\|_{2,\infty}+\big\|\bm{E}\bm{U}^{\star}\big(\bm{\Lambda}^{\star}\big)^{-1}\big\|_{2,\infty}\\
% & \quad\quad\leq\big\|\bm{U}\bm{H}-\bm{U}^{\star}\big\|_{2,\infty}+\mathcal{E}_{3}
%	\leq\frac{4\big(\alpha_{0}+\alpha_{1}+\alpha_{2}\big)}{\lambda_{r}^{\star}}+2\mathcal{E}_{2}+3\mathcal{E}_{3}.
%\end{align*}
%%
%


\subsection{Proof of auxiliary lemmas}
\label{sec:proof-auxiliary-lemmas-general}




%\paragraph{Proof of Lemma~\ref{lem:preliminary-result-general-iid}.}
%
%Under Assumption~\ref{assumption-noise-general},  Theorem~\ref{thm:tighter-spectral-normal-ramon} (in particular \eqref{eq:X-spectral-norm-iid-special-ramon})  reveals the existence of some constant $c_2>0$ such that
%%
%\begin{align}
%	\big\|\bm{E}^{(l)} \big\|\leq \|\bm{E}\|\leq c_2 \sigma \sqrt{n}
%	%\qquad\text{and} \qquad
%	% c_2 \sigma \sqrt{n}
%	\qquad 1\leq l \leq n,
%	\label{eq:noise-E-size-general}
%\end{align}
%%
%holds with probability exceeding $1- O(n^{-7})$. Additionally, repeating the argument in the proof of Corollary~\ref{cor:davis-kahan-conclusion-corollary} reveals that
%%
%\begin{subequations}
%\begin{align}
%	\max_{1\leq j\leq r} | \lambda_j  |  &\geq \lambda_r^{\star} - \|\bm{E}\|,  \quad &  \max_{1\leq j\leq r} | \lambda^{(l)}_j  |  &\geq \lambda_r^{\star} - \|\bm{E}^{(l)}\|, &&\\
%	\max_{j: j> r} | \lambda_j   |  &\leq \|\bm{E}\| ,  & \max_{j: j> r}| \lambda_j^{(l)} |  &\leq \|\bm{E}^{(l)}\|,
%\end{align}
%\end{subequations}
%%
%which taken together with \eqref{eq:noise-E-size-general} validates  \eqref{eq:lambda-size-large-general}-\eqref{eq:lambda-size-small-general}.
%
%
%Regarding \eqref{eq:dist-U-Ustar-general} and \eqref{eq:distp-U-Ustar-general}, apply Corollary~\ref{cor:davis-kahan-conclusion-corollary} to reach
%%
%\begin{align}
%	\mathsf{dist}(\bm{U},\bm{U}^{\star}) \leq  \sqrt{2} \|\sin \bm{\Theta} \|
%	\leq\frac{2\|\bm{E}\|}{\lambda_{r}^{\star}}\leq\frac{2c_{2}\sigma\sqrt{n}}{\lambda_{r}^{\star}}	
%	\label{eq:dist-UUstar-sin-Theta-general-UB123}
%\end{align}
%%
%as long as $\|\bm{E}\|\leq c_2 \sigma \sqrt{n}\leq (1-1/\sqrt{2})\lambda_r^{\star}$, where we have used $\lambda_{r+1}^{\star}=0$. The bound on $\mathsf{dist}\big( \bm{U}^{(l)} , \bm{U}^{\star} \big)$ follows from the same argument.
%%
%
%
%When it comes to $\bm{E}\bm{A}$, we proceed by viewing
%its $l$-th row $\bm{E}_{l,\cdot}\bm{A}$ as a sum of independent random
%vectors as follows
%\[
%\bm{E}_{l,\cdot}\bm{A}=\sum\nolimits_{j}E_{l,j}\bm{A}_{j,\cdot}\eqqcolon\sum\nolimits_{j}\bm{z}_{j},
%\]
%which can be controlled by the matrix Bernstein inequality. Specifically,
%it is seen from Assumption~\ref{assumption-noise-general} that
%%
%\begin{align*}
%v & \coloneqq\sum\nolimits_{j}\mathbb{E}\big[\|\bm{z}_{j}\|_{2}^{2}\big]
%	% = \sum_{j}\mathbb{E}\big[E_{l,j}^{2}\big] \mathbb{E}\big[\|\bm{A}_{j,\cdot}\|_{2}^{2}\big]
%	\leq\sigma^{2}\sum\nolimits_{j}\big\|\bm{A}_{j,\cdot}\big\|_{2}^{2}=\sigma^{2}\|\bm{A}\|_{\mathrm{F}}^{2};\\
%L & \coloneqq\max_{j}\|\bm{z}_{j}\|_{2}=\max_{j}|E_{l,j}|\,\big\|\bm{Z}_{j,\cdot}\big\|_{2}\leq B\|\bm{Z}\|_{2,\infty}.
%\end{align*}
%%
%Invoke the matrix Bernstein inequality (cf.~Corollary~\ref{thm:matrix-Bernstein-friendly})
%and take the union bound to demonstrate that: with probability exceeding
%$1-2n^{-6}$,
%%
%\begin{align*}
%	\big\|\bm{E}_{l,\cdot}\bm{A}\big\|_{2} & \leq4\sqrt{v\log n}+6L\log n
%	\leq  4\sigma\sqrt{\log n}\,\|\bm{A}\|_{\mathrm{F}}+(6B\log n)\|\bm{A}\|_{2,\infty}
%\end{align*}
%%
%holds simultaneously for all $1\leq l\leq n$, thus concluding the proof.
%
%
%
%
%\paragraph{Proof of Lemma~\ref{lem:H-property-summary-general}.}
%
%We shall only prove the result for $\bm{H}$; the proof for  $\bm{H}^{(l)}$ follows from identical arguments and is hence omitted.
%
%
%From our discussion in Section~\ref{sec:introduction-distance-principal-angles},
%one can express the SVD of $\bm{H}=\bm{U}^{\top}\bm{U}^{\star}$ as
%$\bm{H}=\bm{X}(\cos\bm{\Theta})\bm{Y}^{\top}$, where the columns
%of $\bm{X}$ (resp.~$\bm{Y}$) are the left (resp.~right) singular
%vectors of $\bm{H}$, and $\bm{\Theta}$ is a diagonal matrix composed
%of the principal angles between $\bm{U}$ and $\bm{U}^{\star}$. In
%light of this  and the definition (\ref{eq:defn-sgn-Z}),
%we can establish \eqref{eq:H-sgnH-dist-general} as follows
%%
%\begin{align}
%\big\|\bm{H}-\mathsf{sgn}(\bm{H})\big\| & =\big\|\bm{X}\big(\cos\bm{\Theta}-\bm{I}\big)\bm{Y}^{\top}\big\|=\big\|\bm{I}-\cos\bm{\Theta}\big\| \notag\\
% & =2\big\|\sin^{2}(\bm{\Theta}/2)\big\|\leq\|\sin\bm{\Theta}\|^{2} \notag\\
% & \leq \frac{2c_{2}^{2}\sigma^{2}n}{(\lambda_{r}^{\star})^{2}},
%	\label{eq:H-sgnH-diff-UB-123}
%\end{align}
%%
%where the middle line holds since $1-\cos\theta=2\sin^{2}(\theta/2)$
%and $\sin(\theta/2)\leq\frac{\sqrt{2}}{2}\sin\theta$ for any $0\leq\theta\leq\pi/2$, and the last line follows from \eqref{eq:distp-U-Ustar-general}.
%
%Coming back to the claim \eqref{eq:H-inv-norm-bound-general}, it suffices to justify that $\sigma_{\min}(\bm{H})\geq 1/2$.
%Recognizing that $\mathsf{sgn}(\bm{H})=\bm{X}\bm{Y}^{\top}$,
%we see that all singular values of $\mathsf{sgn}(\bm{H})$  equal
% 1. Thus, Weyl's inequality together with (\ref{eq:H-sgnH-diff-UB-123})
%gives
%\[
%\sigma_{\min}(\bm{H})\geq\sigma_{\min}\big(\mathsf{sgn}(\bm{H})\big)-\|\bm{H}-\mathsf{sgn}(\bm{H})\|\geq1-\frac{2c_{2}^{2}\sigma^{2}n}{(\lambda_{r}^{\star})^{2}}\geq\frac{1}{2} ,
%\]
%with the proviso that $2c_{2}\sigma\sqrt{n}\leq\lambda_{r}^{\star}$.
%% This concludes the proof.
%
%
%Additionally, regarding $\bm{U}\bm{H}-\bm{U}^{\star}$ we can derive
%%
%\begin{align}
% & \|\bm{U}\bm{H}-\bm{U}^{\star}\|_{\mathrm{F}}
%   = \|\bm{U}\bm{U}^{\top}\bm{U}^{\star}-\bm{U}^{\star}\bm{U}^{\star\top}\bm{U}^{\star}\|_{\mathrm{F}} \nonumber\\
% & \qquad \quad \leq \|\bm{U}\bm{U}^{\top}-\bm{U}^{\star}\bm{U}^{\star\top}\|\,\|\bm{U}^{\star}\|_{\mathrm{F}}
%   \leq\frac{2c_{2}\sigma \sqrt{rn}}{\lambda_{r}^{\star}},
%\end{align}
%%
%where the last line arises from \eqref{eq:distp-U-Ustar-general}, Lemma~\ref{prop:unitary-norm-property} and $\|\bm{U}^{\star}\|_{\mathrm{F}}=\sqrt{r}$.
%
%
%
\paragraph{Proof of Lemma~\ref{lemma:general-noise-bound}.}

Under Assumption~\ref{assumption-noise-general}, Theorem~\ref{thm:tighter-spectral-normal-ramon}
(in particular \eqref{eq:X-spectral-norm-iid-special-ramon}) reveals
the existence of some constant $c_{2}>0$ such that
\[
\max_{l}\big\|\bm{E}^{(l)}\big\|\leq\|\bm{E}\|\leq c_{2}\sigma\sqrt{n},
\]
holds with probability exceeding $1-O(n^{-7})$.

When it comes to $\bm{E}\bm{A}$, we proceed by viewing its $l$-th
row $\bm{E}_{l,\cdot}\bm{A}$ as a sum of independent random vectors
as follows
\[
\bm{E}_{l,\cdot}\bm{A}=\sum\nolimits _{j}E_{l,j}\bm{A}_{j,\cdot}\eqqcolon\sum\nolimits _{j}\bm{z}_{j},
\]
which can be controlled by the matrix Bernstein inequality. Specifically,
it is seen from Assumption~\ref{assumption-noise-general} that
\begin{align*}
v & \coloneqq\sum\nolimits _{j}\mathbb{E}\big[\|\bm{z}_{j}\|_{2}^{2}\big]%=\sum_{j}\mathbb{E}\big[E_{l,j}^{2}\big]\mathbb{E}\big[\|\bm{A}_{j,\cdot}\|_{2}^{2}\big]
\leq\sigma^{2}\sum\nolimits _{j}\big\|\bm{A}_{j,\cdot}\big\|_{2}^{2}=\sigma^{2}\|\bm{A}\|_{\mathrm{F}}^{2};\\
L & \coloneqq\max_{j}\|\bm{z}_{j}\|_{2}=\max_{j}|E_{l,j}|\,\big\|\bm{Z}_{j,\cdot}\big\|_{2}\leq B\|\bm{Z}\|_{2,\infty}.
\end{align*}
Invoke the matrix Bernstein inequality (cf.~Corollary~\ref{thm:matrix-Bernstein-friendly})
and take the union bound to demonstrate that: with probability exceeding
$1-2n^{-6}$,
\begin{align*}
\big\|\bm{E}_{l,\cdot}\bm{A}\big\|_{2} & \leq4\sqrt{v\log n}+6L\log n\leq4\sigma\sqrt{\log n}\,\|\bm{A}\|_{\mathrm{F}}+(6B\log n)\|\bm{A}\|_{2,\infty}
\end{align*}
holds simultaneously for all $1\leq l\leq n$, thus concluding the
proof.

\paragraph{Proof of Lemma~\ref{lem:preliminary-result-general-iid}.}

Lemma~\ref{lemma:general-noise-bound} tells us that
\begin{equation}
\max_{l}\big\|\bm{E}^{(l)}\big\|\leq\|\bm{E}\|\leq c_{2}\sigma\sqrt{n}\label{eq:noise-E-size-general}
\end{equation}
holds with probability exceeding $1-O(n^{-7})$. Repeating the argument
in the proof of Corollary~\ref{cor:davis-kahan-conclusion-corollary}
reveals that \begin{subequations}
\begin{align}
\max_{1\leq j\leq r}|\lambda_{j}| & \geq\lambda_{r}^{\star}-\|\bm{E}\|,\quad & \max_{1\leq j\leq r}|\lambda_{j}^{(l)}| & \geq\lambda_{r}^{\star}-\|\bm{E}^{(l)}\|,\\
\max_{j:j>r}|\lambda_{j}| & \leq\|\bm{E}\|, & \max_{j:j>r}|\lambda_{j}^{(l)}| & \leq\|\bm{E}^{(l)}\|,
\end{align}
\end{subequations} which taken together with \eqref{eq:noise-E-size-general}
validates \eqref{eq:lambda-size-large-general}-\eqref{eq:lambda-size-small-general}.

Regarding \eqref{eq:dist-U-Ustar-general} and \eqref{eq:distp-U-Ustar-general},
apply Corollary~\ref{cor:davis-kahan-conclusion-corollary} to reach
\begin{align}
\mathsf{dist}(\bm{U},\bm{U}^{\star})\leq\sqrt{2}\|\sin\bm{\Theta}\|\leq\frac{2\|\bm{E}\|}{\lambda_{r}^{\star}}\leq\frac{2c_{2}\sigma\sqrt{n}}{\lambda_{r}^{\star}}\label{eq:dist-UUstar-sin-Theta-general-UB123}
\end{align}
as long as $\|\bm{E}\|\leq c_{2}\sigma\sqrt{n}\leq(1-1/\sqrt{2})\lambda_{r}^{\star}$,
where we have used $\lambda_{r+1}^{\star}=0$. The bound on $\mathsf{dist}\big(\bm{U}^{(l)},\bm{U}^{\star}\big)$
follows from the same argument.

Additionally, regarding $\bm{U}\bm{H}-\bm{U}^{\star}$ we can derive
%
\begin{align}
 & \|\bm{U}\bm{H}-\bm{U}^{\star}\|_{\mathrm{F}}=\|\bm{U}\bm{U}^{\top}\bm{U}^{\star}-\bm{U}^{\star}\bm{U}^{\star\top}\bm{U}^{\star}\|_{\mathrm{F}}\nonumber \\
 & \qquad\quad\leq\|\bm{U}\bm{U}^{\top}-\bm{U}^{\star}\bm{U}^{\star\top}\|\,\|\bm{U}^{\star}\|_{\mathrm{F}} \nonumber\\
 &\qquad\quad=\|\sin\bm{\Theta}\|\,\|\bm{U}^{\star}\|_{\mathrm{F}}\leq\frac{2c_{2}\sigma\sqrt{rn}}{\lambda_{r}^{\star}},
\end{align}
where the last line arises from Lemma~\ref{prop:unitary-norm-property},
\eqref{eq:distp-U-Ustar-general}, and $\|\bm{U}^{\star}\|_{\mathrm{F}}=\sqrt{r}$.




\paragraph{Proof of Lemma~\ref{lem:H-property-summary-general}.}

We shall only prove the result for $\bm{H}$; the proof for $\bm{H}^{(l)}$
follows from identical arguments and is hence omitted.

From our discussion in Section~\ref{sec:introduction-distance-principal-angles},
one can express the SVD of $\bm{H}=\bm{U}^{\top}\bm{U}^{\star}$ as
$\bm{H}=\bm{X}(\cos\bm{\Theta})\bm{Y}^{\top}$, where the columns
of $\bm{X}$ (resp.~$\bm{Y}$) are the left (resp.~right) singular
vectors of $\bm{H}$, and $\bm{\Theta}$ is a diagonal matrix composed
of the principal angles between $\bm{U}$ and $\bm{U}^{\star}$. In
light of this and the definition (\ref{eq:defn-sgn-Z}), we can establish
\eqref{eq:H-sgnH-dist-general} as follows
%
\begin{align}
\big\|\bm{H}-\mathsf{sgn}(\bm{H})\big\| & =\big\|\bm{X}\big(\cos\bm{\Theta}-\bm{I}\big)\bm{Y}^{\top}\big\|=\big\|\bm{I}-\cos\bm{\Theta}\big\|\nonumber \\
 & \leq \big\|\bm{I}-\cos^{2}\bm{\Theta}\big\| = \|\sin\bm{\Theta}\|^{2}\nonumber \\
 & \leq\frac{2c_{2}^{2}\sigma^{2}n}{(\lambda_{r}^{\star})^{2}},
	\label{eq:H-sgnH-diff-UB-123}
\end{align}
%
where the middle line holds since $1-\cos\theta \leq 1 - \cos^2 \theta$,
and the last line follows from \eqref{eq:distp-U-Ustar-general}.

Coming back to the claim \eqref{eq:H-inv-norm-bound-general}, it
suffices to justify that $\sigma_{\min}(\bm{H})\geq1/2$. Recognizing
that $\mathsf{sgn}(\bm{H})=\bm{X}\bm{Y}^{\top}$, we see that all
singular values of $\mathsf{sgn}(\bm{H})$ equal 1. Thus, Weyl's inequality
together with (\ref{eq:H-sgnH-diff-UB-123}) gives
\[
\sigma_{\min}(\bm{H})\geq\sigma_{\min}\big(\mathsf{sgn}(\bm{H})\big)-\|\bm{H}-\mathsf{sgn}(\bm{H})\|\geq1-\frac{2c_{2}^{2}\sigma^{2}n}{(\lambda_{r}^{\star})^{2}}\geq\frac{1}{2},
\]
with the proviso that $2c_{2}\sigma\sqrt{n}\leq\lambda_{r}^{\star}$.


\paragraph{Proof of Lemma~\ref{lem:UH-MUstar-diff-decompose}.}
We start by connecting $\bm{U}\bm{H}\bm{\Lambda}^{\star}$ more explicitly with $\bm{M}\bm{U}^{\star}$ as follows (the invertibility of $\bm{\Lambda}$ can be deduced from \eqref{eq:lambda-size-large-general})
%
\begin{align}
	\bm{U}\bm{H}\bm{\Lambda}^{\star} & =\bm{U}\bm{\Lambda}\bm{\Lambda}^{-1}\bm{H}\bm{\Lambda}^{\star}
	%=\bm{M}\bm{U}\bm{\Lambda}^{-1}\bm{U}^{\top}\bm{U}^{\star}
	=\bm{M}\bm{U}\bm{\Lambda}^{-1}\bm{U}^{\top}\bm{U}^{\star}\bm{\Lambda}^{\star},\label{eq:UH-derivation-1}
\end{align}
%
which relies on the definition (\ref{eq:defn-H-Hl-general})
and the eigendecomposition $\bm{M}\bm{U}=\bm{U}\bm{\Lambda}$. In
addition, the eigendecomposition $\bm{M}^{\star}\bm{U}^{\star}=\bm{U}^{\star}\bm{\Lambda}^{\star}$ gives
%
\begin{align}
	\bm{U}^{\top}\bm{U}^{\star}\bm{\Lambda}^{\star}
	&=\bm{U}^{\top}\bm{M}^{\star}\bm{U}^{\star}=\bm{U}^{\top}\bm{M}\bm{U}^{\star}-\bm{U}^{\top}\bm{E}\bm{U}^{\star} \notag\\
	&= \bm{\Lambda}\bm{U}^{\top}\bm{U}^{\star}-\bm{U}^{\top}\bm{E}\bm{U}^{\star},
	\label{eq:HLambda-Lambda-H-connection}
\end{align}
%
which taken collectively with  (\ref{eq:UH-derivation-1}) demonstrates that
%
\begin{align}
\bm{U}\bm{H}\bm{\Lambda}^{\star} & =\bm{M}\bm{U}\bm{\Lambda}^{-1}\bm{\Lambda}\bm{U}^{\top}\bm{U}^{\star}-\bm{M}\bm{U}\bm{\Lambda}^{-1}\bm{U}^{\top}\bm{E}\bm{U}^{\star} \notag\\
 & =\bm{M}\bm{U}\bm{H}-\bm{M}\bm{U}\bm{\Lambda}^{-1}\bm{U}^{\top}\bm{E}\bm{U}^{\star} \notag\\
 & =\bm{M}\bm{U}^{\star}+\bm{M}\big(\bm{U}\bm{H}-\bm{U}^{\star}\big)-\bm{M}\bm{U}\bm{\Lambda}^{-1}\bm{U}^{\top}\bm{E}\bm{U}^{\star}.
\notag
%\label{eq:UH-derivation-2}
\end{align}
%
Consequently, the difference between $\bm{U}\bm{H}\bm{\Lambda}^{\star}$ and $\bm{M}\bm{U}^{\star}$ obeys
%
\begin{align}
 & \big\|\bm{U}\bm{H}\bm{\Lambda}^{\star}-\bm{M}\bm{U}^{\star}\big\|_{2,\infty}
  \leq \big\|\bm{M}\big(\bm{U}\bm{H}-\bm{U}^{\star}\big)\big\|_{2,\infty}  \notag\\
 %& \qquad\qquad\qquad\qquad\qquad\qquad
	& \qquad\qquad\qquad +\big\|\bm{M}\bm{U}\bm{\Lambda}^{-1}\bm{U}^{\top}\bm{E}\bm{U}^{\star}\big\|_{2,\infty}.
\label{eq:UH-MULambda-UB22}
\end{align}
%


%We then upper bound the two terms on the right-hand side of \eqref{eq:UH-MULambda-UB22} separately. Regarding the first term,
%we observe that
%%
%\begin{align}
%  \big\|\bm{M}\big(\bm{U}\bm{H}-\bm{U}^{\star}\big)\big(\bm{\Lambda}^{\star}\big)^{-1}\big\|_{2,\infty}
% & \leq \big\|\bm{M}\big(\bm{U}\bm{H}-\bm{U}^{\star}\big)\big\|_{2,\infty}\big\|\big(\bm{\Lambda}^{\star}\big)^{-1}\big\| \notag\\
% &=\frac{\big\|\bm{M}\big(\bm{U}\bm{H}-\bm{U}^{\star}\big)\big\|_{2,\infty}}{\lambda_{r}^{\star}} ,
%	\label{eq:M-UH-U-diff-UB1234}
%\end{align}
%%
%which follows from the assumption \eqref{eq:assumption-lambda-r-positive}.
%
Regarding the second term in \eqref{eq:UH-MULambda-UB22}, one can deduce that
%\lambda_{r}^{\star}
\begin{align}
 & \big\|\bm{M}\bm{U}\bm{\Lambda}^{-1}\bm{U}^{\top}\bm{E}\bm{U}^{\star}\big\|_{2,\infty} \notag\\
 & \quad \leq\big\|\bm{M}\bm{U}\big\|_{2,\infty}\big\|\bm{\Lambda}^{-1}\big\|\cdot\big\|\bm{U}\big\|\cdot\big\|\bm{E}\big\|\cdot\big\|\bm{U}^{\star}\big\| \notag\\
 & \quad \overset{(\mathrm{i})}{\leq}\frac{2\big\|\bm{M}\bm{U}\big\|_{2,\infty}\big\|\bm{E}\big\|}{\lambda_{r}^{\star}}\overset{(\mathrm{ii})}{\leq}\frac{4\big\|\bm{M}\bm{U}\bm{H}\big\|_{2,\infty}\big\|\bm{E}\big\|}{\lambda_{r}^{\star}} \notag\\
	& \quad \overset{(\mathrm{iii})}{\leq} \frac{4\big\|\bm{M}\big(\bm{U}\bm{H}-\bm{U}^{\star}\big)\big\|_{2,\infty}\big\|\bm{E}\big\|}{ \lambda_{r}^{\star}  }
	+ \frac{4\big\|\bm{M}\bm{U}^{\star}\big\|_{2,\infty}\big\|\bm{E}\big\|}{\lambda_{r}^{\star}} \notag\\
	& \quad \overset{(\mathrm{iv})}{\leq} \big\|\bm{M}\big(\bm{U}\bm{H}-\bm{U}^{\star}\big)\big\|_{2,\infty}
	+ \frac{4\big\|\bm{M}\bm{U}^{\star}\big\|_{2,\infty}\big\|\bm{E}\big\|}{\lambda_{r}^{\star}} .
	\label{eq:M-UH-U-diff-UB5678}
\end{align}
%
Here, (i) follows from the facts $\|\bm{U}\|=\|\bm{U}^{\star}\|=1$ and $|\lambda_r| \geq \lambda_r^{\star} - c_2\sigma\sqrt{n} \geq  \lambda_r^{\star}/2$ (see Lemma~\ref{lem:preliminary-result-general-iid}),
 (ii) holds due to \eqref{eq:A-2inf-H-Hl-general},
% the observation that $\big\|\bm{M}\bm{U}\big\|_{2,\infty}=\big\|\bm{M}\bm{U}\bm{H}\bm{H}^{-1}\big\|_{2,\infty}\leq\big\|\bm{M}\bm{U}\bm{H}\big\|_{2,\infty}\big\|\bm{H}^{-1}\big\|\leq2\big\|\bm{M}\bm{U}\bm{H}\big\|_{2,\infty}$ (given that $\|\bm{H}^{-1}\|\leq 2$ according to \eqref{eq:H-property-summary-general}),
(iii) invokes the triangle inequality,
whereas (iv) holds true provided that $4\|\bm{E}\|\leq \lambda^{\star}_r$.
Combine (\ref{eq:UH-MULambda-UB22}) and \eqref{eq:M-UH-U-diff-UB5678} to reach
\begin{align*}
 & \big\|\bm{U}\bm{H}\bm{\Lambda}^{\star}-\bm{M}\bm{U}^{\star}\big\|_{2,\infty}
  \leq 2\big\|\bm{M}\big(\bm{U}\bm{H}-\bm{U}^{\star}\big)\big\|_{2,\infty}
	+ \frac{4\big\|\bm{M}\bm{U}^{\star}\big\|_{2,\infty}\big\|\bm{E}\big\|}{\lambda_{r}^{\star}},
\end{align*}
which together with the fact that $\|(\bm{\Lambda}^\star)^{-1}\|=1/\lambda^{\star}_{r}$ and the elementary relation $\|\bm{A}\|_{2,\infty}=\|\bm{A}\bm{\Lambda}^\star(\bm{\Lambda}^\star)^{-1}\|_{2,\infty} \leq \|\bm{A}\bm{\Lambda}^\star\|_{2,\infty}\|(\bm{\Lambda}^\star)^{-1}\|$ yields the desired claim \eqref{eq:UH-MUstarLambdastar-diff-2inf-general}.





When it comes to the second claim \eqref{eq:UH-Ustar-diff-2inf-general}, combining \eqref{eq:UH-MUstarLambdastar-diff-2inf-general} with the triangle inequality
%
\begin{align*}
\big\|\bm{U}\bm{H}-\bm{U}^{\star}\big\|_{2,\infty} & = \big\|\bm{U}\bm{H}-\bm{M}^{\star}\bm{U}^{\star}\big(\bm{\Lambda}^{\star}\big)^{-1}\big\|_{2,\infty}\\
 & \leq\big\|\bm{U}\bm{H}-\bm{M}\bm{U}^{\star}\big(\bm{\Lambda}^{\star}\big)^{-1}\big\|_{2,\infty}+\big\|\bm{E}\bm{U}^{\star}\big(\bm{\Lambda}^{\star}\big)^{-1}\big\|_{2,\infty} \\
 & \leq\big\|\bm{U}\bm{H}-\bm{M}\bm{U}^{\star}\big(\bm{\Lambda}^{\star}\big)^{-1}\big\|_{2,\infty}
	+\big\|\bm{E}\bm{U}^{\star}\big\|_{2,\infty} \,/\, \lambda_r^{\star}
\end{align*}
%
immediately establishes the advertised bound. Here the last relation again arises from the elementary inequality $\|\bm{A}\bm{B}\|_{2,\infty}\leq \|\bm{A}\|_{2,\infty}\|\bm{B}\|$.






\paragraph{Proof of Lemma~\ref{lem:UH-sequence-bounds-E123}.}
First of all, taking Condition (\ref{eq:E1-UB-E11-rho1}) collectively with (\ref{eq:UH-Ustar-diff-2inf-general})
and rearranging terms yield (\ref{eq:UH-Ustar-E11-E2-E3}):
%
\begin{equation}
\big\|\bm{U}\bm{H}-\bm{U}^{\star}\big\|_{2,\infty}\leq\frac{1}{1-\rho_{1}}\big(\mathcal{E}_{1,1}+\mathcal{E}_{2}+\mathcal{E}_{3}\big)\leq2\big(\mathcal{E}_{1,1}+\mathcal{E}_{2}+\mathcal{E}_{3}\big),\label{eq:UH-Ustar-E11-E2-E3-1}
\end{equation}
where the last inequality follows from $\rho\leq 1/2$. Substituting (\ref{eq:E1-UB-E11-rho1}) and (\ref{eq:UH-Ustar-E11-E2-E3})
into (\ref{eq:UH-MUstarLambdastar-diff-2inf-general}) then gives \eqref{eq:UH-MUstar-Lambdastar-E11-E2-E3}:
%
\begin{align}
\big\|\bm{U}\bm{H}-\bm{M}\bm{U}^{\star}\big(\bm{\Lambda}^{\star}\big)^{-1}\big\|_{2,\infty} & \leq\mathcal{E}_{1,1}+2\rho_{1}\big(\mathcal{E}_{1,1}+\mathcal{E}_{2}+\mathcal{E}_{3}\big)+\mathcal{E}_{2}\nonumber \\
 & \leq2\mathcal{E}_{1,1}+2\mathcal{E}_{2}+2\rho_{1}\mathcal{E}_{3},\label{eq:UH-MUstar-Lambdastar-E11-E2-E3-1}
\end{align}
where once again we use the assumption that $\rho\leq 1/2$. 
In addition, the following observation connects $\bm{U}\bm{H}$ with
$\bm{U}\mathsf{sgn}(\bm{H})$:
%
\begin{align}
 & \big\|\bm{U}\bm{H}-\bm{U}\mathsf{sgn}(\bm{H})\big\|_{2,\infty}
   \leq \|\bm{U}\|_{2,\infty}\big\|\bm{H}-\mathsf{sgn}(\bm{H})\big\|\nonumber \\
 & \qquad \overset{(\mathrm{i})}{\leq}\frac{2c_{2}^{2}\sigma^{2}n}{(\lambda_{r}^{\star})^{2}}\big\|\bm{U}\big\|_{2,\infty}\overset{(\mathrm{ii})}{\leq}\frac{4c_{2}^{2}\sigma^{2}n}{(\lambda_{r}^{\star})^{2}}\big\|\bm{U}\bm{H}\big\|_{2,\infty}\nonumber \\
 & \qquad \overset{(\mathrm{iii})}{\leq}\frac{4c_{2}^{2}\sigma^{2}n}{(\lambda_{r}^{\star})^{2}}\big\|\bm{U}\bm{H}-\bm{U}^{\star}\big\|_{2,\infty}+\frac{4c_{2}^{2}\sigma^{2}n}{(\lambda_{r}^{\star})^{2}}\sqrt{\frac{\mu r}{n}},
 \label{eq:UH-UsgnH-bound-iterate}
\end{align}
%
where (i) results from (\ref{eq:H-sgnH-dist-general}), (ii) relies on (\ref{eq:A-2inf-H-general}), and (iii) comes from the triangle
inequality $\|\bm{U}\bm{H}\|_{2,\infty}\leq \|\bm{U}\bm{H}-\bm{U}^{\star}\|_{2,\infty} + \|\bm{U}^{\star} \|_{2,\infty}$
and the definition (\ref{eq:defn-Ustar-incoherence}).

The preceding bound
 together with the triangle inequality gives \eqref{eq:UsgnH-Ustar-Lambdastar-E11-E2-E3}:
%
\begin{align*}
\big\|\bm{U}\mathsf{sgn}(\bm{H})-\bm{U}^{\star}\big\|_{2,\infty} & \leq\big\|\bm{U}\bm{H}-\bm{U}^{\star}\big\|_{2,\infty}+\big\|\bm{U}\bm{H}-\bm{U}\mathsf{sgn}(\bm{H})\big\|_{2,\infty}\\
 & \leq\Big(1+\frac{4c_{2}^{2}\sigma^{2}n}{(\lambda_{r}^{\star})^{2}}\Big)\big\|\bm{U}\bm{H}-\bm{U}^{\star}\big\|_{2,\infty}+\frac{4c_{2}^{2}\sigma^{2}n}{(\lambda_{r}^{\star})^{2}}\sqrt{\frac{\mu r}{n}}\\
 & \leq4\big(\mathcal{E}_{1,1}+\mathcal{E}_{2}+\mathcal{E}_{3}\big)+\frac{4c_{2}^{2}\sigma^{2}\sqrt{\mu rn}}{(\lambda_{r}^{\star})^{2}},
\end{align*}
%
where the last inequality holds as long as $4c_{2}^{2}\sigma^{2}n\leq(\lambda_{r}^{\star})^{2}$. Additionally, the inequalities \eqref{eq:UH-MUstar-Lambdastar-E11-E2-E3-1} and \eqref{eq:UH-UsgnH-bound-iterate} further allow us to deduce \eqref{eq:UsgnH-MUstar-Lambdastar-E11-E2-E3}:
%
\begin{align*}
 & \big\|\bm{U}\mathsf{sgn}(\bm{H})-\bm{M}\bm{U}^{\star}\big(\bm{\Lambda}^{\star}\big)^{-1}\big\|_{2,\infty} \\
 & \quad \leq\big\|\bm{U}\bm{H}-\bm{M}\bm{U}^{\star}\big(\bm{\Lambda}^{\star}\big)^{-1}\big\|_{2,\infty}+\big\|\bm{U}\bm{H}-\bm{U}\mathsf{sgn}(\bm{H})\big\|_{2,\infty}\\
 % & \quad\leq2\mathcal{E}_{1,1}+2\mathcal{E}_{2}+2\rho_{1}\mathcal{E}_{3}+\frac{4c_{2}^{2}\sigma^{2}n}{(\lambda_{r}^{\star})^{2}}\big\|\bm{U}\bm{H}-\bm{U}^{\star}\big\|_{2,\infty}+\frac{4c_{2}^{2}\sigma^{2}n}{(\lambda_{r}^{\star})^{2}}\sqrt{\frac{\mu r}{n}}\\
 & \quad\leq2\mathcal{E}_{1,1}+2\mathcal{E}_{2}+2\rho_{1}\mathcal{E}_{3}+\frac{8c_{2}^{2}\sigma^{2}n}{(\lambda_{r}^{\star})^{2}}\big(\mathcal{E}_{1,1}+\mathcal{E}_{2}+\mathcal{E}_{3}\big)+\frac{4c_{2}^{2}\sigma^{2}\sqrt{\mu rn}}{(\lambda_{r}^{\star})^{2}}\\
 & \quad\leq3\mathcal{E}_{1,1}+3\mathcal{E}_{2}+\Big(2\rho_{1}+\frac{8c_{2}^{2}\sigma^{2}n}{(\lambda_{r}^{\star})^{2}}\Big)\mathcal{E}_{3}+\frac{4c_{2}^{2}\sigma^{2}\sqrt{\mu rn}}{(\lambda_{r}^{\star})^{2}} .
\end{align*}
%
Here, the penultimate line combines (\ref{eq:UH-Ustar-E11-E2-E3}), (\ref{eq:UH-MUstar-Lambdastar-E11-E2-E3})
and (\ref{eq:UH-UsgnH-bound-iterate}), while the last inequality relies on the
assumption $8 c_{2}^{2}\sigma^{2}n\leq(\lambda_{r}^{\star})^{2}$.



%
%\begin{align}
% & \big\|\bm{U}\bm{H}-\bm{M}\bm{U}^{\star}\big(\bm{\Lambda}^{\star}\big)^{-1}\big\|_{2,\infty} \notag \\
%	& \qquad \leq\frac{2\big\|\bm{M}\big(\bm{U}\bm{H}-\bm{U}^{\star}\big)\big\|_{2,\infty}}{\lambda_{r}^{\star}}+\frac{4\big\|\bm{M}\bm{U}^{\star}\big\|_{2,\infty}\big\|\bm{E}\big\|}{(\lambda_{r}^{\star})^{2}},
%% & \qquad\leq \frac{2\big\|\bm{E}\big(\bm{U}\bm{H}-\bm{U}^{\star}\big)\big\|_{2,\infty}}{\lambda_{r}^{\star}}  + \frac{2\big\|\bm{M}^{\star}\big(\bm{U}\bm{H}-\bm{U}^{\star}\big)\big\|_{2,\infty}}{\lambda_{r}^{\star}} \notag\\
%%	& \qquad\qquad\quad +\frac{4\Big\{ \sqrt{\frac{\mu r}{n}}\big|\lambda_{1}^{\star}\big|+(4+6c_{1})\sigma\sqrt{r\log n} \Big\} \cdot c_2 \sigma\sqrt{n}}{(\lambda_{r}^{\star})^{2}},
%\label{eq:UB-UH-MULambda-UB3456}
%\end{align}

%
%
%
%
%Substituting \eqref{eq:M-UH-U-diff-UB1234} and \eqref{eq:M-UH-U-diff-UB5678} into \eqref{eq:UH-MULambda-UB22}, we arrive at
%%
%\begin{align}
% & \big\|\bm{U}\bm{H}-\bm{M}\bm{U}^{\star}\big(\bm{\Lambda}^{\star}\big)^{-1}\big\|_{2,\infty}  \leq\frac{2\big\|\bm{M}\big(\bm{U}\bm{H}-\bm{U}^{\star}\big)\big\|_{2,\infty}}{\lambda_{r}^{\star}}+\frac{4\big\|\bm{M}\bm{U}^{\star}\big\|_{2,\infty}\big\|\bm{E}\big\|}{(\lambda_{r}^{\star})^{2}}\nonumber\\
% & \qquad\leq \frac{2\big\|\bm{E}\big(\bm{U}\bm{H}-\bm{U}^{\star}\big)\big\|_{2,\infty}}{\lambda_{r}^{\star}}  + \frac{2\big\|\bm{M}^{\star}\big(\bm{U}\bm{H}-\bm{U}^{\star}\big)\big\|_{2,\infty}}{\lambda_{r}^{\star}} \notag\\
%	& \qquad\qquad\quad +\frac{4\Big\{ \sqrt{\frac{\mu r}{n}}\big|\lambda_{1}^{\star}\big|+(4+6c_{1})\sigma\sqrt{r\log n} \Big\} \cdot c_2 \sigma\sqrt{n}}{(\lambda_{r}^{\star})^{2}},
%\label{eq:UB-UH-MULambda-UB12345}
%\end{align}
%%
%%
%%\begin{align}
%% & \big\|\bm{U}\bm{H}-\bm{M}\bm{U}^{\star}\big(\bm{\Lambda}^{\star}\big)^{-1}\big\|_{2,\infty}\nonumber \\
%% & \leq\frac{2\big\|\bm{M}\big(\bm{U}\bm{H}-\bm{U}^{\star}\big)\big\|_{2,\infty}}{\lambda_{r}^{\star}}+\frac{4\big\|\bm{M}\bm{U}^{\star}\big\|_{2,\infty}\big\|\bm{E}\big\|}{(\lambda_{r}^{\star})^{2}}\nonumber \\
%% &  \leq\frac{2\big\|\bm{M}\big(\bm{U}\bm{H}-\bm{U}^{\star}\big)\big\|_{2,\infty}}{\lambda_{r}^{\star}}
%%	+\frac{4\big\{\sqrt{\frac{\mu r}{n}} \|\bm{\Lambda}^{\star}\|+\big\|\bm{E}\bm{U}^{\star}\big\|_{2,\infty}\big\}\big\|\bm{E}\big\|}{(\lambda_{r}^{\star})^{2}},
%% \label{eq:UB-UH-MULambda-UB12345}
%%\end{align}
%%
%
%






\input{chapters/analysis_entrywise.tex}
\section{Appendix C: Proof of Theorem~\ref{thm:2-inf-asymm}}
\label{sec:proof-inf-rect}

\paragraph{A symmetrization trick. }

As alluded to previously, the proof is built  on a ``symmetric dilation'' trick that helps symmetrize a general matrix. We start with the following definition.  

\begin{definition}[Symmetric dilation]
\label{defn:sym-dilation}
For any matrix $\bm{A}\in\mathbb{R}^{n_{1}\times n_{2}}$, its symmetric dilation $\mathcal{S}(\bm{A})\in\mathbb{R}^{(n_{1}+n_{2})\times(n_{1}+n_{2})}$ is defined to be 
%
\[
\mathcal{S}(\bm{A})=\left[\begin{array}{cc}
\bm{0} & \bm{A}\\
\bm{A}^{\top} & \bm{0}
\end{array}\right].
\]
%
 \end{definition}
 %
Apart from the symmetry of $\mathcal{S}(\bm{A})$,
which is immediate from its definition, the main benefit of the symmetric
dilation lies in the correspondence between the eigendecomposition
of $\mathcal{S}(\bm{A})$ and the singular value decomposition of
$\bm{A}$. More specifically, let $\bm{U}\bm{\Sigma}\bm{V}^{\top}$
be the SVD of $\bm{A}$. Then one has the following eigendecomposition
for $\mathcal{S}(\bm{A})$: 
\begin{equation}
\mathcal{S}(\bm{A})=\frac{1}{\sqrt{2}}\left[\begin{array}{cc}
\bm{U} & \bm{U}\\
\bm{V} & -\bm{V}
\end{array}\right]\cdot\left[\begin{array}{cc}
\bm{\Sigma} & \bm{0}\\
\bm{0} & -\bm{\Sigma}
\end{array}\right]\cdot\frac{1}{\sqrt{2}}\left[\begin{array}{cc}
\bm{U} & \bm{U}\\
\bm{V} & -\bm{V}
\end{array}\right]^{\top}. \label{eq:dilation-eigen}
\end{equation}
%
Here, the columns of $\frac{1}{\sqrt{2}}\left[{\scriptsize \begin{array}{cc}
\bm{U} & \bm{U}\\
\bm{V} & -\bm{V}
\end{array}}\right]$
%
are orthonormal and represent the eigenvectors of $\mathcal{S}(\bm{A})$, whereas
$\left[{\scriptsize \begin{array}{cc}
\bm{\Sigma} & \bm{0}\\
\bm{0} & -\bm{\Sigma}
\end{array}}\right]$ contains all (non-zero) eigenvalues
of $\mathcal{S}(\bm{A})$. 

Utilizing this ``symmetric dilation'' trick, we can translate the observation
model $\bm{M}=\bm{M}^{\star}+\bm{E}$ into the following equivalent
form
\[
\mathcal{S}(\bm{M})=\mathcal{S}(\bm{M}^{\star})+\mathcal{S}(\bm{E}),
\]
which is in line with the symmetric observation model stated in
(\ref{eq:M-noisy-copy-general}). 



\paragraph{Verifying conditions. }

% (cf.~Theorem~\ref{thm:UsgnH-Ustar-MUstar-general} and Corollary~\ref{cor:entrywise-error-general})
To invoke the general theory in Section \ref{sec:setup-general-theory-independent},
one is required to first examine the spectral properties of $\mathcal{S}(\bm{M}^{\star})$
and $\mathcal{S}(\bm{M})$, as well as the  assumptions on the noise part
$\mathcal{S}(\bm{E})$. 

%from Section~\ref{sec:general-theory-setup-asymm} 

Recall that $\bm{M}^{\star}=\bm{U}^{\star}\bm{\Sigma}^{\star}\bm{V}^{\star\top}$, which together with the relation (\ref{eq:dilation-eigen}) reveals that: (i) $\mathcal{S}(\bm{M}^{\star})$ has rank $2r$ and condition number $\kappa$; (ii) the nonzero eigenvalues of $\mathcal{S}(\bm{M}^{\star})$ and the corresponding eigenvectors are reflected respectively in the  matrices
	%
\begin{equation}
\overline{\bm{\Lambda}}^{\star}\coloneqq\left[\begin{array}{cc}
\bm{\Sigma}^{\star} & \bm{0}\\
\bm{0} & -\bm{\Sigma}^{\star}
\end{array}\right]	
\quad\text{and}\quad
\overline{\bm{U}}^{\star}\coloneqq\frac{1}{\sqrt{2}}\left[\begin{array}{cc}
\bm{U}^{\star} & \bm{U}^{\star}\\
\bm{V}^{\star} & -\bm{V}^{\star}
\end{array}\right]
	\label{eq:defn-spectral-S-M-star}
\end{equation}
%
Similarly, given that the SVD of $\bm{M}$ is $\bm{M}=\bm{U}\bm{\Sigma}\bm{V}^{\top}+\bm{U}_{\perp}\bm{\Sigma}_{\perp}\bm{V}_{\perp}^{\top}$, we see that the $2r$-leading eigenvalues of $\mathcal{S}(\bm{M})$ and the corresponding eigenvectors are represented respectively by the matrices
\[
\overline{\bm{\Lambda}}\coloneqq\left[\begin{array}{cc}
\bm{\Sigma} & \bm{0}\\
\bm{0} & -\bm{\Sigma}
\end{array}\right]
\quad\text{and}\quad
\overline{\bm{U}}\coloneqq\frac{1}{\sqrt{2}}\left[\begin{array}{cc}
\bm{U} & \bm{U}\\
\bm{V} & -\bm{V}
\end{array}\right]. 
\]
%
 Further, the incoherence parameter $\overline{\mu}$
of $\mathcal{S}(\bm{M}^{\star})$ (cf.~(\ref{eq:defn-Ustar-incoherence}))
obeys
\begin{align}
\overline{\mu} & \coloneqq\frac{(n_{1}+n_{2})}{2r}\|\overline{\bm{U}}^{\star}\|_{2,\infty}^{2}
	\overset{(\text{i})}{=}  \frac{(n_{1}+n_{2})}{2r}\max\big\{\|\bm{U}^{\star}\|_{2,\infty}^{2},\|\bm{V}^{\star}\|_{2,\infty}^{2}\big\} \notag\\
 & \overset{(\text{ii})}{\leq}\frac{(n_{1}+n_{2})}{\min\{n_{1},n_{2}\}}\frac{\mu}{2}\overset{(\text{iii})}{=}\frac{(n_{1}+n_{2})}{2n_{1}}\mu.
	\label{eq:mu-bar-general}
\end{align}
Here, the relation (i) is based on the definition (\ref{eq:defn-spectral-S-M-star}),
the  inequality (ii) follows from the incoherence of $\bm{M}^{\star}$
(cf.~Definition~\ref{assump:mc-incoherence}), while the last one (iii)
holds under the assumption $n_{1}\leq n_{2}$. 

When it comes to the ``symmetrized'' noise part, it is straightforward to verify
that under Assumption~\ref{assump:noise-general-rect}, 
the matrix
$\mathcal{S}(\bm{E})$ satisfies Assumption~\ref{assumption-noise-general}
with precisely the quantities $\sigma, B$ and $c_{\mathsf{b}}$. 



%
%(resp.~$\bm{M}=\bm{U}\bm{\Sigma}\bm{V}^{\top}+\bm{U}_{\perp}\bm{\Sigma}_{\perp}\bm{V}_{\perp}^{\top}$)
%is the SVD of $\bm{M}^{\star}$ (resp.~$\bm{M}$), which together
%with the relation (\ref{eq:dilation-eigen}) reveals the eigendecompositions
%of $\mathcal{S}(\bm{M}^{\star})$ and $\mathcal{S}(\bm{M})$: 
%\begin{align*}
%\mathcal{S}(\bm{M}^{\star}) & =\frac{1}{\sqrt{2}}\left[\begin{array}{cc}
%\bm{U}^{\star} & \bm{U}^{\star}\\
%\bm{V}^{\star} & -\bm{V}^{\star}
%\end{array}\right]\cdot\left[\begin{array}{cc}
%\bm{\Sigma}^{\star} & \bm{0}\\
%\bm{0} & -\bm{\Sigma}^{\star}
%\end{array}\right]\cdot\frac{1}{\sqrt{2}}\left[\begin{array}{cc}
%\bm{U}^{\star} & \bm{U}^{\star}\\
%\bm{V}^{\star} & -\bm{V}^{\star}
%\end{array}\right]^{\top};\\
%\mathcal{S}(\bm{M}) & =\frac{1}{\sqrt{2}}\left[\begin{array}{cccc}
%\bm{U} & \bm{U} & \bm{U}_{\perp} & \bm{U}_{\perp}\\
%\bm{V} & -\bm{V} & \bm{V}_{\perp} & -\bm{V}_{\perp}
%\end{array}\right]\cdot\left[\begin{array}{cccc}
%\bm{\Sigma}\\
% & -\bm{\Sigma}\\
% &  & \bm{\Sigma}_{\perp}\\
% &  &  & -\bm{\Sigma}_{\perp}
%\end{array}\right]\\
% & \quad\cdot\frac{1}{\sqrt{2}}\left[\begin{array}{cccc}
%\bm{U} & \bm{U} & \bm{U}_{\perp} & \bm{U}_{\perp}\\
%\bm{V} & -\bm{V} & \bm{V}_{\perp} & -\bm{V}_{\perp}
%\end{array}\right]^{\top}.
%\end{align*}
%In words, $\mathcal{S}(\bm{M}^{\star})$ has rank $2r$, and 
%%
%
%are the leading eigenvectors and eigenvalues (in magnitude) of $\mathcal{S}(\bm{M}^{\star})$,
%respectively. Similarly, 
%\[
%\overline{\bm{U}}\coloneqq\frac{1}{\sqrt{2}}\left[\begin{array}{cc}
%\bm{U} & \bm{U}\\
%\bm{V} & -\bm{V}
%\end{array}\right],\qquad\text{and}\qquad\overline{\bm{\Lambda}}\coloneqq\left[\begin{array}{cc}
%\bm{\Sigma} & \bm{0}\\
%\bm{0} & -\bm{\Sigma}
%\end{array}\right]
%\]
%are the leading eigenvectors and eigenvalues (in magnitude) of $\mathcal{S}(\bm{M})$,
%respectively. Here the last assertion arises since $\bm{U}\bm{\Sigma}\bm{V}^{\top}$
%is the top-$r$ SVD of $\bm{M}$, and hence the singular values in
%$\bm{\Sigma}$ are larger than those in $\bm{\Sigma}_{\perp}$. In
%addition, the condition number of $\mathcal{S}(\bm{M}^{\star})$ (see
%the definition in (\ref{eq:defn-kappa})) corresponds exactly to that
%of $\bm{M}^{\star}$.




\paragraph{$\ell_{2,\infty}$ and $\ell_{\infty}$ guarantees. }

With the above preparations in place, 
%we are ready to invoke Theorem~\ref{thm:UsgnH-Ustar-MUstar-general} and Corollary~\ref{cor:entrywise-error-general}
%to demonstrate finer performance guarantees for the spectral
%estimates $\bm{U},\bm{V}$.
%
apply Theorem~\ref{thm:UsgnH-Ustar-MUstar-general} (more specifically (\ref{eq:UH-MUstar-Lambda-star-2inf}))
 to demonstrate that
 %
 \begin{align}
\|\overline{\bm{U}}\,\overline{\bm{U}}^{\top}\overline{\bm{U}}^{\star}-\overline{\bm{U}}^{\star}\|_{2,\infty} & \lesssim\frac{\sigma\kappa\sqrt{\overline{\mu}r}+\sigma\sqrt{r\log n}}{\sigma_{r}^{\star}} \notag\\
 & \lesssim \frac{\sigma\kappa\sqrt{n_{2}\mu r/n_{1}}+\sigma\sqrt{r\log n}}{\sigma_{r}^{\star}}, 
	 \label{eq:mc-inf-first-step}
\end{align}
%
where the last relation follows from \eqref{eq:mu-bar-general}. Further, note that 
%
\begin{align*}
\overline{\bm{U}}\,\overline{\bm{U}}^{\top}\overline{\bm{U}}^{\star}-\overline{\bm{U}}^{\star} & =\frac{1}{\sqrt{2}}\left[\begin{array}{cc}
\bm{U}\bm{H}_{\bm{U}}-\bm{U}^{\star}, & \bm{U}\bm{H}_{\bm{U}}-\bm{U}^{\star}\\
\bm{V}\bm{H}_{\bm{V}}-\bm{V}^{\star}, & -(\bm{V}\bm{H}_{\bm{V}}-\bm{V}^{\star})
\end{array}\right],
%\overline{\bm{U}}\overline{\bm{U}}^{\top}\overline{\bm{U}}^{\star}-\overline{\bm{U}}^{\star}  =\frac{1}{\sqrt{2}}\left[\begin{array}{cc}
%\bm{U}\bm{U}^{\top}\bm{U}^{\star}-\bm{U}^{\star}, & \bm{U}\bm{U}^{\top}\bm{U}^{\star}-\bm{U}^{\star}\\
%\bm{V}\bm{V}^{\top}\bm{V}^{\star}-\bm{V}^{\star}, & -(\bm{V}\bm{V}^{\top}\bm{V}^{\star}-\bm{V}^{\star})
%\end{array}\right],
\end{align*}
%
where $\bm{H}_{\bm{U}}\coloneqq \bm{U}^{\top} \bm{U}^{\star}$ and $\bm{H}_{\bm{V}}\coloneqq \bm{V}^{\top} \bm{V}^{\star}$. Combining this with the upper bound (\ref{eq:mc-inf-first-step}) then yields
%
\[
	\max\big\{\|\bm{U}\bm{H}_{\bm{U}}-\bm{U}^{\star}\|_{2,\infty},\|\bm{V}\bm{H}_{\bm{V}}-\bm{V}^{\star}\|_{2,\infty} \big\}
	\lesssim \frac{\sigma\kappa\sqrt{n_{2}\mu r/n_{1}}+\sigma\sqrt{r\log n}}{\sigma_{r}^{\star}}.
\]
%
We can then repeat the same analysis  as in the proof of Lemma~\ref{lem:UH-sequence-bounds-E123}
to connect $\|\bm{U}\bm{H}_{\bm{U}}-\bm{U}^{\star}\|_{2,\infty}$
(resp.~$\|\bm{V}\bm{H}_{\bm{V}}-\bm{V}^{\star}\|_{2,\infty}$)
with $\|\bm{U}\mathsf{sgn}(\bm{H}_{\bm{U}})-\bm{U}^{\star}\|_{2,\infty}$
(resp.~$\|\bm{V}\mathsf{sgn}(\bm{H}_{\bm{V}})-\bm{V}^{\star}\|_{2,\infty}$),
and obtain the desired bound in (\ref{eq:main-result-2-infty-general-asymm}); the details are omitted here for the sake of conciseness. 

We now proceed to the second claim (\ref{eq:main-result-infty-general-asymm}).
Towards this, invoke Corollary~\ref{cor:entrywise-error-general} and the inequality \eqref{eq:mu-bar-general} to obtain 
\[
	\big\|\overline{\bm{U}}\,\overline{\bm{\Lambda}}\,\overline{\bm{U}}^{\top}-\mathcal{S}(\bm{M}^{\star}) \big\|_{\infty}\lesssim\sigma\kappa^{2}\overline{\mu}r\sqrt{\frac{\log n}{n}}
	\lesssim \sigma\kappa^{2}\mu r\sqrt{\frac{n\log n}{n_{1}^{2}}}. 
\]
This in conjunction with the following observations 
%
\begin{align*}
	\overline{\bm{U}}\,\overline{\bm{\Lambda}}\,\overline{\bm{U}}^{\top} & =\mathcal{S}\big(\bm{U}\bm{\Sigma}\bm{V}^{\top} \big) \\
	\big\|\mathcal{S} \big(\bm{U}\bm{\Sigma}\bm{V}^{\top}\big)- \mathcal{S}\big(\bm{M}^{\star}\big) \big\|_{\infty} &=  \big\|\bm{U}\bm{\Sigma}\bm{V}^{\top}-\bm{M}^{\star} \big\|_{\infty}
\end{align*}
%
immediately establishes the second claim.


\section{Appendix D: Proof of Theorem~\ref{thm:matrix-distribution-complete}}
\label{sec:proof-thm:matrix-distribution-complete}




As before (see \eqref{eq:assumption-lambda-r-positive}), we assume throughout the proof that $\lambda_r^{\star}>0$ for the purpose of simplifying notation. 




\subsection{Proof outline}


We now outline the proof of our distributional guarantees in Theorem~\ref{thm:matrix-distribution-complete}. 
The first step consists of justifying the heuristic first-order approximation in \eqref{eq:approx-ULambdaU-Mstar-discuss}. This is stated in the lemma below, with the proof deferred to Section~\ref{sec:sub-proof-thm:matrix-distribution}. 
%
\begin{lemma}[First-order approximation]
\label{thm:matrix-distribution}
%
Suppose that the assumptions of Theorem~\ref{thm:UsgnH-Ustar-MUstar-general} hold. Then with probability at least $1-O(n^{-5})$, one can write
%
\begin{subequations}
\label{eq:U-M-decomposition-thm}
\begin{align}
\bm{U}\mathsf{sgn}(\bm{H})-\bm{U}^{\star} & =\underset{\eqqcolon\,\bm{Z}}{\underbrace{\bm{E}\bm{U}^{\star}\big(\bm{\Lambda}^{\star}\big)^{-1}}}+\bm{\Psi},
	\label{eq:U-Ustar-decomposition-Z-Psi-thm}\\
\bm{M}-\bm{M}^{\star} & =\underset{\eqqcolon\,\bm{W}}{\underbrace{\bm{E}\bm{U}^{\star}\bm{U}^{\star\top}+\bm{U}^{\star}\bm{U}^{\star\top}\bm{E}}}+\bm{\Phi}
	\label{eq:M-Mstar-diff-decomposition-thm}
\end{align}
\end{subequations}
%
for some matrices $\bm{\Psi}$ and $\bm{\Phi}$ obeying
%
\begin{subequations}
\label{eq:Psi-Phi-bound-thm}
\begin{align}
\|\bm{\Psi}\|_{2,\infty} & \lesssim\frac{\sigma^{2}\kappa\sqrt{\mu rn\log^{2}n}}{(\lambda_{r}^{\star})^{2}}+\frac{\kappa\sigma^{2}\sqrt{\mu rn}}{(\lambda_{r}^{\star})^{2}}+\frac{\sigma}{\lambda_{r}^{\star}}\sqrt{\frac{\mu r^{2}\log n}{n}},\\
\big\|\bm{\Phi}\big\|_{\infty} & \lesssim\frac{\sigma^{2}\mu r\kappa^{2}\log n}{\lambda_{r}^{\star}}+\frac{\sigma\kappa\mu\sqrt{r^{3}\log n}}{n}.
\end{align}
\end{subequations}
%
\end{lemma}
%
\begin{remark}
	In addition to quantifying the goodness of the approximation \eqref{eq:M-Mstar-diff-decomposition-thm}, Lemma~\ref{thm:matrix-distribution} also delivers a more refined characterization for the first-order approximation $\bm{U}\mathsf{sgn}(\bm{H})-\bm{U}^{\star}\approx \bm{E}\bm{U}^{\star}\big(\bm{\Lambda}^{\star}\big)^{-1}$ in comparison to Theorem~\ref{thm:UsgnH-Ustar-MUstar-general}.  As it turns out, this result \eqref{eq:U-Ustar-decomposition-Z-Psi-thm} also assists in performing statistical inference on the low-rank factors $\bm{U}^{\star}$. The interested reader is referred  to \citet{yan2021inference} for details. 
\end{remark}


In turn, Lemma~\ref{thm:matrix-distribution} motivates one to pin down the distribution of the matrix $\bm{W}$ in \eqref{eq:M-Mstar-diff-decomposition-thm}. This can be accomplished by invoking the Berry-Esseen Theorem (e.g., \citet[Theorem 3.7]{chen2010normal}), which gives rise to the following distributional characterization. The proof of this lemma can be found in Section~\ref{sec:proof-lem:normality-spectral}. 
%
\begin{lemma}[Gaussian approximation]
	\label{lem:normality-spectral}
	Suppose that the assumptions of Theorem~\ref{thm:UsgnH-Ustar-MUstar-general} hold, and that 
	%
	%assume that
%
\begin{align}
	\label{eq:Uj-Ui-lower-bound-lem-2}
	\frac{\big\|\bm{U}_{j,\cdot}^{\star}\big\|_{2}^{2}+\big\|\bm{U}_{i,\cdot}^{\star}\big\|_{2}^{2}}{\big\|\bm{U}^{\star}\big\|_{\mathrm{F}}^{2}}
	& \gtrsim \frac{B^{2}\kappa^2 \mu^{2}r^2\log ^2 n}{\sigma_{\min}^{2}n^{2}}
	+ \frac{\sigma^{4}\mu^{2}r\kappa^{4}\log^{3}n}{\sigma_{\min}^{2}(\lambda_{r}^{\star})^{2}} .
\end{align}
%
%Condition~\ref{eq:Uj-Ui-lower-bound-lem-2} holds. 
%
	Let $\bm{W}=\bm{E}\bm{U}^{\star}\bm{U}^{\star\top} + \bm{U}^{\star}\bm{U}^{\star\top}\bm{E}$. For any $1\leq i,j\leq n$, one has
	%
	\begin{subequations}
	\begin{align}
		\sup_{z\in\mathbb{R}}\left|\mathbb{P}\left( W_{i,j}\leq z \sqrt{v_{i,j}^{\star}}\right) -\Phi(z)\right| &=o(1) , \label{eq:eqn-normality-spectral}\\
		\big\|\bm{\Phi}\big\|_{\infty} &= o\left( \sqrt{ v_{i,j}^{\star} } \right), 
		\label{eq:Phi-inf-norm-o-vij}
	\end{align}
	\end{subequations}
	%
	where $v_{i,j}^{\star}$ is defined in \eqref{eq:variance-formula-Mij-lem}, and $\Phi(\cdot)$ denotes the CDF of the standard Gaussian distribution.  
\end{lemma}
%



To finish up, invoke Lemma~\ref{thm:matrix-distribution} and \eqref{eq:Phi-inf-norm-o-vij} to yield
%
\[
	M_{i,j}-M_{i,j}^{\star}=W_{i,j}+\delta_{\phi}\sqrt{v_{i,j}^{\star}}\qquad\text{with }\delta_{\phi}=o(1).
\]
%
A little algebra further gives
%
\begin{align*}
 & \left|\mathbb{P}\left( M_{i,j}-M_{i,j}^{\star}\leq z\sqrt{v_{i,j}^{\star}}\right) -\Phi(z)\right|=\left|\mathbb{P}\left( W_{i,j}\leq\left(z-\delta_{\phi}\right)\sqrt{v_{i,j}^{\star}}\right) -\Phi(z)\right|\\
 & \qquad \leq \left|\mathbb{P}\left( W_{i,j}\leq\left(z-\delta_{\phi}\right)\sqrt{v_{i,j}^{\star}}\right)  -\Phi(z-\delta_{\phi})\right|+\left|\Phi(z-\delta_{\phi})-\Phi(z)\right|\\
 & \qquad\leq o(1)+\delta_{\phi}=o(1),
\end{align*}
%
where the last line invokes Lemma~\ref{lem:normality-spectral} and the fact that $|\Phi(u)-\Phi(v)|\leq |u-v|$ for any $u,v\in \mathbb{R}$. 
This completes the proof of  Theorem~\ref{thm:matrix-distribution-complete}, as long as Lemmas~\ref{thm:matrix-distribution} and \ref{lem:normality-spectral} can be established. 
The rest of this section is thus devoted to proving Lemmas~\ref{thm:matrix-distribution} and \ref{lem:normality-spectral}. 






\subsection{Proof of Lemma~\ref{thm:matrix-distribution}}
\label{sec:sub-proof-thm:matrix-distribution}



Before proceeding to the proof, we make note of several preliminary facts that are all direct consequences of the analysis of Theorem~\ref{thm:UsgnH-Ustar-MUstar-general} and Corollary~\ref{cor:entrywise-error-general}. The proof of these preliminary results can be found in Section~\ref{sec:proof-auxiliary-distribution}.
%
\begin{lemma}
\label{lem:bound-Delta1-Delta2}
%
With probability exceeding $1-O(n^{-5})$, one can write
%
\begin{subequations}\label{eq:decomp_soda}
\begin{align}
\bm{U}\bm{\Lambda}\mathsf{sgn}(\bm{H}) & =\bm{M}\bm{U}^{\star}+\bm{\Delta}_{1} ,
			\label{eq:U-Lambda-sgnH-decompose}\\
\mathsf{sgn}(\bm{H})\bm{\Lambda}^{\star}\mathsf{sgn}(\bm{H})^{\top} & =\bm{\Lambda}+\bm{\Delta}_{2}
			\label{eq:sgnH-Lambdastar-decompose}
\end{align}
\end{subequations}
%
for some matrices $\bm{\Delta}_{1}$ and $\bm{\Delta}_{2}$ obeying
%
\begin{subequations}
\begin{align}
  \big\| \bm{\Delta}_1 \big\|_{2,\infty} & \lesssim 
\frac{\sigma^{2}\kappa\sqrt{\mu rn\log^{2}n}}{\lambda_{r}^{\star}}, \label{eq:Delta-1-U-R-Lambda}\\
  \big\|\bm{\Delta}_2\big\|
%\\
% & \qquad\leq\big\|\mathsf{sgn}(\bm{H})\bm{\Lambda}^{\star}\mathsf{sgn}(\bm{H})^{\top}-\bm{H}\bm{\Lambda}^{\star}\bm{H}^{\top}\big\|+\big\|\bm{H}\bm{\Lambda}^{\star}\bm{H}^{\top}-\bm{\Lambda}\big\|\\
 %& \qquad
  &\lesssim\frac{\kappa\sigma^{2}n}{\lambda_{r}^{\star}}+\sigma\sqrt{r\log n}.
   \label{eq:sgn-H-Lambdastar-sgn-H-Lambda}
\end{align}
\end{subequations}
%
\end{lemma}
%



We are now ready to embark on the proof of Lemma~\ref{thm:matrix-distribution}. 
In order to analyze the behavior of $\bm{U}\mathsf{sgn}(\bm{H})$, 
we first point out the following decomposition: 
%
\begin{align*}
\bm{U}\mathsf{sgn}(\bm{H})\bm{\Lambda}^{\star} & =\bm{U}\bm{\Lambda}\mathsf{sgn}(\bm{H})+\bm{U}\big(\mathsf{sgn}(\bm{H})\bm{\Lambda}^{\star}-\bm{\Lambda}\mathsf{sgn}(\bm{H})\big)\\
 & =\bm{M}\bm{U}^{\star}+\bm{\Delta}_{1}+\bm{U}\big(\mathsf{sgn}(\bm{H})\bm{\Lambda}^{\star}\mathsf{sgn}(\bm{H})^{\top}-\bm{\Lambda}\big)\mathsf{sgn}(\bm{H})\\
 & =\bm{M}^{\star}\bm{U}^{\star}+\bm{E}\bm{U}^{\star}+\bm{\Delta}_{1}+\bm{U}\bm{\Delta}_{2} \, \mathsf{sgn}(\bm{H})\\
 & =\bm{U}^{\star}\bm{\Lambda}^{\star}+\bm{E}\bm{U}^{\star}+\bm{\Delta}_{1}+\bm{U}\bm{\Delta}_{2}\, \mathsf{sgn}(\bm{H}),
\end{align*}
%
where the second and the third identities result from Lemma~\ref{lem:bound-Delta1-Delta2} (cf.~\eqref{eq:decomp_soda}). 
The key point of this decomposition is to establish a connection between $\bm{U}\mathsf{sgn}(\bm{H})\bm{\Lambda}^{\star}$ and $\bm{U}^{\star}\bm{\Lambda}^{\star} + \bm{E}\bm{U}^{\star}$, with the assistance of the matrices $\bm{\Delta}_1$ and $\bm{\Delta}_2$ studied in Lemma~\ref{lem:bound-Delta1-Delta2}. 
A little algebra then yields
%
\begin{align}
\bm{U}\mathsf{sgn}(\bm{H})-\bm{U}^{\star}=\underset{\eqqcolon\,\bm{Z}}{\underbrace{\bm{E}\bm{U}^{\star}\big(\bm{\Lambda}^{\star}\big)^{-1}}} & +\underset{\eqqcolon\,\bm{\Psi}}{\underbrace{\bm{\Delta}_{1}\big(\bm{\Lambda}^{\star}\big)^{-1}+\bm{U}\bm{\Delta}_{2}\,\mathsf{sgn}(\bm{H})\big(\bm{\Lambda}^{\star}\big)^{-1}}}.
\label{eq:U-Ustar-decomposition-Z-Psi}
\end{align}
%
In view of Lemma~\ref{lem:bound-Delta1-Delta2}, the residual matrix
$\bm{\Psi}$ obeys
%
\begin{align}
\|\bm{\Psi}\|_{2,\infty} & \leq\big\|\bm{\Delta}_{1}\big\|_{2,\infty}\big\|\big(\bm{\Lambda}^{\star}\big)^{-1}\big\|+\big\|\bm{U}\big\|_{2,\infty}\big\|\bm{\Delta}_{2}\big\|\,\big\|\mathsf{sgn}(\bm{H})\big\|\,\big\|\big(\bm{\Lambda}^{\star}\big)^{-1}\big\|\nonumber \\
 & \lesssim\frac{1}{\lambda_{r}^{\star}}\big\|\bm{\Delta}_{1}\big\|_{2,\infty}+\frac{1}{\lambda_{r}^{\star}}\sqrt{\frac{\mu r}{n}}\big\|\bm{\Delta}_{2}\big\|\nonumber \\
 & \lesssim\frac{\sigma^{2}\kappa\sqrt{\mu rn\log^{2}n}}{(\lambda_{r}^{\star})^{2}}+\frac{\kappa\sigma^{2}\sqrt{\mu rn}}{(\lambda_{r}^{\star})^{2}}+\frac{\sigma}{\lambda_{r}^{\star}}\sqrt{\frac{\mu r^{2}\log n}{n}}
\label{eq:Psi-two-inf-norm-bound}
\end{align}
%
as claimed. 



The next step lies in analyzing the matrix estimator $\bm{M}=\bm{U}\bm{\Lambda}\bm{U}^{\top}$.
Towards this, we make the observation that
%
\begin{align}
\bm{M}-\bm{M}^{\star} & =\bm{U}\bm{\Lambda}\bm{U}^{\top}-\bm{M}^{\star}\nonumber \\
 & =\bm{U}\big(\mathsf{sgn}(\bm{H})\bm{\Lambda}^{\star}\mathsf{sgn}(\bm{H})^{\top}\big)\bm{U}^{\top}-\bm{U}\bm{\Delta}_{2}\bm{U}^{\top}-\bm{M}^{\star}\nonumber \\
 & =(\bm{U}^{\star}+\bm{Z}+\bm{\Psi})\bm{\Lambda}^{\star}(\bm{U}^{\star}+\bm{Z}+\bm{\Psi})^{\top}-\bm{U}\bm{\Delta}_{2}\bm{U}^{\top}-\bm{U}^{\star}\bm{\Lambda}^{\star}\bm{U}^{\star\top}\nonumber \\
 & =\bm{Z}\bm{\Lambda}^{\star}\bm{U}^{\star\top}+\bm{U}^{\star}\bm{\Lambda}^{\star}\bm{Z}^{\top}+\bm{\Phi}\nonumber \\
 & =\bm{E}\bm{U}^{\star}\bm{U}^{\star\top}+\bm{U}^{\star}\bm{U}^{\star\top}\bm{E}+\bm{\Phi},\label{eq:M-Mstar-diff-decomposition}
\end{align}
%
where the second line relies on Lemma~\ref{lem:bound-Delta1-Delta2} (cf.~\eqref{eq:sgnH-Lambdastar-decompose}),
the third identity makes use of (\ref{eq:U-Ustar-decomposition-Z-Psi}),
and the residual matrix $\bm{\Phi}$ is defined as
%
\begin{equation}
\bm{\Phi}\coloneqq\bm{\Psi}\bm{\Lambda}^{\star}\big(\bm{U}\mathsf{sgn}(\bm{H})\big)^{\top}+\bm{U}^{\star}\bm{\Lambda}^{\star}\bm{\Psi}^{\top}+\bm{Z}\bm{\Lambda}^{\star}\big(\bm{U}\mathsf{sgn}(\bm{H})-\bm{U}^{\star}\big)^{\top}-\bm{U}\bm{\Delta}_{2}\bm{U}^{\top}.\label{eq:defn-Phi-residual}
\end{equation}
%
In addition, it is seen from (\ref{eq:EU-2-inf-bound-general}) that
%
\begin{equation}
\|\bm{Z}\|_{2,\infty}\leq\big\|\bm{E}\bm{U}^{\star}\big(\bm{\Lambda}^{\star}\big)^{-1}\big\|_{2,\infty}\leq\frac{1}{\lambda_{r}^{\star}}\big\|\bm{E}\bm{U}^{\star}\big\|_{2,\infty}\lesssim\frac{\sigma\sqrt{r\log n}}{\lambda_{r}^{\star}}.\label{eq:Z-two-inf-norm-bound}
\end{equation}
%
Consequently, one can deduce that
%
\begin{align*}
\big\|\bm{\Phi}\big\|_{\infty} & \leq\big\|\bm{\Lambda}^{\star}\big\|\left\{ \|\bm{Z}\|_{2,\infty}^{2}+\|\bm{Z}\|_{2,\infty}\big\|\bm{U}\mathsf{sgn}(\bm{H})-\bm{U}^{\star}\big\|_{2,\infty}\right\} \\
 & \quad+\big\|\bm{\Lambda}^{\star}\big\|\|\bm{\Psi}\|_{2,\infty}\left(\big\|\bm{U}\big\|_{2,\infty}+\|\bm{U}^{\star}\|_{2,\infty}\right)+\|\bm{U}\|_{2,\infty}^{2}\|\bm{\Delta}_{2}\|\\
 & \lesssim\big|\lambda_{1}^{\star}\big|\left\{ \frac{\sigma^{2}r\log n}{(\lambda_{r}^{\star})^{2}}+\frac{\sigma\sqrt{r\log n}}{\lambda_{r}^{\star}}\cdot\frac{\sigma\kappa\sqrt{\mu r\log n}}{\lambda_{r}^{\star}}\right\} \\
 & \quad+\big|\lambda_{1}^{\star}\big|\sqrt{\frac{\mu r}{n}}\left\{ \frac{\sigma^{2}\kappa\sqrt{\mu rn\log^{2}n}}{(\lambda_{r}^{\star})^{2}}+\frac{\kappa\sigma^{2}\sqrt{\mu rn}}{(\lambda_{r}^{\star})^{2}}+\frac{\sigma}{\lambda_{r}^{\star}}\sqrt{\frac{\mu r^{2}\log n}{n}}\right\} \\
 & \quad+\frac{\mu r}{n}\left\{ \frac{\kappa\sigma^{2}n}{\lambda_{r}^{\star}}+\sigma\sqrt{r\log n}\right\} \\
 & \lesssim\frac{\sigma^{2}\mu r\kappa^{2}\log n}{\lambda_{r}^{\star}}+\frac{\sigma\kappa\mu\sqrt{r^{3}\log n}}{n},
\end{align*}
%
where the second inequality follows from (\ref{eq:Psi-two-inf-norm-bound}),
(\ref{eq:Z-two-inf-norm-bound}), \eqref{eq:incoherence-U-estimate}, Theorem~\ref{thm:UsgnH-Ustar-MUstar-general}, and Lemma~\ref{lem:bound-Delta1-Delta2}. 

%






\subsection{Proof of Lemma~\ref{lem:normality-spectral}}
\label{sec:proof-lem:normality-spectral}

In what follows, we shall only focus on the case with $i\neq j$. The case with $i=j$ can be analyzed in an analogous manner; we omit it for the sake of brevity.
Before proceeding to the proof, we make note of a couple of basic facts about $\bm{P}^{\star}$ that will prove useful.  
The first property asserts that, for any $1\leq j\leq n$, 
%
\begin{align}
\sum_{l=1}^{n}P_{j,l}^{\star2} & =\sum_{l=1}^{n}P_{l,j}^{\star2}=\big\|\big(\bm{U}^{\star}\bm{U}^{\star\top}\big)_{\cdot,j}\big\|_{2}^{2}=\big\|\bm{U}^{\star}\big(\bm{U}_{j,\cdot}^{\star}\big)^{\top}\big\|_{2}^{2}=\bm{U}_{j,\cdot}^{\star}\bm{U}^{\star\top}\bm{U}^{\star}\big(\bm{U}_{j,\cdot}^{\star}\big)^{\top}\nonumber \\
 & =\big\|\bm{U}_{j,\cdot}^{\star}\big\|_{2}^{2}=P_{j,j}^{\star}.
	\label{eq:P-jl-square-sum}
\end{align}
%
The second property is concerned with the term $\|\bm{P}^{\star}\|_{\infty}$:  
%
\begin{equation}
	\|\bm{P}^{\star}\|_{\infty}=\|\bm{U}^{\star}\bm{U}^{\star\top}\|_{\infty}\leq\|\bm{U}^{\star}\|_{2,\infty}^{2}\leq\frac{\mu r}{n},
	\label{eq:P-inf-norm-UB}
\end{equation}
%
where the last inequality follows from the incoherence assumption. 




In view of the definition \eqref{eq:defn-P-UU-T} of $\bm{P}^{\star}$, we can express $\bm{W}=\bm{E}\bm{P}^{\star}+\bm{P}^{\star}\bm{E}$, which reveals that 
%
\begin{equation}
	W_{i,j}
	%=\sum_{l=1}^{n}E_{i,l}P_{l,j}+\sum_{l=1}^{n}P_{i,l}E_{l,j}
	=\sum_{l:\,l\neq j}E_{i,l}P^{\star}_{l,j}+\sum_{l:\,l\neq i}P^{\star}_{i,l}E_{l,j}+E_{i,j}(P^{\star}_{i,i}+P^{\star}_{j,j}).
\end{equation}
%
In other words,  $W_{i,j}$ can be viewed as a weighted sum of independent random variables $\{E_{i,l}\mid l\neq j\} \cup \{E_{l,j}\mid l\neq i\} \cup \{E_{i,j}\}$. 
To pin down the distribution of $W_{i,j}$, we resort to a non-asymptotic
version of the celebrated Berry-Esseen Theorem; 
see \citet[Theorem 3.7]{chen2010normal} for a proof using Stein's method.
%
\begin{theorem}[The Berry-Esseen bound]\label{thm:berry-esseen}
Let $\xi_{1},\ldots,\xi_{n}$ be independent zero-mean random variables satisfying $\sum_{i=1}^{n}\mathsf{Var}(\xi_{i})= v$.
Then the quantity $S= \frac{1}{\sqrt{v}} \sum_{i=1}^{n}\xi_{i}$ satisfies 
%
\[
	\sup_{z\in\mathbb{R}}\big|\,\mathbb{P}\left(S\leq z\right)-\Phi\left(z\right)\big|\leq10\gamma,\qquad\text{where}\ \gamma
	=\sum_{i=1}^{n} \frac{\mathbb{E}\big[|\xi_{i}|^{3}\big] }{ v^{3/2} }.
\]
%
%where $\Phi(\cdot)$ represents the CDF of a standard Gaussian distribution. 
\end{theorem}





According to the Berry-Esseen bound (cf.~Theorem~\ref{thm:berry-esseen}), 
proving the approximate Gaussianity of $W_{i,j}$ boils down to characterizing the second and the third moments of these random variables under consideration. 




Let us start with the variance statistics. 
Given that $\{E_{i,j}\mid i\geq j\}$ are independently generated, 
%the variance  $v_{i,j}^{\star}  \coloneqq\mathsf{Var}(W_{i,j}) $ 
we can straightforwardly see that
%
\begin{align}
	v_{i,j}^{\star}   %\coloneqq\mathsf{Var}(W_{i,j})
	&=\sum_{l:\,l\neq j}\sigma_{i,l}^{2}P_{l,j}^{\star 2}+\sum_{l:\,l\neq i}P_{i,l}^{\star 2}\sigma_{l,j}^{2}+\sigma_{i,j}^{2}(P_{i,i}^{\star}+P_{j,j}^{\star})^{2} \notag\\
	&= \mathsf{Var}(W_{i,j}).
	\label{eq:variance-formula-Mij}
\end{align}
%
We now develop a lower bound on this variance term.  Given
that $\bm{P}^{\star} \succeq \bm{0}$, one has $P_{i,i}^{\star},P_{j,j}^{\star}\geq0$,
%and $(P^{\star}_{i,i}+P^{\star}_{j,j})^{2}\geq P_{i,i}^{\star 2}+P_{j,j}^{\star 2}$,
which combined with \eqref{eq:P-jl-square-sum} reveals that
%
\begin{align}
v_{i,j}^{\star} & 
%\geq\sigma_{\min}^{2}\Bigg\{\sum_{l:\,l\neq i}P_{l,j}^{\star 2}+\sum_{l:\,l\neq j}P_{i,l}^{\star 2}+P_{i,i}^{\star 2}+P_{j,j}^{\star 2}\Bigg\}
\geq \sigma_{\min}^{2}\Bigg\{\sum_{l=1}^{n}P_{l,j}^{\star 2}+\sum_{l=1}^{n}P_{i,l}^{\star 2}\Bigg\} 
%\notag\\
% & =\sigma_{\min}^{2}\Big\{\big\|\big(\bm{U}^{\star}\bm{U}^{\star\top}\big)_{\cdot,j}\big\|_{2}^{2}+\big\|\big(\bm{U}^{\star}\bm{U}^{\star\top}\big)_{i,\cdot}\big\|_{2}^{2}\Big\} \notag\\
% & =\sigma_{\min}^{2}\Big\{\big\|\bm{U}^{\star}\big(\bm{U}_{j,\cdot}^{\star}\big)^{\top}\big\|_{2}^{2}+\big\|\bm{U}_{i,\cdot}^{\star}\bm{U}^{\star\top}\big\|_{2}^{2}\big\} \notag\\
% & 
=\sigma_{\min}^{2}\Big\{\big\|\bm{U}_{j,\cdot}^{\star}\big\|_{2}^{2}+\big\|\bm{U}_{i,\cdot}^{\star}\big\|_{2}^{2}\Big\}.
	\label{eq:vij-lower-bound-12}
\end{align}
%

Next, we move on to bound the third moments. 
Utilizing the independence of $\{E_{i,j}\mid i\geq j\}$ once again gives
%
\begin{align}
	\gamma & \coloneqq\frac{\sum_{l:\,l\neq j}\mathbb{E}\big[\big|E_{i,l}\big|^{3}\big]\big|P_{l,j}^{\star} \big|^{3}+\sum_{l:\,l\neq i}\big|P_{i,l}\big|^{3}\mathbb{E}\big[\big|E_{l,j}\big|^{3}\big]+\mathbb{E}\big[\big|E_{i,j}\big|^{3}\big]\big|P_{i,i}^{\star} +P_{j,j}^{\star} \big|^{3}}{\big(v_{i,j}^{\star}\big)^{3/2}} \notag\\
 & \leq\frac{2B\|\bm{P}^{\star} \|_{\infty}}{\big(v_{i,j}^{\star}\big)^{3/2}}\left\{ \sum_{l:\,l\neq j}\mathbb{E}\big[E_{i,l}^{2}\big] P_{l,j}^{\star 2}+\sum_{l:\,l\neq i} P_{i,l}^{\star 2}\mathbb{E}\big[E_{l,j}^{2}\big]+\mathbb{E}\big[E_{i,j}^{2}\big]\big|P_{i,i}^{\star} +P_{j,j}^{\star} \big|^{2}\right\} \notag\\
 & =\frac{2B\|\bm{P}^{\star} \|_{\infty}}{\big(v_{i,j}^{\star}\big)^{3/2}}\cdot v_{i,j}^{\star}=\frac{2B\|\bm{P}^{\star} \|_{\infty}}{\big(v_{i,j}^{\star}\big)^{1/2}}, 
	\label{eq:gamma-UB-123}
\end{align}
% 
where we have used the assumption that $|E_{i,j}|\leq B$, and the last line arises from the expression \eqref{eq:variance-formula-Mij}. 
Substituting \eqref{eq:P-inf-norm-UB} and \eqref{eq:vij-lower-bound-12} into \eqref{eq:gamma-UB-123} and invoking the elementary identity $\|\bm{U}^{\star} \|_{\mathrm{F}}^2=r$, 
we arrive at
%
\begin{align*}
\gamma & \leq\frac{2B\mu r}{n\sigma_{\min}\sqrt{\big\|\bm{U}_{j,\cdot}^{\star}\big\|_{2}^{2}+\big\|\bm{U}_{i,\cdot}^{\star}\big\|_{2}^{2}}}=\frac{2B\mu\sqrt{r}\,\big\|\bm{U}^{\star}\big\|_{\mathrm{F}}}{n\sigma_{\min}\sqrt{\big\|\bm{U}_{j,\cdot}^{\star}\big\|_{2}^{2}+\big\|\bm{U}_{i,\cdot}^{\star}\big\|_{2}^{2}}}
	= o(1),
\end{align*}
%
provided that Condition~\eqref{eq:Uj-Ui-lower-bound-lem-2} holds. 


With the above calculations in place, invoking Theorem~\ref{thm:berry-esseen} immediately leads to 
%
\[
\sup_{z\in\mathbb{R}}\left|\mathbb{P}\left( W_{i,j}\leq\sqrt{v_{i,j}^{\star}}z\right) -\Phi(z)\right|\leq10\gamma=o(1). 
\]
%
as claimed in \eqref{eq:eqn-normality-spectral}. 



Finally, we turn to proving the bound \eqref{eq:Phi-inf-norm-o-vij}. 
By virtue of Lemma~\ref{thm:matrix-distribution} and \eqref{eq:vij-lower-bound-12}, we know that 
%
\begin{align*}
\big\|\bm{\Phi}\big\|_{\infty} & \lesssim\frac{\sigma^{2}\mu r\kappa^{2}\log n}{\lambda_{r}^{\star}}+\frac{\sigma\kappa\mu\sqrt{r^{3}\log n}}{n}\\
 & =o\left(\sigma_{\min}\sqrt{\big\|\bm{U}_{j,\cdot}^{\star}\big\|_{2}^{2}+\big\|\bm{U}_{i,\cdot}^{\star}\big\|_{2}^{2}}\right)\leq o\big(\sqrt{v_{i,j}^{\star}}\big) , 
\end{align*}
%
with the proviso that
%
\begin{align}
	& \sqrt{\big\|\bm{U}_{j,\cdot}^{\star}\big\|_{2}^{2}+\big\|\bm{U}_{i,\cdot}^{\star}\big\|_{2}^{2}}  \gtrsim\frac{\sigma^{2}\mu r\kappa^{2}\log^{3/2}n}{\sigma_{\min}\lambda_{r}^{\star}}+\frac{\sigma\kappa\mu\sqrt{r^{3}}\log n}{\sigma_{\min}n} \notag\\
 & \qquad =\left(\frac{\sigma^{2}\mu\sqrt{r}\kappa^{2}\log^{3/2}n}{\sigma_{\min}\lambda_{r}^{\star}}+\frac{\sigma\kappa\mu r\log n}{\sigma_{\min}n}\right)\big\|\bm{U}^{\star}\big\|_{\mathrm{F}}. 
	\label{eq:condition-7878}
\end{align}
%
Recognizing the trivial bound $\sigma^2 = \max_{i,j}\mathbb{E}[E_{i,j}^2]\leq B^2$, we know that Condition~\eqref{eq:condition-7878} holds as long as
%
\begin{align}
	\frac{\big\|\bm{U}_{j,\cdot}^{\star}\big\|_{2}^{2}+\big\|\bm{U}_{i,\cdot}^{\star}\big\|_{2}^{2}}{\big\|\bm{U}^{\star}\big\|_{\mathrm{F}}^{2}}
	& \gtrsim \frac{\sigma^{4}\mu^{2}r\kappa^{4}\log^{3}n}{\sigma_{\min}^{2}(\lambda_{r}^{\star})^{2}} + \frac{B^{2}\kappa^2 \mu^{2}r^2\log ^2 n}{\sigma_{\min}^{2}n^{2}}, 
	%+\frac{\sigma^{2}\kappa^{2}\mu^{2}r^{2}\log^{2}n}{\sigma_{\min}^{2}n^{2}}.
	%\label{eq:Uj-Ui-lower-bound-lem-2}
\end{align}
%
which is precisely Condition~\eqref{eq:Uj-Ui-lower-bound-lem-2}. This concludes the proof of Lemma~\ref{lem:normality-spectral}. 

%as claimed. 


%
%\[
%\frac{B\|\bm{P}\|_{\infty}}{\big(v_{i,j}^{\star}\big)^{1/2}}\lesssim\frac{\sigma_{\max}\sqrt{n/(\mu^{2}\log n)}\cdot\frac{\mu r}{n}}{\sigma_{\min}\Big\{\Big\|\bm{U}_{j,\cdot}^{\star}\Big\|_{2}+\Big\|\bm{U}_{i,\cdot}^{\star}\Big\|_{2}\Big\}}\asymp\frac{1}{\sqrt{\log n}}\frac{\|\bm{U}^{\star}\|_{\mathrm{F}}}{\sqrt{n}\Big\{\big\|\bm{U}_{j,\cdot}^{\star}\big\|_{2}+\big\|\bm{U}_{i,\cdot}^{\star}\big\|_{2}\Big\}}
%\]
%


%And we know that 
%\[
%B\lesssim\sigma_{\max}\sqrt{n/(\mu^{2}\log n)}.
%\]






\subsection{Proof of Lemma~\ref{lem:bound-Delta1-Delta2}}
\label{sec:proof-auxiliary-distribution}



%\paragraph{Proof of Lemma~\ref{lem:bound-Delta1-Delta2}.}
To begin with, let us begin by proving \eqref{eq:Delta-1-U-R-Lambda}. 
From the definition of the quantity $\mathcal{E}_{1}$ (see Lemma~\ref{lem:UH-MUstar-diff-decompose}), we have
%
\begin{align*}
 & \big\|\bm{M}(\bm{U}\bm{H}-\bm{U}^{\star})\big\|_{2,\infty}=\lambda_{r}^{\star}\mathcal{E}_{1}/2\\
 & \qquad\lesssim\left(\alpha_{0}+\alpha_{1}+\alpha_{2}\right)+\left(\sigma\sqrt{n}+B\log n\right)\big\|\bm{U}\bm{H}-\bm{U}^{\star}\big\|_{2,\infty}\\
 & \qquad\lesssim\left(\alpha_{0}+\alpha_{1}+\alpha_{2}\right)+\left(\sigma\sqrt{n}+B\log n\right)\frac{\sigma\kappa\sqrt{\mu r\log n}}{\lambda_{r}^{\star}}\\
 & \qquad\asymp\frac{\sigma^{2}\kappa\sqrt{\mu rn\log^{2}n}}{\lambda_{r}^{\star}},
\end{align*}
%
where the first inequality comes from \eqref{eq:E1-UB-general-1} and \eqref{eq:defn-E1-rho1},
the second inequality is a consequence of Theorem~\ref{thm:UsgnH-Ustar-MUstar-general}, and 
the last line relies on our previous bounds on $\alpha_0,\alpha_1,\alpha_2$ (see \eqref{eq:defn-alpha-0-general}, \eqref{eq:alpha1-bound-1234} and \eqref{eq:Mstar-UH-Ustar-2inf}) and holds as long as $B\lesssim\sigma\sqrt{n/(\mu\log n)}$. 
%
Additionally, from the elementary identity $\bm{M}\bm{U}\bm{H}=\bm{U}\bm{\Lambda}\bm{H}$, we obtain
%
\begin{align*}
\big\|\bm{U}\bm{\Lambda}\mathsf{sgn}(\bm{H})-\bm{M}\bm{U}\bm{H}\big\|_{2,\infty} & =\big\|\bm{U}\bm{\Lambda}\mathsf{sgn}(\bm{H})-\bm{U}\bm{\Lambda}\bm{H}\big\|_{2,\infty}\\
 & \leq\big\|\bm{U}\big\|_{2,\infty}\big\|\bm{\Lambda}\big\|\,\big\|\mathsf{sgn}(\bm{H})-\bm{H}\big\|\\
 & \lesssim\sqrt{\frac{\mu r}{n}} |\lambda_{1}^{\star} | \cdot\frac{\sigma^{2}n}{(\lambda_{r}^{\star})^{2}}=\frac{\kappa\sigma^{2}\sqrt{\mu rn}}{\lambda_{r}^{\star}},
\end{align*}
%
where the last line results from Lemma~\ref{lem:H-property-summary-general}, the fact \eqref{eq:incoherence-U-estimate}, and the following inequality
%
\begin{equation}
	\|\bm{\Lambda}\| \leq \|\bm{\Lambda}^{\star}\| + \|\bm{E}\|\leq |\lambda_1^{\star}| + O(\sigma\sqrt{n}) \leq 2|\lambda_1^{\star}|. 
	\label{eq:Lambda-norm-upper-bound}
\end{equation}
%
Taking together the above bounds and applying the triangle inequality 
immediately establish~\eqref{eq:Delta-1-U-R-Lambda}. 







Next, we turn to the proof of the bound \eqref{eq:sgn-H-Lambdastar-sgn-H-Lambda}. 
Note that it has been shown in (\ref{eq:H-Lambda-star-H-Lambda-dist}) that
%
\[
\big\|\bm{H}\bm{\Lambda}^{\star}\bm{H}^{\top}-\bm{\Lambda}\big\|\lesssim\frac{\kappa\sigma^{2}n}{\lambda_{r}^{\star}}+\sigma\sqrt{r\log n}.
\]
%
In addition, the triangle inequality leads to
%
\begin{align*}
 & \big\|\mathsf{sgn}(\bm{H})\bm{\Lambda}^{\star}\mathsf{sgn}(\bm{H})^{\top}-\bm{H}\bm{\Lambda}^{\star}\bm{H}^{\top}\big\|\\
 & \qquad\leq\big\|\mathsf{sgn}(\bm{H})\bm{\Lambda}^{\star}\big(\mathsf{sgn}(\bm{H})-\bm{H}\big)^{\top}\big\|+\big\|\big(\mathsf{sgn}(\bm{H})-\bm{H}\big)\bm{\Lambda}^{\star}\bm{H}^{\top}\big\|\\
 & \qquad\leq\left(\big\|\mathsf{sgn}(\bm{H})\big\|+\big\|\bm{H}\big\|\right)\,\big\|\bm{\Lambda}^{\star}\big\|\,\big\|\mathsf{sgn}(\bm{H})-\bm{H}\big\|\\
 & \qquad\leq  2 |\lambda_{1}^{\star}| \, \big\|\mathsf{sgn}(\bm{H})-\bm{H}\big\|\\
 & \qquad\lesssim |\lambda_{1}^{\star}| \frac{\sigma^{2}n}{\big(\lambda_{r}^{\star}\big)^{2}}=\frac{\kappa\sigma^{2}n}{\lambda_{r}^{\star}},
\end{align*}
%
where the third line follows since $\big\|\mathsf{sgn}(\bm{H})\big\|=1$
and $\|\bm{H}\|\leq\|\bm{U}\|\|\bm{U}^{\star}\|=1$, and the last
inequality comes from (\ref{eq:H-sgnH-diff-UB-123}). Combining
the above two results and invoking the triangle inequality lead to
the advertised bound \eqref{eq:sgn-H-Lambdastar-sgn-H-Lambda}. 





\section{Appendix E: Proof of Theorem~\ref{thm:CI-general}} 
\label{sec:proof-thm:CI-general}


With the distributional guarantees in Theorem~\ref{thm:matrix-distribution-complete} in place, 
the only remaining task boils down to verifying the statistical accuracy of the variance estimator $\widehat{v}_{i,j}$. 
This can be achieved via the following lemma, whose proof is provided in Section~\ref{sec:proof-lem:variance-estimation-guarantee}.  
%
\begin{lemma}
	\label{lem:variance-estimation-guarantee}
	%
	Suppose that the assumptions of Theorem~\ref{thm:UsgnH-Ustar-MUstar-general} hold. In addition, assume that $\kappa^{4}\mu^{2}r^{2}\log n\leq n$, $\sigma\sqrt{n}\lesssim|\lambda_{r}^{\star}|/\kappa$ and
	%
	\begin{align}
		\big\|\bm{U}_{j,\cdot}^{\star}\big\|_{2}^{2}+\big\|\bm{U}_{i,\cdot}^{\star}\big\|_{2}^{2}\gtrsim\frac{B\sigma\kappa^{2}\mu^{2}r^{2}\log n}{\sigma_{\min}^2 n^{3/2}}+\frac{\sigma^{3}\kappa\mu^{2}r^{2}\log n}{\sigma_{\min}^2|\lambda_{r}^{\star}|\sqrt{n}}. 
		\label{eq:Uj-Ui-lower-bound-lem}
	\end{align}
	%
	With probability exceeding $1-O(n^{-5})$, one has
	%
	\begin{align}
		\big|\widehat{v}_{i,j}-v_{i,j}^{\star}\big| =o\big( v_{i,j}^{\star} \big). 
	\end{align}
	%
\end{lemma}
%

Lemma~\ref{lem:variance-estimation-guarantee} essentially enables us to express
%
\[
	\widehat{v}_{i,j} = (1 + \zeta_v)^2 v_{i,j}^{\star} \qquad \text{with } \zeta_v = o(1). 
\]
%
As a consequence, we can further demonstrate that
%
\begin{align*}
 & \left|\mathbb{P}\left( M_{i,j}^{\star}\in\mathsf{CI}_{i,j}^{1-\alpha}\right) -(1-\alpha)\right|=\left|\mathbb{P}\left( \big| \widehat{M}_{i,j}-M_{i,j}^{\star}\big|\leq z_{\alpha/2}\sqrt{\widehat{v}_{i,j}}\right) -(1-\alpha)\right|\\
 & =\left|\mathbb{P}\left( \big| \widehat{M}_{i,j}-M_{i,j}^{\star}\big|\leq(1+\zeta_{v})z_{\alpha/2}\sqrt{v_{i,j}^{\star}}\right) -(1-\alpha)\right|\\
 & \leq\left|\Phi\big((1+\zeta_{v})z_{\alpha/2}\big)-\Phi\big(-(1+\zeta_{v})z_{\alpha/2}\big)-(1-\alpha)\right|+o(1)\\
 & \leq\left|\Phi\big(z_{\alpha/2}\big)-\Phi\big(-z_{\alpha/2}\big)-(1-\alpha)\right|+2\left|\Phi\big((1+\zeta_{v})z_{\alpha/2}\big)-\Phi\big(z_{\alpha/2}\big)\right|+o(1)\\
 & \leq2\zeta_{v}z_{\alpha/2}=o(1),
\end{align*}
%
where the first inequality follows from Theorem~\ref{thm:matrix-distribution-complete} and $z_{\alpha/2}:= \Phi^{-1}(1-\alpha/2)$, the second inequality applies the triangle inequality, 
and the validity of the last line can be seen from the basic fact $|\Phi(u)-\Phi(v)| \leq |u-v|$.   


When $\sigma_{\min}\asymp \sigma$, Condition~\eqref{eq:Uj-Ui-lower-bound-lem} simplifies to
%
\begin{align}
	\frac{\big\|\bm{U}_{j,\cdot}^{\star}\big\|_{2}^{2}+\big\|\bm{U}_{i,\cdot}^{\star}\big\|_{2}^{2}}{\big\|\bm{U}^{\star}\big\|_{\mathrm{F}}^{2}}
	\gtrsim\frac{B\kappa^{2}\mu^{2}r\log n}{\sigma n^{3/2}}+\frac{\sigma\kappa\mu^{2}r\log n}{|\lambda_{r}^{\star}|\sqrt{n}}.		 
		\label{eq:Uj-Ui-lower-bound-lem-31}
\end{align}
%
We still need to ensure that Condition~\eqref{eq:Ui-Uj-condition-thm-distribution-1} is satisfied.  
It is seen that 
%
\begin{align*}
	\frac{B^{2}\kappa^{2}\mu^{2}r^{2}\log^{2}n}{\sigma^{2}n^{2}}+\frac{\sigma^{2}\mu^{2}r\kappa^{4}\log^{3}n}{(\lambda_{r}^{\star})^{2}}&\lesssim\frac{B\kappa^{2}\mu^{2}r^{2}\log^{2}n}{\sigma n^{3/2}}+\frac{\sigma\mu^{2}r\kappa^{3}\log^{3}n}{|\lambda_{r}^{\star}|\sqrt{n}},
\end{align*}
%
%where the first inequality has made use of the fact that $B\geq \sigma$, and the second inequality 
which holds if $\sigma\kappa \sqrt{n} \lesssim |\lambda_r^{\star}|$ and $B\lesssim \sigma\sqrt{n}$. 
As a result, if Condition~\eqref{eq:Uj-Ui-lower-bound-thm} holds, then both \eqref{eq:Uj-Ui-lower-bound-lem-31} and \eqref{eq:Ui-Uj-condition-thm-distribution-1}
are satisfied. 
This finishes the proof, as long as  Lemma~\ref{lem:variance-estimation-guarantee} can be established. 



%%
%\begin{align*}
% & \frac{B^{2}\mu^{2}r\log n}{n^{2}\sigma^{2}}+\frac{\sigma^{2}\mu^{2}r\kappa^{4}\log^{3}n}{(\lambda_{r}^{\star})^{2}}+\frac{\kappa^{2}\mu^{2}r^{2}\log^{2}n}{n^{2}}\\
% & \lesssim\frac{B^{2}\kappa^{2}\mu^{2}r^{2}\log^{2}n}{n^{2}\sigma^{2}}+\frac{\sigma^{2}\mu^{2}r\kappa^{4}\log^{3}n}{(\lambda_{r}^{\star})^{2}}\lesssim\frac{B\kappa^{2}\mu^{2}r^{2}\log^{2}n}{\sigma n^{3/2}}+\frac{\sigma\mu^{2}r\kappa^{3}\log^{3}n}{|\lambda_{r}^{\star}|\sqrt{n}},
%\end{align*}
%%


\subsection{Proof for Lemma~\ref{lem:variance-estimation-guarantee}}
\label{sec:proof-lem:variance-estimation-guarantee}

As before, we shall only present the proof for the case with $i\neq j$ for the sake of conciseness. 
In order to justify the goodness of the estimator $\widehat{v}_{i,j}$, 
we find it convenient to first look at the surrogate estimator introduced in \eqref{eq:v-ij-plugin-surrogate-123}, i.e., 
%
\begin{equation}
	\widetilde{v}_{i,j}=\sum_{l=1}^{n}E_{i,l}^{2}P_{l,j}^{\star 2}+\sum_{l=1}^{n}P_{i,l}^{\star 2}E_{l,j}^{2}+2E_{i,j}^{2}P_{i,i}^{\star}P_{j,j}^{\star}, 
	\label{eq:v-surrogate}
\end{equation}
%
% which can also be viewed as a plug-in estimator if an oracle reveals all information about $\bm{E}$ and $\bm{P}^{\star}$. 
In the sequel, our proof consists of two main steps: 
%
\begin{itemize}
	\item Show that the surrogate $\widetilde{v}_{i,j}$ is a reliable estimate of the truth ${v}^{\star}_{i,j}$, namely, $\widetilde{v}_{i,j}\approx {v}^{\star}_{i,j}$. 
	\item Show that the estimator in use and the surrogate estimator are sufficiently close, namely, $\widehat{v}_{i,j}\approx \widetilde{v}_{i,j}$.
\end{itemize}


\subsubsection{Step 1: show that $\widetilde{v}_{i,j}\approx {v}^{\star}_{i,j}$} 
%
Firstly, the fact that the $E_{i,j}$'s are zero-mean random variables
immediately reveals that $\widetilde{v}_{i,j}$ is an unbiased estimate
of $v_{i,j}^{\star}$, that is, 
%
\[
\mathbb{E}\big[\widetilde{v}_{i,j}\big]=v_{i,j}^{\star}.
\]
%
Secondly, given that the $E_{i,j}$'s are statistically independent,
we intend to invoke the Bernstein inequality to control the difference $\widetilde{v}_{i,j}-v_{i,j}^{\star}=\widetilde{v}_{i,j}-\mathbb{E}\big[\widetilde{v}_{i,j}\big]$.
To do so, one first calculates that
%
\begin{align*}
L_{0} & \coloneqq\max\left\{ \max_{l}E_{i,l}^{2}P_{l,j}^{\star 2},\,\max_{l}E_{l,j}^{2}P_{i,l}^{\star 2},\,E_{i,j}^{2}(P_{i,i}^\star +P_{j,j}^\star )^{2}\right\} \leq4B^{2}\|\bm{P}^\star \|_{\infty}^{2}
\end{align*}
%
and
\begin{align*}
V_{0} & \coloneqq\sum_{l:\,l\neq j}\mathsf{Var}\big(E_{i,l}^{2}\big)P_{l,j}^{\star 4}+\sum_{l:\,l\neq i}\mathsf{Var}\big(E_{l,j}^{2}\big)P_{i,l}^{\star 4}+\mathsf{Var}\big(E_{i,j}^{2}\big)(P_{i,i}^\star +P_{j,j}^\star )^{4}\\
 & \lesssim\sum_{l=1}^{n}\mathbb{E}\big[E_{i,l}^{4}\big]P_{l,j}^{\star 4}+\sum_{l=1}^{n}\mathbb{E}\big[E_{l,j}^{4}\big]P_{i,l}^{\star 4}+\mathbb{E}\big[E_{i,j}^{4}\big]P_{i,i}^{\star 2}P_{j,j}^{\star 2}\\
 & \lesssim B^{2}\|\bm{P}^\star \|_{\infty}^{2}\left\{ \sum_{l=1}^{n}\mathbb{E}\big[E_{i,l}^{2}\big]P_{l,j}^{2}+\sum_{l=1}^{n}\mathbb{E}\big[E_{l,j}^{2}\big]P_{i,l}^{\star 2}+\mathbb{E}\big[E_{i,j}^{2}\big]P_{i,i}^\star P_{j,j}^\star \right\} \\
 & \lesssim \sigma^{2} B^{2} \|\bm{P}^\star \|_{\infty}^{2}\left\{ \sum_{l=1}^{n}P_{l,j}^{\star 2}+\sum_{l=1}^{n}P_{i,l}^{\star 2}\right\} \\
 & \lesssim \sigma^{2} B^{2} \|\bm{P}^\star \|_{\infty}^{3},
\end{align*}
%
where the last line follows from \eqref{eq:P-jl-square-sum}. 
%since, for any $1\leq j\leq n$, 
%%
%\[
%\sum_{l=1}^{n}P_{l,j}^{2}=\big\|\bm{U}^{\star}\big(\bm{U}_{j,\cdot}^{\star}\big)^{\top}\big\|_{2}^{2}=\big\|\bm{U}_{j,\cdot}^{\star}\big\|_{2}^{2}=P_{j,j}\leq\|\bm{P}\|_{\infty}.
%\]
%
Invoking the Bernstein inequality (cf.~Corollary~\ref{thm:matrix-Bernstein-friendly}) reveals that with probability exceeding $1-O(n^{-5})$, 
%
\begin{align}
\big|\widetilde{v}_{i,j}-v_{i,j}^{\star}\big| & =\left|\widetilde{v}_{i,j}-\mathbb{E}\left[\widetilde{v}_{i,j}\right]\right|\lesssim\sqrt{V_{0}\log n}+L_{0}\log n \notag\\
 & \lesssim\sigma B\sqrt{\|\bm{P}^\star \|_{\infty}^{3}\log n}+B^{2}\|\bm{P}^\star \|_{\infty}^{2}\log n \notag\\
 & \lesssim\frac{\sigma B\mu^{3/2}r^{3/2}\sqrt{\log n}}{n^{3/2}}+\frac{\mu^{2}r^{2}B^{2}\log n}{n^{2}} ,
	\label{eq:tilde-v-star-v-diff}
\end{align}
%
where the last inequality results from \eqref{eq:P-inf-norm-UB}. 



\subsubsection{Step 2: show that $\widehat{v}_{i,j}\approx \widetilde{v}_{i,j}$} 
%
In order to accomplish this, we are in need of controlling the difference between $E_{i,j}$ (resp.~$P^\star_{i,j}$) and $\widehat{E}_{i,j}$ (resp.~$\widehat{P}_{i,j}$). 
To this end, apply the entrywise estimation guarantees in Corollary~\ref{cor:entrywise-error-general} to yield
%
\begin{align}
\big\|\widehat{\bm{E}}-\bm{E}\big\|_{\infty} & =\big\|\big(\bm{M}-\bm{U}\bm{\Lambda}\bm{U}^{\top}\big)-\big(\bm{M}-\bm{M}^{\star}\big)\big\|_{\infty} \notag\\
 & =\big\|\bm{U}\bm{\Lambda}\bm{U}^{\top}-\bm{M}^{\star}\big\|_{\infty}\lesssim\sigma\kappa^{2}\mu r\sqrt{\frac{\log n}{n}},
	\label{eq:entrywise-estimation-E}
\end{align}
%
and as a result,
%
\begin{equation}
\big\|\widehat{\bm{E}}\big\|_{\infty}\leq\big\|\bm{E}\big\|_{\infty}+\big\|\widehat{\bm{E}}-\bm{E}\big\|_{\infty}
	\lesssim B+\sigma\kappa^{2}\mu r\sqrt{\frac{\log n}{n}}
	\asymp B,
	\label{eq:entrywise-estimation-E-1}
\end{equation}
%
provided that $B\gtrsim \sigma\kappa^{2}\mu r\sqrt{\frac{\log n}{n}}$ (which is trivially satisfied if $\kappa^{2}\mu r\sqrt{\frac{\log n}{n}}\leq 1$). Moving on to the error term $\widehat{P}_{i,j}-P_{i,j}^\star $, we observe that
%
\begin{align*}
 & \big\|\widehat{\bm{P}}-\bm{P}^\star \big\|_{\infty} =\big\|\bm{U}\bm{U}^{\top}-\bm{U}^{\star}\bm{U}^{\star\top}\big\|_{\infty}\\
 & \qquad \leq\big\|\big(\bm{U}\mathsf{sgn}(\bm{H})-\bm{U}^{\star}\big)\bm{U}^{\star\top}\big\|_{\infty}+\big\|\bm{U}\mathsf{sgn}(\bm{H})\big(\bm{U}\mathsf{sgn}(\bm{H})-\bm{U}^{\star}\big)^{\top}\big\|_{\infty}\\
 & \qquad \leq\big\|\bm{U}\mathsf{sgn}(\bm{H})-\bm{U}^{\star}\big\|_{2,\infty}\left\{ \big\|\bm{U}^{\star}\big\|_{2,\infty}+\big\|\bm{U}\mathsf{sgn}(\bm{H})\big\|_{2,\infty}\right\} \\	
 & \qquad \leq\big\|\bm{U}\mathsf{sgn}(\bm{H})-\bm{U}^{\star}\big\|_{2,\infty}\left\{ \big\|\bm{U}^{\star}\big\|_{2,\infty}+\big\|\bm{U}\big\|_{2,\infty}\right\} .
\end{align*}
%
Taking this  together with the bound \eqref{eq:UsgnH-Ustar-bound-theorem-general} in Theorem~\ref{thm:UsgnH-Ustar-MUstar-general}, 
the incoherence assumption, and the inequality \eqref{eq:incoherence-U-estimate}, we arrive at
%
\begin{align}
	\big\|\widehat{\bm{P}}-\bm{P}^\star \big\|_{\infty} & \lesssim \frac{(\sigma\kappa\sqrt{\mu r}+\sigma\sqrt{r\log n})}{|\lambda_{r}^{\star}|} \sqrt{\frac{\mu r}{n}}. \label{eq:P-hat-Linf-norm-diff}
\end{align}
%
This taken together with \eqref{eq:P-inf-norm-UB} indicates that
%
\begin{align}
\big\|\widehat{\bm{P}}\big\|_{\infty} & \leq\big\|\bm{P}^\star \big\|_{\infty}+\big\|\widehat{\bm{P}}-\bm{P}^\star \big\|_{\infty}\nonumber \\
 & \lesssim \frac{\mu r}{n}+\frac{(\sigma\kappa\sqrt{\mu r}+\sigma\sqrt{r\log n})}{|\lambda_{r}^{\star}|}\sqrt{\frac{\mu r}{n}}\asymp\frac{\mu r}{n},\label{eq:P-hat-Linf-norm}
\end{align}
%
with the proviso that $\sigma\sqrt{n}\lesssim|\lambda_{r}^{\star}|/\kappa$ and $\sigma\sqrt{n\log n}\lesssim|\lambda_{r}^{\star}|$. 



Armed with the preceding bounds, we are now positioned to control $\widehat{v}_{i,j} - \widetilde{v}_{i,j}$. From the definition of $\widetilde{v}_{i,j}$ and $v_{i,j}^{\star}$, we recognize that
%
\begin{align}
\big| \widehat{v}_{i,j} - \widetilde{v}_{i,j} \big| & \leq\underset{\eqqcolon\,\alpha_{1}}{\underbrace{\left|\sum_{l=1}^{n}\left(\widehat{E}_{i,l}^{2}\widehat{P}_{l,j}^{2}-E_{i,l}^{2}P_{l,j}^{\star 2}\right)\right|}}+\underset{\eqqcolon\,\alpha_{2}}{\underbrace{\left|\sum_{l=1}^{n}\left(\widehat{P}_{i,l}^{2}\widehat{E}_{l,j}^{2}-P_{i,l}^{\star 2}E_{l,j}^{2}\right)\right|}}\nonumber \\
 & \qquad+2\, \underset{\eqqcolon\,\alpha_{3}}{\underbrace{\left|\widehat{E}_{i,j}^{2}\widehat{P}_{i,i}\widehat{P}_{j,j}-E_{i,j}^{2}P_{i,i}^\star P_{j,j}^\star \right|}},
	\label{eq:v-tilde-bound-alpha-123}
\end{align}
%
leaving us with three terms to cope with. 
Regarding the first term $\alpha_1$ on the right-hand side of \eqref{eq:v-tilde-bound-alpha-123}, it can be easily verified that
%
\begin{align}
\alpha_{1} & \leq
\left|\sum_{l=1}^{n}\left(\widehat{E}_{i,l}^{2}\widehat{P}_{l,j}^{2}-E_{i,l}^{2}\widehat{P}_{l,j}^{2}\right)\right|+\left|\sum_{l=1}^{n}\left(E_{i,l}^{2}\widehat{P}_{l,j}^{2}-E_{i,l}^{2}P_{l,j}^{\star2}\right)\right| \notag\\
 & \leq \left(\max_{l}\left|\widehat{E}_{i,l}^{2}-E_{i,l}^{2}\right|\right)\sum_{l=1}^{n}\widehat{P}_{l,j}^{2}+\left(\max_{l}\left| \widehat{P}_{l,j}^{2} - P_{l,j}^{\star 2} \right|\right)\sum_{l=1}^{n}E_{i,l}^{2} \notag\\
	& \overset{\mathrm{(i)}}{\leq} \left(\|\bm{E}\|_{\infty}+\|\widehat{\bm{E}}\|_{\infty}\right)\big\|\widehat{\bm{E}}-\bm{E}\big\|_{\infty}\widehat{P}_{j,j} \notag\\
	&\qquad +\left(\|\bm{P}^\star \|_{\infty}+\|\widehat{\bm{P}}\|_{\infty}\right)\big\|\widehat{\bm{P}}-\bm{P}^\star \big\|_{\infty}\|\bm{E}\|^{2} \notag\\
	& \overset{\mathrm{(ii)}}{\lesssim} B\sigma\kappa^{2}\mu r\sqrt{\frac{\log n}{n}}\cdot\frac{\mu r}{n}+\frac{\mu r}{n}\cdot\frac{(\sigma\kappa\sqrt{\mu r}+\sigma\sqrt{r\log n})}{|\lambda_{r}^{\star}|}\sqrt{\frac{\mu r}{n}}\cdot\sigma^{2}n.  \notag
	%\label{eq:alpha1-ii-UB}
\end{align}
%
Here, (i) makes use of \eqref{eq:P-jl-square-sum} (with $\bm{P}^{\star}$ replaced by $\widehat{\bm{P}}$), 
whereas (ii) holds true due to \eqref{eq:entrywise-estimation-E}, \eqref{eq:entrywise-estimation-E-1}, \eqref{eq:P-inf-norm-UB}, \eqref{eq:P-hat-Linf-norm-diff}, \eqref{eq:P-hat-Linf-norm}, and Lemma~\ref{lemma:general-noise-bound}.  
The second term $\alpha_2$ on the right-hand side of \eqref{eq:v-tilde-bound-alpha-123} can be bounded in the same manner and we omit it here for brevity. 
When it comes to the last term $\alpha_3$ on the right-hand side of \eqref{eq:v-tilde-bound-alpha-123}, one has
%
\begin{align*}
\alpha_{3} & \leq\left|\widehat{E}_{i,j}^{2}-E_{i,j}^{2}\right|\left|\widehat{P}_{i,i}\widehat{P}_{j,j}\right|+ E_{i,j}^{2} \left|\widehat{P}_{i,i}\widehat{P}_{j,j}-P_{i,i}^\star P_{j,j}^\star \right|\\
 & \lesssim\big\|\widehat{\bm{E}}-\bm{E}\big\|_{\infty}\left\{ \big\|\bm{E}\big\|_{\infty}+\big\|\widehat{\bm{E}}\big\|_{\infty}\right\} \big\|\widehat{\bm{P}}\big\|_{\infty}^{2} \notag\\
	& \qquad +B^{2}\big\|\widehat{\bm{P}}- \bm{P}^{\star} \big\|_{\infty}\left\{ \big\|\bm{P}^\star \big\|_{\infty}+\big\|\widehat{\bm{P}}\big\|_{\infty}\right\} \\
 & \lesssim\sigma\kappa^{2}\mu r\sqrt{\frac{\log n}{n}}\cdot B\cdot\left(\frac{\mu r}{n}\right)^{2}+B^{2}\cdot\frac{(\sigma\kappa\sqrt{\mu r}+\sigma\sqrt{r\log n})}{|\lambda_{r}^{\star}|}\sqrt{\frac{\mu r}{n}}\cdot\frac{\mu r}{n} \\
 & \lesssim B\sigma\kappa^{2}\mu r\sqrt{\frac{\log n}{n}}\cdot\frac{\mu r}{n}+\frac{\mu r}{n}\cdot\frac{(\sigma\kappa\sqrt{\mu r}+\sigma\sqrt{r\log n})}{|\lambda_{r}^{\star}|}\sqrt{\frac{\mu r}{n}}\cdot\sigma^{2}n, 
\end{align*}
%
where the penultimate inequality is a consequence of \eqref{eq:entrywise-estimation-E}, \eqref{eq:entrywise-estimation-E-1}, \eqref{eq:P-inf-norm-UB}, \eqref{eq:P-hat-Linf-norm-diff} and \eqref{eq:P-hat-Linf-norm},
and the last line holds true as long as $\mu r \leq n$ (cf.~\eqref{eq:mu-constraint-general}) and $B\lesssim \sigma \sqrt{n}$. 
Combining the above inequalities allows one to reach
%
\begin{align*}
 & \big| \widehat{v}_{i,j} - \widetilde{v}_{i,j} \big|  \leq\alpha_{1}+\alpha_{2}+\alpha_{3}\\
 & \qquad \lesssim  
	 B\sigma\kappa^{2}\mu r\sqrt{\frac{\log n}{n}}\cdot\frac{\mu r}{n}+\frac{\mu r}{n}\cdot\frac{\sigma\kappa\sqrt{\mu r}+\sigma\sqrt{r\log n}}{|\lambda_{r}^{\star}|}\sqrt{\frac{\mu r}{n}}\cdot\sigma^{2}n \\
 & \qquad \asymp \frac{B\sigma\kappa^{2}\mu^{2}r^{2}\sqrt{\log n}}{n^{3/2}}+\frac{\sigma^{3}\kappa\mu^{2}r^{2}+\sigma^{3}\mu^{3/2}r^{2}\sqrt{\log n}}{|\lambda_{r}^{\star}|\sqrt{n}}
\end{align*}
%
with probability at least $1-O(n^{-5})$. 
%provided that $\mu r\leq n$. 




\subsubsection{Step 3: combining the above bounds} 

Putting together the results in the previous steps, we can readily derive
%
\begin{align}
\big|\widehat{v}_{i,j}-v_{i,j}^{\star}\big| & \leq\big|\widetilde{v}_{i,j}-v_{i,j}^{\star}\big|+ \big| \widehat{v}_{i,j} - \widetilde{v}_{i,j} \big|  \notag\\
 & \lesssim\left(\frac{\sigma B\mu^{3/2}r^{3/2}\sqrt{\log n}}{n^{3/2}}+\frac{\mu^{2}r^{2}B^{2}\log n}{n^{2}}\right) \notag\\
 & \qquad+\left(\frac{B\sigma\kappa^{2}\mu^{2}r^{2}\sqrt{\log n}}{n^{3/2}}+\frac{\sigma^{3}\kappa\mu^{2}r^{2}+\sigma^{3}\mu^{3/2}r^{2}\sqrt{\log n}}{|\lambda_{r}^{\star}|\sqrt{n}}\right) \notag\\
	& \asymp \frac{B\sigma\kappa^{2}\mu^{2}r^{2}\sqrt{\log n}}{n^{3/2}}+\frac{\sigma^{3}\kappa\mu^{2}r^{2}\sqrt{\log n}}{|\lambda_{r}^{\star}|\sqrt{n}},  \label{eq:vij-vhat-ij-gap-bound}
 %& \lesssim\frac{1}{\sqrt{\log n}}v_{i,j}^{\star},
\end{align}
%
where the last relation is guaranteed as long as $B\lesssim \sigma \sqrt{n/\log n}$. 
Consequently, if Condition~\eqref{eq:Uj-Ui-lower-bound-lem} holds, 
%$\frac{B\sigma\kappa^{2}\mu^{2}r^{2}\log n}{n^{3/2}}+\frac{\sigma^{3}\kappa\mu^{2}r^{2}\log n}{|\lambda_{r}^{\star}|\sqrt{n}}
%\lesssim\sigma_{\min} ^2 \big\{ \|\bm{U}_{j,\cdot}^{\star}\|_{2}^{2}+\|\bm{U}_{i,\cdot}^{\star}\|_{2}^{2} \big\} $, 
then it follows from \eqref{eq:vij-vhat-ij-gap-bound} and the lower bound \eqref{eq:vij-lower-bound-12} that
%
\begin{align*}
	\big|\widehat{v}_{i,j}-v_{i,j}^{\star}\big| \lesssim \frac{1}{\sqrt{\log n}} v_{i,j}^{\star}, 
\end{align*}
%
thus concluding the proof. 





 
%\end{subappendices}







\section{Notes}







\paragraph{Leave-one-out analysis.} The core idea of leave-one-out analysis,  which drops a small amount of randomness to decouple complicated statistical dependency, is deeply rooted in the probability and statistics literature. For instance, an idea of this kind was invoked by \citet{stein1972a} to help establish normal approximation,   was paired with the Stieltjes transform to establish the  limiting spectral law of random matrices (see, e.g., \citep[Section 2.4.3]{Tao2012RMT}), and bears some resemblance to the cavity method in statistical physics \citep{mezard2009information}.
When it comes to statistical estimation, a prominent series of work that unveiled the striking effectiveness of leave-one-out analysis was \citet{el2013robust,el2015impact}, which characterized rigorously the sharp statistical performance (including pre-constants) of M-estimators in high dimension (i.e., a challenging regime where the number of samples is comparable to the number of unknown parameters).
The deep analysis framework developed in these papers inspired much of the follow-up work presented in this chapter. Particularly worth mentioning are: (1) \citet{zhong2017near}: which was the first to determine the entrywise behavior of the generalized projected power method; (2) \citet{abbe2020entrywise,chen2017spectral}: which extended the leave-one-out analysis idea to establish entrywise eigenvector perturbation; and (3) \citet{ma2017implicit,chen2019gradient}: which characterized tight convergence guarantees for nonconvex optimization algorithms with the aid of leave-one-out ideas. For readers' reference, we list below several topics for which leave-one-out analyses prove useful:
%
\begin{itemize}
	\item Maximum likelihood estimation and M-estimation: \citet{el2013robust,el2015impact,lei2018asymptotics,sur2019likelihood,sur2019modern,chen2017spectral,chen2020partial};
	\item spectral methods: \citet{chen2017spectral,abbe2020entrywise,ma2017implicit,cai2019subspace,lei2019unified,abbe2020ell_p,ling2020near,chen2020partial};
	\item nonconvex optimization for statistical estimation: \citet{ma2017implicit,chen2019gradient, li2019nonconvex,chen2020nonconvex,cai2019tensor,dong2018nonconvex,chen2020convex,wang2021entrywise};
	\item semidefinite relaxation for low-rank matrix factorization: \citet{zhong2017near,ding2020leave,chen2020noisy,chen2020bridging,chen2020convex};
	\item uncertainty quantification and  confidence intervals: \citet{javanmard2018debiasing,chen2019inference,cai2020uncertainty,yan2021inference};
	\item cross validation: \citet{xu2019consistent};
	\item reinforcement learning: \citet{agarwal2020model,li2020breaking,pananjady2020instance,zhang2020model,cui2020minimax,
		wang2021sample}.
\end{itemize}



\paragraph{$\ell_{\infty}$ and $\ell_{2,\infty}$  perturbation theory.}

The $\ell_{\infty}$ and $\ell_{2,\infty}$ perturbation theory for eigenspace and singular subspaces have been investigated in the literature \citep{fan2016ell,cape2019two,eldridge2018unperturbed}, but only scatteredly until very recently.
A modern and systematic framework was established recently,  empowered by the leave-one-out analysis idea.
Its efficacy and tightness  were first demonstrated by \citet{abbe2020entrywise} in a setting that subsumes the one presented herein (with applications to SBMs, phase synchronization and matrix completion), and  by \citet{chen2017spectral} in an asymmetric setting (with application to
 top-$K$ ranking).
 The theoretical framework has subsequently been extended in several aspects.
 For instance, (1) \citet{cai2019subspace} investigated the ``unbalanced'' scenario where the column dimension far exceeds the row dimension of the matrix, resulting in near-optimal fine-grained guarantees for PCA, bi-clustering and tensor completion; (2) \citet{lei2019unified}  expanded the setting by accounting for more flexible noise distributions (including the ones exhibiting certain dependency structure), leading to tight guarantees for, e.g., spectral clustering with more than two communities, and  hierarchical clustering;  (3)  \citet{abbe2020ell_p} explored a more general $\ell_{p}$ perturbation theory that subsumes the $\ell_{\infty}$ perturbation theory as special cases.
Another plausible  approach to study $\ell_{\infty}$ eigenvector perturbation
is to analyze instead the dynamics of an iterative procedure (e.g., the power method) that converges to the leading eigenvector \citep{zhong2017near}, again using the leave-one-out ideas.
This iterative approach offers a perspective complementary to the analysis framework presented herein,
while at the same time playing a pivotal role when studying nonconvex optimization algorithms (see \citet{chi2019nonconvex}).
Moving beyond the leave-one-out analysis framework, $\ell_{\infty}$ and $\ell_{2,\infty}$ eigenspace perturbation theory has been derived via other powerful tools as well, e.g., the Neumann trick \citep{eldridge2018unperturbed,chen2018asymmetry,cheng2020tackling,}, the Procrustes analysis \citep{cape2019two}, and more specialized techniques tailored to Gaussian ensembles \citep{koltchinskii2016perturbation,koltchinskii2016asymptotics,koltchinskii2020efficient}. Perturbations of linear forms and bilinear forms of eigenvectors have also been investigated  in the literature
\citep{koltchinskii2016perturbation,koltchinskii2016asymptotics,chen2018asymmetry,cheng2020tackling,fan2020asymptotic,koltchinskii2020efficient},
which are beyond the scope of this work.
Finally, distributional theory and uncertainty quantification for spectral methods, which
have been recently studied by \citet{xia2019normal, cheng2020tackling, fan2020asymptotic, yan2021inference,agterberg2021entrywise}, 
are still in their infancy. The results presented in Section~\ref{sec:distribution-theory} follow the analysis framework developed in \citet{yan2021inference}. 


